\pdfoutput=1
\documentclass[10pt]{amsart}
%%%%%%%%%%%%%%%%%%%%%%%%%%%%%%%%%%%%%%%%%%%%%%
%%              Packages                    %%
%%%%%%%%%%%%%%%%%%%%%%%%%%%%%%%%%%%%%%%%%%%%%%
\RequirePackage{amsthm,amsmath,amsfonts,amssymb}
\RequirePackage[numbers]{natbib}
%\RequirePackage[authoryear]{natbib} %% uncomment this for author-year bibliography
\RequirePackage[section]{algorithm}
\RequirePackage{algorithmicx, algpseudocode}
\RequirePackage{color, xcolor}
\RequirePackage{graphicx}
\RequirePackage[toc,page]{appendix}

\usepackage{math_macros}
\usepackage{bbm}
\usepackage{comment}
\usepackage{enumitem}
\usepackage{fullpage}
\usepackage{subcaption}
\usepackage{hyperref}
\hypersetup{colorlinks,                              linkcolor=red,
            citecolor=blue,
            urlcolor=blue}
\usepackage{cleveref}
%%%%%%%%%%%%%%%%%%%%%%%%%%%%%%%%%%%%%%%%%%%%%%
%%                                          %%
%% Uncomment next line to change            %%
%% the type of equation numbering           %%
%%                                          %%
%%%%%%%%%%%%%%%%%%%%%%%%%%%%%%%%%%%%%%%%%%%%%%
\numberwithin{equation}{section}
%%%%%%%%%%%%%%%%%%%%%%%%%%%%%%%%%%%%%%%%%%%%%%
%%                                          %%
%% For Axiom, Claim, Corollary, Hypothezis, %%
%% Lemma, Theorem, Proposition              %%
%% use \theoremstyle{plain}                 %%
%%                                          %%
%%%%%%%%%%%%%%%%%%%%%%%%%%%%%%%%%%%%%%%%%%%%%%
\theoremstyle{plain}
\newtheorem{theorem}{Theorem}[section]
\newtheorem{axiom}{Axiom}
\newtheorem{claim}[axiom]{Claim}
\newtheorem{corollary}[theorem]{Corollary}
\newtheorem{lemma}[theorem]{Lemma}
\newtheorem{proposition}[theorem]{Proposition}
\newtheorem{definition}[theorem]{Definition}
\newtheorem{remark}[theorem]{Remark}
\newtheorem{conjecture}[theorem]{Conjecture}
%%%%%%%%%%%%%%%%%%%%%%%%%%%%%%%%%%%%%%%%%%%%%%
%%                                          %%
%% For Assumption, Definition, Example,     %%
%% Notation, Property, Remark, Fact         %%
%% use \theoremstyle{remark}                %%
%%                                          %%
%%%%%%%%%%%%%%%%%%%%%%%%%%%%%%%%%%%%%%%%%%%%%%
\theoremstyle{remark}
\newtheorem{assumption}[theorem]{Assumption}
\newtheorem*{example}{Example}
\newtheorem*{fact}{Fact}
\newtheorem{notation}[theorem]{Notation}
%%%%%%%%%%%%%%%%%%%%%%%%%%%%%%%%%%%%%%%%%%%%%%
%% Please put your definitions here:        %%
%%%%%%%%%%%%%%%%%%%%%%%%%%%%%%%%%%%%%%%%%%%%%%
\newcommand{\todo}[1]{\textcolor{magenta}{[#1]}}

\newcommand{\vphi}{\varphi}
\renewcommand{\epsilon}{\varepsilon}
\renewcommand{\algorithmicrequire}{ \textbf{Input:}} % Use Input in the format of Algorithm  
\renewcommand{\algorithmicensure}{ \textbf{Output:}}
\renewcommand{\algorithmiccomment}[1]{\hfill$\triangleright$\textit{#1}}


\begin{document}

\title{Partial recovery and weak consistency in the non-uniform hypergraph Stochastic Block Model}
% Author information

\author{Ioana Dumitriu}
\address{Department of Mathematics, University of California, San Diego, La Jolla, CA 92093}
\email{idumitriu@ucsd.edu}
% \thanks{}

\author{Haixiao Wang}
\address{Department of Mathematics, University of California, San Diego, La Jolla, CA 92093}
\email{h9wang@ucsd.edu}
% \thanks{}

\author{Yizhe Zhu}
\address{Department of Mathematics, University of California, San Diego, La Jolla, CA 92093}
\email{yiz084@ucsd.edu}
% \thanks{}
%  \thanks will become a 1st page footnote.

\date{\today}

\begin{abstract}

 We consider the community detection problem in sparse random hypergraphs under the non-uniform hypergraph stochastic block model (HSBM), a  general model of random networks with community structure and  higher-order interactions. When the random hypergraph has bounded expected degrees, we provide a  spectral algorithm that outputs a partition with at least a $\gamma$ fraction of the vertices  classified correctly, where $\gamma\in (0.5,1)$ depends on the signal-to-noise ratio (SNR) of the model. When the SNR  grows slowly as the number of vertices goes to infinity, our algorithm achieves weak consistency, which improves the previous  results in  %Ghoshdastidar and Dukkipati (2017)
 \cite{ghoshdastidar2017consistency} for non-uniform HSBMs.
 
 Our spectral algorithm consists of  three major steps: (1) Hyperedge selection: select hyperedges of certain sizes to provide the maximal signal-to-noise ratio for the induced sub-hypergraph; (2) Spectral partition: construct a regularized adjacency matrix and obtain an approximate partition based on singular vectors; (3) Correction and merging: incorporate the hyperedge information from adjacency tensors to upgrade the error rate guarantee. The theoretical analysis of our algorithm relies on the concentration and regularization of the adjacency matrix for sparse non-uniform random hypergraphs, which can be of independent interest.
\end{abstract}

\maketitle

%\tableofcontents

\input{introduction}
\input{multiple}
%----------------------------------
% Bibliography
%----------------------------------
\bibliographystyle{plain}
\bibliography{ref.bib}

\appendix
\addcontentsline{toc}{section}{Appendices}
\input{concentration}
\input{appendix}
\input{binary}
\end{document}

%------------------------------------------------------------------------------
% End of journal.tex
%------------------------------------------------------------------------------
	 
\section{Introduction}


Clustering is one of the central problems in network analysis and machine learning \cite{newman2002random,shi2000normalized,ng2002spectral}.  Many
clustering algorithms make use of graph models, which represent pairwise relationships among data.  A well-studied probabilistic model is the stochastic block model (SBM), which was first introduced in \cite{holland1983stochastic} as a random graph model that generates community structure with  given ground truth for clusters so that one can study algorithm accuracy. The past decades have brought many notable results in the analysis of different algorithms and fundamental limits for community detection in SBMs in different settings. A major breakthrough was the proof of phase transition behaviors of community detection algorithms in various connectivity regimes \cite{massoulie2014community,bordenave2018nonbacktracking,mossel2015reconstruction,mossel2018proof,mossel2016consistency,abbe2016exact,abbe2018proof}. See the survey \cite{abbe2018community} for more references.

Hypergraphs can represent more complex relationships among data \cite{benson2016higher,battiston2020networks},  including recommendation systems \cite{bu2010music,li2013news}, computer vision \cite{govindu2005tensor,wen2019learning}, and biological networks \cite{michoel2012alignment,tian2009hypergraph}, and they have been shown empirically to have advantages over graphs \cite{zhou2007learning}.  Besides community detection problems,
sparse hypergraphs and their spectral theory have also found applications in data science \cite{jain2014provable,zhou2021sparse,harris2021}, combinatorics \cite{dumitriu2021spectra,friedman1995second,soma2019spectral}, and statistical physics \cite{caceres2021sparse,sen2018optimization}.  

With the motivation given by this broad set of applications,  many efforts have been made in recent years to study community detection on random hypergraphs.
The hypergraph stochastic block model (HSBM), as a generalization of graph SBM, was  first introduced and studied in \cite{ghoshdastidar2014consistency}. In this model, we observe a random uniform hypergraph where each  hyperedge appears independently with some given probability depending on the community structure of the vertices in the hyperedge. 

Succinctly put, the HSBM recovery problem is to find the ground truth clusters either approximately or exactly, given a sample hypergraph and estimates of model parameters. We may ask the following questions about the quality of the solutions (see \cite{abbe2018community} for further details in the graph case).
\begin{enumerate}
	\item \textbf{Strong consistency (a.k.a. exact recovery):} With high probability, find all clusters exactly (up to permutation).
	\item \textbf{Weak consistency:} With high probability, find a partition of the vertex set such that at most $o(n)$ vertices are misclassified.
	\item \textbf{Partial recovery:} Given a fixed $\gamma\in (0.5,1)$, with high probability, find a partition of the vertex set such that at least a fraction $\gamma$ of the vertices are clustered correctly.
	\item \textbf{Detection:} With high probability, find a partition  correlated with the true partition.
\end{enumerate}




% A faithful representation of a hypergraph is its adjacency tensor. However, most of the computations involving tensors are NP-hard \cite{hillar2013most}. Instead,
% they considered spectral algorithms for exact recovery using hypergraph Laplacians. Subsequently, the analysis was extended to sparse, non-uniform hypergraphs \cite{ghoshdastidar2015provable,ghoshdastidar2015spectral,ghoshdastidar2017consistency}. 





For exact recovery of uniform HSBMs, it was shown that the phase transition occurs in the regime of logarithmic expected  degrees in \cite{lin2017fundamental,chien2018minimax,chien2018community} and the exact threshold was given in \cite{kim2018stochastic,zhang2021exact}, by a generalization of the techniques in \cite{abbe2016exact,abbe2015community}. 
%  In \cite{ghoshdastidar2014consistency}, the authors used spectral clustering based on the hypergraph's Laplacian to recover HSBMs that are dense and uniform. Subsequently, they extended their results to sparse, non-uniform hypergraphs  in \cite{ghoshdastidar2015provable,ghoshdastidar2015spectral,ghoshdastidar2017consistency}. 
 Spectral methods were considered in \cite{chien2018community,ahn2016community,cole2020exact,zhang2021exact,yuan2021community}, while semidefinite programming methods were analyzed in \cite{kim2018stochastic,lee2020robust}.  Weak consistency for HSBMs was studied in \cite{chien2018community,chien2018minimax,ghoshdastidar2017consistency,ghoshdastidar2017uniform,ke2020community}.  For detection of the HSBM, the authors of \cite{angelini2015spectral} proposed a conjecture that the phase transition occurs in the regime of constant expected degrees, and the positive part of the conjecture for the binary case was solved in \cite{Pal_2021}.  
 
 
 



\subsection{Non-uniform hypergraph stochastic block model}
The non-uniform HSBM  was first studied in \cite{ghoshdastidar2017consistency}, which removed the uniform hypergraph assumption in previous works, and it is a more realistic model to study higher-order interaction on networks  \cite{lung2018hypergraph,wen2019learning}. It can be seen as a superposition of several uniform HSBMs with different  model parameters.  We first define the uniform HSBM in our setting and extend it to non-uniform hypergraphs. 

\begin{definition}[Uniform HSBM]
Let $V=\{V_1,\dots V_k\}$ be a partition of the set $[n]$ into $k$ blocks of size $n/k$ (assuming $n$ is divisible by $k$). For any set of $m$ distinct vertices $i_1,\dots i_m$, a hyperedge $\{i_1,\dots i_m\}$ is generated with probability $a_m/ \binom{n}{m-1}$ if the vertices $i_1,\dots i_m$ are in the same block; otherwise with probability $b_m / \binom{n}{m-1}$. We denote this distribution on the set of $m$-uniform hypergraphs as 
\begin{equation}\label{eqn:uniform_HSBM}
 H_m\sim  \textnormal{HSBM}_{m} \left( \frac{n}{k}, \frac{a_m}{\binom{n}{m-1}}, \frac{b_m}{\binom{n}{m-1}} \right)\,, \quad  2 \leq m \leq M\,,
\end{equation}
where $a_m/ \binom{n}{m-1}$ and $b_m/ \binom{n}{m-1} $ denote the connecting probabilities within and across blocks, respectively. 
\end{definition}

\begin{definition}[Non-uniform HSBM]\label{def:non_uni_hyper}
Let $H = (V, E)$ be a non-uniform random hypergraph, which can be considered as a collection of $m$-uniform hypergraphs, i.e., $H = \bigcup_{m=2}^{M}H_m$ with each $H_m$ from \Cref{eqn:uniform_HSBM}.
\end{definition}

Examples of $2$-uniform and $3$-uniform \textnormal{HSBM}, and  an example of non-uniform \textnormal{HSBM} with $M=3$ and $k=3$ can be seen in \Cref{fig:non_uniform}.
\begin{figure}
     \centering
     \begin{subfigure}{0.25\textwidth}
         \centering
        \includegraphics[width=\textwidth]{fig/2_uniform.png}
        %\subcaption{2-uniform HSBM}
        %\label{fig:2_uniform}
     \end{subfigure}
     \begin{subfigure}{0.25\textwidth}
         \centering
        \includegraphics[width=\textwidth]{fig/3_uniform.png}
        %\subcaption{3-uniform HSBM}
        %\label{fig:3_uniform}
     \end{subfigure}
     \begin{subfigure}{0.25\textwidth}
         \centering
        \includegraphics[width=\textwidth]{fig/non_uniform.png}
        %\subcaption{Non-uniform HSBM}
        %\label{fig:non_uniform}
     \end{subfigure}
     \caption{Uniform \textnormal{HSBM} and Non-uniform \textnormal{HSBM}.}
     \label{fig:non_uniform}
\end{figure}
    
\subsection{Main results}
In this paper, we consider \textit{partial recovery} and \textit{weak consistency} in the non-uniform HSBM given in \Cref{def:non_uni_hyper}. The accuracy of the recovery is measured  as follows.

\begin{definition}[$\gamma$-correctness]
Suppose we have $k$ disjoint blocks $V_1,\dots, V_k$. A collection of subsets $\widehat{V}_1,\dots, \widehat{V}_k$ of $V$ is $\gamma$-correct if $|V_i\cap \widehat{V}_i|\geq \gamma |V_i|$ for all $i\in [k]$. Note that we do not require $\widehat{V}_1,\dots, \widehat{V}_k$ to be a partition of $V$.
\end{definition}


The following theorem provides an algorithm that outputs a $\gamma$-correct partition of a non-uniform HSBM with high probability, when $k\geq 3$ is fixed.

\begin{theorem}[$k\geq 3$]\label{thm:mainresult_k}
Given any $\nu\in (0.5, 1)$,  assume $a_m \geq b_m$ are constants independent of $n$ for all $m\in \mathcal{M}$,  where $\mathcal{M}\subset \{2, \dots, M\}$ is obtained from \Cref{alg:parameter_preprocessing}, and
\begin{align}
    d:=\sum_{m\in \mathcal M} (m-1)a_m &\geq C, \notag\\
    \sum_{m\in\mathcal M} (m-1)(a_m-b_m) &\geq C_{\nu} \cdot k^{\mathcal M_{\max}-1}\sqrt{\log\Big(\frac{k}{1-\nu}\Big)d \,\,}\label{eqn:SNRcondition}
\end{align}
for some constants $C, C_{\nu}$.
Then with probability $1-O(n^{-2})$, for sufficiently large $n$, Algorithm \ref{alg:k_block_partition} outputs a $\gamma$-correct partition with $\gamma = \max\{\nu,\, 1 - k\rho\}$ and $\rho:=\exp\left( -C_{\mathcal M} \cdot \mathrm{SNR}_{\mathcal{M}} \right)$, where
\begin{align}\label{eqn:SNR_CM}
    \mathrm{SNR}_{\mathcal{M}}:= \frac{ \left[\sum_{m\in \mathcal{M}} (m-1)\left(\frac{a_m - b_m}{k^{m-1}} \right) \right]^2 }{\sum_{m\in \mathcal{M}} (m-1)\left(\frac{a_m - b_m}{k^{m-1}} + b_m \right)}\,, \quad C_{\mathcal M}:=\frac{[\nu^{\mathcal{M}_{\max}-1} - (1-\nu)^{\mathcal{M}_{\max}-1}]^2 }{(\mathcal{M}_{\max} - 1)^2\cdot 2^{2\mathcal{M}_{\max} + 3}}\,.
\end{align}
\end{theorem}
\begin{remark}
Taking $M=2$, \Cref{thm:mainresult_k} can  be reduced to  \cite[Lemma 9]{chin2015stochastic} for the graph case. The failure probability $O(n^{-2})$ can be decreased to $O(n^{-p})$ for any $p>0$, as long as one is willing to pay a price in making the graph denser (increasing d). 


Our algorithm for  \Cref{thm:mainresult_k} requires the input of model parameters $a_m,b_m, m\in\mathcal M$, and they can be estimated by counting cycles in hypergraphs as shown in \cite{mossel2015reconstruction,yuan2021testing}. Estimation of the number of blocks can be done by counting the outliers in the spectrum of the   non-backtracking operator \cite{le2015estimating,angelini2015spectral}.
\end{remark}



Our \Cref{alg:k_block_partition} can be summarized in 3 steps:
\begin{enumerate}
    \item \textbf{Hyperedge selection}: select hyperedges of certain sizes to provide the maximal signal-to-noise ratio (SNR) for the induced sub-hypergraph.
    \item \textbf{Spectral partition}: construct a regularized  adjacency matrix and obtain an approximate partition based on singular vectors (first approximation). 
    \item \textbf{Correction and merging}: incorporate the hyperedge information from adjacency tensors to upgrade the error rate guarantee (second, better approximation). 
\end{enumerate}

\begin{algorithm}
    \caption{\textbf{Partition}}\label{alg:k_block_partition}
    \begin{algorithmic}[1]	
        \Require The adjacency tensors $\bcT_m$, $k$, $a_m$, $b_m$ for $m\in \{2, \cdots, M\}$. 
        \Ensure The corrected sets $\widehat{V}_{1}, \cdots, \widehat{V}_k$.
        
		\Function{\textbf{Partition}}{$\bA$, $a_m$, $b_m$ for $m\in \mathcal{M}$}
		\State Run \textbf{\Cref{alg:parameter_preprocessing} Pre-processing} to obtain the subset $\mathcal{M}$ which achieves maximal $\mathrm{SNR}$.
		
		\State Randomly color the hyperedges red or blue with equal probability.
		
		\State Randomly partition $V$ into $2$ disjoint subsets $Y$ and $Z$ by assigning $+1$ or $-1$ to each vertex with equal probability.
		
		\State Let $\bB$ denote the adjacency matrix of the bipartite hypergraph between $Y$ and $Z$ consisting only of the red hyperedges, with rows indexed by $Z$ and columns indexed by $Y$.

		\State Run \textbf{\Cref{alg:k_block_spectral_partition} Spectral Partition} on the red hypergraph and output $U_{1}^{\prime}, \cdots, U_{k}^{\prime}$.  \newline \textit{This step only uses the red hyperedges between vertices in $Y$ and $Z$ and outputs approximate clusters for $U_{i}:= V_{i}\cap Z$, with $i=1,\ldots, k$}
		
		\State Run \textbf{\Cref{alg:k_block_correction} Correction} on the red hypergraph and output $\widehat{U}_{1}, \cdots, \widehat{U}_k$. %\Comment{This procedure only uses the red hyperedges that are internal to $Z$ and improves the clustering in $Z$.}
		
		\State Run \textbf{\Cref{alg:k_block_merging} Merging} on the blue hypergraph and output $\widehat{V}_{1}, \cdots, \widehat{V}_k$. \newline\textit{This step only uses the blue hyperedges  between vertices in $Y$ and $Z$ and assigns the vertices in $Y$ to an appropriate approximate cluster.}
	    \EndFunction
	    \end{algorithmic} 
\end{algorithm}

We generalize the innovative graph algorithm in \cite{chin2015stochastic} to uniform hypergraphs with  a more detailed analysis. In particular, the random partitioning step in \Cref{alg:2_block_partition} was explained only in the case when  the output is an equi-partition; in general, the output sets $\widehat{V}_1, \ldots, \widehat{V}_k$ are only approximately equal in size,  and we give complete explanations of the general case here.

A non-uniform HSBM can be seen as a collection of noisy observations for the same underlying community structure through several uniform HSBMs of different orders. A possible issue is that some uniform hypergraphs with small SNR might not be informative (if we observe an $m$-uniform hypergraph with parameters $a_m=b_m$,  including  hyperedge information from it ultimately increases the noise). To improve our error rate guarantees, we start by adding a pre-processing step for hyperedge selection according to SNR,  and then apply the algorithm on the sub-hypergraph with maximal SNR.

When the number of blocks is $2$, we give a simpler algorithm (\Cref{alg:2_block_partition}) for partial recovery with a better error rate guarantee, as shown in the following \Cref{thm:mainresult_2}. \Cref{alg:2_block_partition} does not need the merging routine in \Cref{alg:k_block_partition}, since once we cluster one block, the other one is clustered automatically.

\begin{theorem}[$k=2$]\label{thm:mainresult_2}
Given $\nu\in(0.5, 1)$, assume $a_m \geq b_m$ are constants independent of $n$ for all $m\in \mathcal{M}$, 
\begin{subequations}
\begin{align*}
    d:=\sum_{m\in \mathcal M} (m-1)a_m &\geq C\,,\\
    \sum_{m\in\mathcal M}(m-1)(a_m-b_m) &\geq C_{\nu} \cdot 2^{\mathcal M_{\max}+2}\sqrt{d}\,,
  \end{align*}
\end{subequations}
 for some constants $C, C_{\nu}$. Then with probability $1 - O(n^{-2})$, for sufficiently large $n$, \Cref{alg:2_block_partition} outputs a $\gamma$-correct partition with $\gamma = \max\{\nu,\, 1 - 2\rho\}$ and $\rho = 2\exp( -C_{\mathcal M} \cdot \mathrm{SNR}_{\mathcal{M}})$, where 
\begin{align}\label{eqn:SNR_CM_2}
    \mathrm{SNR}_{\mathcal{M}}:= \frac{ \left[\sum_{m\in \mathcal{M}} (m-1)\left(\frac{a_m - b_m}{2^{m-1}} \right) \right]^2 }{\sum_{m\in \mathcal{M}} (m-1)\left(\frac{a_m - b_m}{2^{m-1}} + b_m \right)}\,, \quad C_{\mathcal M}:=\frac{[(\nu)^{\mathcal{M}_{\max}-1} - (1 - \nu)^{\mathcal{M}_{\max}-1}]^2 }{8(\mathcal{M}_{\max} - 1)^2}\,.
\end{align}
\end{theorem}



\begin{algorithm}[ht]
	\caption{\textbf{Binary Partition}} \label{alg:2_block_partition}
    \begin{algorithmic}[1]
        \Require The adjacency tensors $\bcT^{(m)}$, $a_m$, $b_m$ for $m\in \{2, \cdots, M\}$. 
        \Ensure The corrected sets $\widehat{V}_{1}, \widehat{V}_2$.
		\Function{\textbf{Partition}}{$\bcT^{(m)}$, $a_m$, $b_m$ for $m\in \{2, \cdots, M\}$}
		\State Run \textbf{\Cref{alg:parameter_preprocessing} Pre-processing} to obtain the subset $\mathcal{M}$ which achieves maximal $\mathrm{SNR}$.
		
		\State Randomly color the hyperedges red or blue with equal probability.
		
		\State Run \textbf{\Cref{alg:2_block_spectral_partition} Spectral Partition} on the red hypergraph and output $V_{1}^{\prime}, V_{2}^{\prime}$.
		\State Run \textbf{\Cref{alg:2_block_correction} Correction} on the blue hypergraph and output $\widehat{V}_{1}, \widehat{V}_2$.
		\State \Return{the corrected sets  $\widehat{V}_{1}, \widehat{V}_2$.}
	    \EndFunction
	   \end{algorithmic}  
\end{algorithm} 

Throughout the proofs for Theorems \ref{thm:mainresult_k} and \ref{thm:mainresult_2}, we make only one assumption on the growth or finiteness of $d$ and $\mathrm{SNR}_{\mathcal{M}}$, and that happens in estimating the failure probability in the proof of \Cref{lem:discard_half_blue} as noted in Remark \ref{rem:SNRlogn}. As a consequence, we can also give the following corollary, which covers the case when $d$ and  $\mathrm{SNR}_{\mathcal{M}}$ grow with $n$.

\begin{corollary}[Weak consistency]\label{cor:weak_consistency}
For fixed $M$ and $k$, if
\[
\mathrm{SNR}_{\mathcal{M}}:= \frac{ \left[\sum_{m\in \mathcal{M}} (m-1)\left(\frac{a_m - b_m}{k^{m-1}} \right) \right]^2 }{\sum_{m\in \mathcal{M}} (m-1)\left(\frac{a_m - b_m}{k^{m-1}} + b_m \right)} \to\infty\,,
\]
and $\mathrm{SNR}_{\mathcal{M}} = o(\log n)$, then with probability $1- O(n^{-2})$, \Cref{alg:k_block_partition} outputs a partition with only $o(n)$  misclassified vertices.
 \end{corollary}
Here, the failure probability is of order $O(\exp[-n \exp(-\mathrm{SNR}_{\mathcal{M}})])$, and it fails to go to zero if $\mathrm{SNR}_{\mathcal{M}}$ is not smaller than $o(\log(n))$.

\subsection{Comparison with existing results} Although many algorithms and theoretical results have been developed for hypergraph community detection, most of them are restricted to uniform hypergraphs, and few results are known for non-uniform ones. We will discuss the most relevant ones here and compare our results to theirs. 

In \cite{ke2020community}, the authors studied the degree-corrected HSBM with general connection probability parameters by using a tensor power iteration algorithm and Tucker decomposition. Their algorithm achieves weak consistency for uniform hypergraphs when the average degree is $\omega(\log^2 n)$, which is the regime complementary to the regime we studied here. They discussed a way to generalize the algorithm to non-uniform hypergraphs, but the theoretical analysis remains open.  The recent paper \cite{zhen2021community} analyzed non-uniform hypergraph community detection by using hypergraph embedding and optimization algorithms and obtained weak consistency when the expected degrees are of $\omega(\log n)$, again a complementary regime to ours. Results on spectral norm concentration of sparse random tensors  were obtained in \cite{cooper2020adjacency,nguyen2015tensor,jain2014provable,lei2020consistent,zhou2021sparse}, but no provable tensor algorithm in the bounded expected degree is known.   Testing for the community structure for non-uniform hypergraphs was studied in \cite{yuan2021testing,jin2021sharp}, which is a problem different from community detection.

In our approach, we relied on knowing the tensors for each uniform hypergraph. However, in computations, we only ran the spectral algorithm on the adjacency matrix of the entire hypergraph,  since the stability of tensor algorithms does not yet come with  guarantees due to the lack of concentration, and for non-uniform hypergraphs, $M-1$ adjacency tensors would be needed. This approach presented the challenge that, unlike for graphs, the adjacency matrix of a random non-uniform hypergraph has dependent entries, and the concentration properties of such a random matrix were previously unknown. We were able to overcome this issue and prove concentration bounds from scratch, down to the bounded degree regime. Similar to \cite{feige2005spectral,le2017concentration}, we provided here a regularization analysis by removing rows in the adjacency matrix with large row sums (suggestive of large degree vertices) and proving a concentration result for the regularized matrix down to the bounded expected degree regime (see \Cref{thm:regularization}).

%We now discuss the closest results existing in the literature. 

In terms of partial recovery for hypergraphs, our results are new, even in the uniform case. In \cite[Theorem 1]{ahn2018hypergraph}, for uniform hypergraphs, the authors showed detection (not partial recovery) is possible when the average degree is $\Omega(\frac{1}{n})$; in addition, the error rate is not exponential in the model parameters.  In the graph case, it was shown in \cite{zhang2016minimax} that the error rate in \Cref{eqn:SNR_CM} is optimal up to a constant in the exponent. It's an interesting open problem to extend the analysis in \cite{zhang2016minimax} to obtain a minimax error rate for hypergraphs. 

In \cite{ghoshdastidar2017consistency}, the authors considered  weak consistency in  a non-uniform HSBM model with a spectral algorithm based on the hypergraph Laplacian matrix, and  showed that  weak consistency is achievable if the expected degree is of $\Omega(\log^2 n)$  with high probability \cite[Theorem 4.2]{ghoshdastidar2014consistency}. Their algorithm can't be applied to sparse regime straightforwardly since the normalized Laplacian is not well-defined due to the existence of isolated vertices in the bounded degree case. In addition, our weak consistency results obtained here are valid as long as the expected degree is $\omega(1)$ and $o(\log n)$, which is the entire set of problems on which weak consistency is expected. By contrast, in \cite{ghoshdastidar2017consistency}, weak consistency is shown only when the  expected degree is $\Omega(\log^2(n))$, which is a regime complementary to ours and where exact recovery should (in principle) be possible: for example, this is known to be an exact recovery regime in the uniform case \cite{chien2018minimax,kim2018stochastic,lee2020robust,zhang2021exact}.

Finally, although our analysis does not provide a way to identify the exact threshold for the phase transition of detection in the non-uniform hypergraph, we conjecture that (similar to \cite{angelini2015spectral} in the uniform case) $ \mathrm{SNR}_{\mathcal{M}}=1$ in \Cref{eqn:SNR} is the exact threshold for detection.



\subsection{Organization of the paper}
In \Cref{sec:prelim}, we include the definitions of adjacency matrices of hypergraphs. The concentration results for the adjacency matrices are provided in \Cref{sec:concentration}. The algorithms for partial recovery are presented in Section \ref{sec:spec_algorithm}. The proof for correctness of our algorithms for Theorem \ref{thm:mainresult_k} and \Cref{cor:weak_consistency} are given in Section \ref{sec:analysis_k}. The proof of Theorem \ref{thm:mainresult_2} as well as the proofs of many auxiliary lemmas, and useful lemmas in the literature are provided in the Appendices.






% \todo{this is standard, we don't need to write}
% \begin{enumerate}
%     \item $f(n)=O(g(n))$ or $f(n)\lesssim g(n)$ means that there exists a constant $c > 0$ such that $|f (n)| \leq c|g(n)|$ holds for all sufficiently large $n$;
%     \item $f(n)=o(g(n))$ indicates that $|f(n)|/|g(n)| \to 0$ as $n\to \infty$;
%     \item $f(n) = \Omega(g(n))$ or $f(n)\gtrsim g(n)$ means that there exists a constant $c > 0$ such that $|f (n)| \geq c|g(n)|$ holds for all sufficiently large $n$;
%     \item $f(n) \asymp g(n)$ means that there exist constants $c_1, c_2 > 0$ such that $c_1|g(n)|\leq |f(n)| \geq c_2|g(n)|$ holds for all sufficiently large $n$;
%     \item We sometimes use $f(n) \ll g(n)$ (resp. $f(n) \gg g(n)$) to indicate that there exists some very small (resp. large) universal constant $c>0$ such that $|f (n)| \leq c |g(n)|$ (resp. $|f (n)| \geq c |g(n)|$).
% \end{enumerate}



\section{Preliminaries}\label{sec:prelim}

% \subsubsection*{Notation}
% The lowercase boldface letters ($\bu$, $\bv$) and uppercase boldface letters ($\bA$, $\bU$), are used to denote vectors and matrices, respectively. The boldface calligraphic letters ($\bcT$) are used to denote tensors. The $\mathbbm{1}$ denotes the indicator function.

\begin{definition}[Adjacency tensor]\label{def:adjacency_tensor}
Given an $m$-uniform hypergraph $H_m=([n], E_m)$, we can associate to it an order-$m$ adjacency tensor $\bcT^{(m)}$. For any $m$-hyperedge $e = \{ i_1, \dots, i_m \}$, let $\bcT^{(m)}_e$ denote the corresponding entry $\bcT_{[i_{1},\dots,i_{m}]}^{(m)}$, such that
\begin{equation}\label{eqn:adjacency_tensor}
    \bcT_{e}^{(m)} := \bcT_{[i_1,\dots, i_m]}^{(m)} = \mathbbm{1}_{\{ e \in E_m\} }\,. %\left\{
        %\begin{aligned}
            %&1& \textnormal{ if } e = \{i_1,\dots, %i_m\}\in E_m,\\
            %&0\,, & \textnormal{ otherwise }.
        %\end{aligned}
    %\right.
\end{equation}
\end{definition}

\begin{definition}[Adjacency matrix]
For the non-uniform hypergraph $H$ (Definition \ref{def:non_uni_hyper}), let $\bcT^{(m)}$ be the order-$m$ adjacency tensor of the underlying $m$-uniform hypergraph for each $m\in\{2, \dots, M\}$. The adjacency matrix $\bA := [\bA_{ij}]_{n \times n}$ of the non-uniform hypergraph $H$ is defined by
\begin{equation}\label{eqn:adjacency_matrix_entry}
    \bA_{ij} = \mathbbm{1}_{\{ i \neq j\} } \cdot \sum_{m=2}^{M}\,\,\sum_{\substack{e\in E_m\\ \{i,j\}\subset e} }\bcT^{(m)}_{e}\,. %\left\{
        %\begin{aligned}
            %&\sum_{m=2}^{M}\,\,\sum_{\substack{e\in E_m\\ \{i,j\}\subset e} }\bcT^{(m)}_{e}\,,& \textnormal{if } i\neq j,\\
            %&0\,, & \textnormal{if } i=j.
        %\end{aligned}
    %\right.
\end{equation}
\end{definition}

We compute the  expectation of $\bA$ first. In each $H_m$, we pick two vertices $i,j\in V$ arbitrarily, then the expected number of $m$-hyperedges containing $i$ and $j$ can be computed as follows. Since our model does not allow for loops, in the below, we only consider the case $i \neq j$. Recall that $n/k \in \N$.
\begin{itemize}
    \item If $i$ and $j$ are from the same block, the $m$-hyperedge can either be sampled with probability $a_m /\binom{n}{m-1}$ when the other $m-2$ vertices are from the same block as $i$ and $j$, otherwise they are sampled with probability $b_m /\binom{n}{m-1}$. Then
    \begin{align}
        \alpha_m:=\E\bA_{ij} = \binom{\frac{n}{k} -2}{m-2}  \frac{a_m}{\binom{n}{m-1}} + \left[ \binom{n-2}{m-2} - \binom{\frac{n}{k} - 2}{m-2} \right]\frac{b_m}{\binom{n}{m-1}} \,.\notag
    \end{align}
    \item If $i$ and $j$ are not from the same block, we  sample the $m$-hyperedge with probability $b_m /\binom{n}{m-1}$, and
   \begin{align}
         \beta_m :=\E\bA_{ij} = \binom{n-2}{m-2} \frac{b_m}{\binom{n}{m-1}}\,.\notag
    \end{align}
\end{itemize}
By assumption $a_m \geq b_m$,  then $\alpha_m \geq \beta_m$ for each $m\in \{2, \dots, M\}$. Summing over $m$, the \textit{expected adjacency} matrix under the $k$-block non-uniform $\textnormal{HSBM}$ can written as 
\begin{align}\label{eqn:expected_adjacency_k}
  \E \bA = \begin{bmatrix}
     \alpha \bJ_{\frac{n}{k}}  &  \beta \bJ_{\frac{n}{k}} & \cdots & \beta \bJ_{\frac{n}{k}} \\
    \beta \bJ_{\frac{n}{k}}  & \alpha \bJ_{\frac{n}{k}} & \cdots & \beta \bJ_{\frac{n}{k}} \\
    \vdots  &  \vdots & \ddots & \vdots \\
    \beta  \bJ_{\frac{n}{k}} & \beta \bJ_{\frac{n}{k}} & \cdots & \alpha \bJ_{\frac{n}{k}}
  \end{bmatrix} - \alpha \bI_{n}\,,
\end{align}
with
\begin{align}\label{eqn:alpha_beta}
    \alpha:= \sum_{m=2}^{M} \alpha_m\,, \quad \beta:= \sum_{m=2}^{M} \beta_m\,,
\end{align}
where $\bJ_{\frac{n}{k}}\in \R^{\frac{n}{k} \times \frac{n}{k}}$ denotes the all-one matrix. The \Cref{lem:eigenvalues_HSBM} can be verified through direct computation. 
\begin{lemma}\label{lem:eigenvalues_HSBM}
The eigenvalues of $\E \bA$ are given below:
\begin{subequations}
    \begin{align*}
        \lambda_1(\E \bA) =&\, \frac{n}{k}(\alpha + (k-1)\beta) - \alpha\,,\\
        \lambda_i(\E \bA) =&\, \frac{n}{k}(\alpha - \beta) -\alpha \,, \quad \quad \quad \quad \quad\quad 2\leq i \leq k\,,\\
        \lambda_i(\E \bA) =&\, - \alpha\,, \quad \quad \quad \quad\quad \quad \quad \quad \quad k+1\leq i \leq n\,.
    \end{align*}
\end{subequations}
\end{lemma}

In our proof, \Cref{lem:general_eigen_value_approximation} is used for more approximately equal partitions, meaning that we can approximate the eigenvalues of $\widetilde{\E \bA}$ by eigenvalues $\E \bA$ when $n$ is sufficiently large.
\begin{lemma}\label{lem:general_eigen_value_approximation}
For any partition $(V_1, \dots, V_k)$ of $V$, consider the following matrix
\begin{align*}
\widetilde{\E \bA} = \begin{bmatrix}
     \alpha \bJ_{n_1} &  \beta \bJ_{n_1 \times n_2} & \cdots & \beta \bJ_{n_1 \times n_{k-1}} & \beta \bJ_{n_1 \times n_k} \\
    \beta \bJ_{n_2 \times n_1}  & \alpha \bJ_{n_2} & \cdots & \beta \bJ_{n_2 \times n_{k-1}} & \beta \bJ_{n_2 \times n_k} \\
    \vdots  &  \vdots & \ddots & \vdots & \vdots\\
    \beta \bJ_{n_{k-1} \times n_1}  & \beta \bJ_{n_{k-1} \times n_2} & \cdots & \alpha \bJ_{n_{k-1}} & \beta \bJ_{n_{k-1} \times n_k} \\
    \beta  \bJ_{n_k \times n_1} & \beta \bJ_{n_k \times n_2} & \cdots & \beta \bJ_{n_k \times n_{k-1}} & \alpha \bJ_{n_k}
  \end{bmatrix} - \alpha \bI_{n}\,,
\end{align*}
where $n_i := |V_i|$. Assume that $n_i = (n/k) + O(\sqrt{n} \log n)$ for all $i\in [k]$. Then, for all $1\leq i \leq k$, 
\begin{align*}
    \frac{|\lambda_i( \widetilde{\E \bA} ) - \lambda_i(\E \bA)|}{|\lambda_i(\E \bA)| } = O\Big(n^{-1/4}\log^{1/2}(n)\Big)\,.
\end{align*}
Note that both $(\,\widetilde{\E \bA} + \alpha \bI_{n})$ and $(\,\E \bA + \alpha \bI_{n})$ are rank $k$ matrices, then $\lambda_i(\widetilde{\E \bA}) = \lambda_i(\E \bA) = -\alpha$ for all $(k+1) \leq i \leq n$.
\end{lemma}



\begin{definition}[signal-to-noise ratio, $\mathrm{SNR}$] We define the \textit{signal-to-noise ratio} as 
\begin{align}\label{eqn:SNR}
    \mathrm{SNR} :=\frac{ \left[\sum_{m=2}^{M} (m-1)\left(\frac{a_m - b_m}{k^{m-1}} \right) \right]^2 }{\sum_{m=2}^{M} (m-1)\left(\frac{a_m - b_m}{k^{m-1}} + b_m \right)}. 
\end{align}
\end{definition}
$\mathrm{SNR}$ is related to the following quantity
\begin{equation}
    \begin{aligned}
&\, \frac{[\lambda_2(\E\bA)]^{2}}{\lambda_1(\E\bA)} = \frac{[(n-k)\alpha - n \beta]^2}{k[(n - k)\alpha + n(k-1)\beta]}
    %=&\, \frac{ \left[ \left(\frac{n}{k} - 1 \right)\sum_{m=2}^{M}\binom{n/k-2}{m-2}  \frac{a_m - b_m}{\binom{n}{m-1}} - \sum_{m=2}^{M}\binom{n-2}{m-2}  \frac{b_m}{\binom{n}{m-1}} \right]^2 }{ \left(\frac{n}{k} - 1 \right)\sum_{m=2}^{M}\binom{n/k-2}{m-2}  \frac{a_m - b_m}{\binom{n}{m-1}} +(n-1) \sum_{m=2}^{M}\binom{n-2}{m-2}  \frac{b_m}{\binom{n}{m-1}} }\\
    = \frac{ \left[\sum_{m=2}^{M} (m-1)\left(\frac{a_m - b_m}{k^{m-1}} \right) \right]^2 }{\sum_{m=2}^{M} (m-1)\left(\frac{a_m - b_m}{k^{m-1}} + b_m \right)}(1+o(1)) \notag.
    \end{aligned}
\end{equation}
When $M=2$ and $k$ is fixed, \eqref{eqn:SNR} is equal to $\frac{(a - b)^2}{k[a + (k-1)b]}$, which corresponds to the $\mathrm{SNR}$ for the undirected graph in \cite{chin2015stochastic}, see also \cite[Section 6]{abbe2018community}.

\section{Spectral norm concentration}\label{sec:concentration}


The correctness of \Cref{alg:k_block_partition} and \Cref{alg:2_block_partition} relies on the concentration of the adjacency matrix of $H$. We include the following two concentration results for general random hypergraphs. The proofs are given in \Cref{sec:proof_concentration}.

\begin{theorem}\label{thm:concentration}
Let $H=\bigcup_{m=2}^M H_m$, where $H_m = ([n], E_m)$ is an Erd\H{o}s-R\'{e}nyi inhomogeneous hypergraph of order $m$ with $2\leq m\leq M$ with a probability tensor $\bcT^{(m)}$ such that  $\bcT^{(m)}_{[i_1,\dots, i_m]} = d_{[i_1,\dots,i_m]} / \binom{n}{m-1}$ and $d_m=\max d_{[i_1,\dots, i_m]}$.
Suppose for some constant $c>0$,
\begin{align}\label{eq:assumption_d}
    d:=\sum_{m=2}^M (m-1)\cdot d_m \geq c\log n\,.
\end{align}
 Then for any $K>0$, there exists a constant 
 $
    C= 512M(M-1)(K+6)\left[ 2 + (M-1)(1+K)/c \right]
 $
 such that with probability at least $1-2n^{-K}-2e^{-n}$, the adjacency matrix $\bA$ of $H$ satisfies  
\begin{align}\label{eq:concentrateA}
    \|\bA - \E \bA\|\leq C \sqrt{d}\,.
\end{align}
\end{theorem}

Taking $M=2$, \Cref{eq:concentrateA} is the result for graph case obtained in \cite{feige2005spectral,lei2015consistency}. Result for uniform hypergraph is obtained in \cite{lee2020robust}. Note that $d$ is a fixed constant in our community detection problem, thus \Cref{eq:assumption_d} does not hold and \Cref{thm:concentration} cannot be directly applied. However, we can still prove a concentration bound for a regularized version of $\bA$, following the strategy of the proof for Theorem \ref{thm:concentration}. 
\begin{definition}[Regularized matrix]\label{def:regularize}
Given any $n\times n$ matrix $\bA$ and an index set $\cI$, let $\bA_{\cI}$ be the $n\times n$ matrix obtained from $\bA$ by zeroing out the rows and columns not indexed by $\cI$. Namely,
\begin{equation}\label{eqn:restricted_adjacency_matrix}
     (\bA_{\cI})_{ij} = \mathbbm{1}_{ \{i, j\in \cI\} } \cdot \bA_{ij}\,. %\left\{
        %\begin{aligned}
            %&\bA_{ij} & \textnormal{if } i,j\in \cI, \\
            %&0\,, & \textnormal{otherwise}
        %\end{aligned}
    %\right.
\end{equation}
\end{definition}
Since every hyperedge of size $m$ containing $i$ is counted $(m-1)$ times in the $i$-th row sum of $\bA$, the $i$-th row sum of $\bA$ is given by
\begin{align*}
    \textnormal{row}(i) :=\sum_{j}  \bA_{ij} := \sum_{j} \mathbbm{1}_{\{ i \neq j\} } \sum_{m=2}^{M}\,\,\sum_{\substack{e\in E_m\\ \{i,j\}\subset e} }\bcT^{(m)}_{e} = \sum_{m=2}^{M}\,(m-1) \sum_{e\in E_m:\, i\in e} \bcT_e^{(m)}.
\end{align*}

\Cref{thm:regularization} is the concentration result for the regularized $\bA_{\cI}$, by zeroing out rows and columns corresponding to vertices with high row sums.

\begin{theorem}\label{thm:regularization}
%Let $H_m$ be Erd\H{o}s-R\'{e}nyi inhomogeneous hypergraph of order $m$ with $2\leq m\leq M$ with a probability tensor $\bcT^{(m)}= \frac{1}{\binom{n}{m-1} } (d_{i_1,\dots,i_m})$. 
Following all the notations in \Cref{thm:concentration},  for any constant $\tau >1$, define \[\cI=\{i\in [n]: \textnormal{row}(i)\leq \tau d\}.\] Let $\bA_{\cI}$ be the regularized version of $\bA$, as in \Cref{def:regularize}. Then for any $K>0$, there exists a constant $C_{\tau}$ depending on $M, \tau, K$, such that
$
     \| (\bA-\E \bA)_{\cI}\|\leq C_{\tau}\sqrt{d}
$
with probability at least $1-2(e/2)^{-n} -n^{-K}$.
\end{theorem}


\section{Spectral algorithms}\label{sec:spec_algorithm}

In this section, we present the algorithmic blocks that we use to construct our main partition algorithm (\Cref{alg:k_block_partition}): pre-processing (\Cref{alg:parameter_preprocessing}), first attempt at partition (\Cref{alg:k_block_spectral_partition}), correction of blemishes via majority rule (\Cref{alg:k_block_correction}), and merging (\Cref{alg:k_block_merging}), together with their proofs of correctness. We start by giving the algorithms. 

\begin{algorithm}
    \caption{ \textbf{Pre-processing}}\label{alg:parameter_preprocessing}
     \begin{algorithmic}[1]
        \Function{\textbf{Pre-processing}}{Parameters $a_m$, $b_m$ for all $m\in\{2, \cdots, M \}$}
	    \State For each subset $\mathcal{S} \subset \{2, \cdots, M\}$, let $H_{\mathcal{S}} = \bigcup_{m \in \mathcal{S}}H_m$ denote the restriction of the  non-uniform hypergraph $H$ on $\mathcal{S}$. Compute $\mathrm{SNR}$ of $H_{\mathcal{S}}$, denoted by
	    \[
	        \mathrm{SNR}_{\mathcal{S}}:= \frac{ \left[\sum_{m\in \mathcal{S}} (m-1)\left(\frac{a_m - b_m}{k^{m-1}} \right) \right]^2 }{\sum_{m\in \mathcal{S}} (m-1)\left(\frac{a_m - b_m}{k^{m-1}} + b_m \right)}\,.
	    \]
	    \State Among all the $\mathcal{S}$, find the subset $\mathcal{M}$ such that
	    \[
	        \mathcal{M}:= \arg\max_{\mathcal{S} \subset \{2, \cdots, M\}} \,\mathrm{SNR}_{\mathcal{S}}\,,
	    \]
	    with $\mathcal{M}_{\max}$ denoting its maximal element.
	    \State \Return{$\mathcal{M}$.}
	    \EndFunction
     \end{algorithmic}
\end{algorithm}

\subsection{Three or more blocks ($k\geq 3$)}

\begin{algorithm}
        \caption{\textbf{Spectral Partition}}\label{alg:k_block_spectral_partition}
	    \begin{algorithmic}[1]
	    \Function{\textbf{Spectral Partition}}{$\bB$, $k$, $a_m$, $b_m$ for $m\in \mathcal{M}$}
	    \State Randomly label vertices in $Y$ with $+1$ and $-1$ sign with equal probability, and partition $Y$ into $2$ disjoint subsets $Y_1$ and $Y_2$. 
	    \State Let $\bB_{1}$ (resp. $\bB_2$) denote the adjacency matrices with all vertices  in $Z \cup Y_1$, with rows indexed by $Z$ and columns indexed by  $Y_1$ (resp. $Y_2$). Pad $\bB_{1}$, $\bB_{2}$ with zeros to obtain the $n\times n$ matrices $\bA_{1}$ and $\bA_{2}$.
	    
	    \State Let $d:= \sum_{m\in \mathcal{M}}(m-1)a_m$. Zero out all the rows and columns of $\bA_{1}$ corresponding to vertices whose row sum is bigger than $20\mathcal{M}_{\max} d$, to obtain the matrix $(\bA_1)_{\cI_1}$.
	    
	    \State Find the space $\bU$, spanned by the first $k$ left singular vectors of $(\bA_1)_{\cI_1}$.
	    
	    \State  Randomly sample $s = 2k\log^2 n$ vertices from $Y_2$ without replacement. Denote the corresponding columns in $\bA_2$ by $\ba_{i_1},\dots, \ba_{i_s}$. For each $i\in \{i_1, \cdots, i_s\}$, project $\ba_{i}- \overline{\ba}$ onto $\bU$, where the elements in vector $\overline{\ba}\in \R^{n}$ is defined by $\overline{\ba}(j) = \mathbbm{1}_{j\in Z}\cdot (\overline{\alpha} + \overline{\beta})/2$
        %\begin{equation}
            %\tilde{\ba}(j) = \left\{
            %\begin{aligned}
                %&(\alpha + \beta)/2, & \textnormal{if } j\in Z \\
                %&0, & \textnormal{if } j\in Y
            %\end{aligned}
            %\right.
        %\end{equation}
	    \State For each projected vector $P_{\bU}(\ba_{i}- \overline{\ba})$, identify the top $n/(2k)$ coordinates in value as a set $U_i$.  Discard half of the $s$ sets $U_i$, those with the lowest blue hyperedge density in them.
	    
	    \State Sort the remaining sets according to blue hyperedge density and identify first $k$ distinct subsets $U^{\prime}_1, \cdots, U^{\prime}_k$ such that $|U^{\prime}_i \cap U^{\prime}_j| < \lceil(1 - \nu)n/k \rceil$ if $i\neq j$.
	    
	    \State \Return{ $U_{1}^{\prime}, \cdots, U_{k}^{\prime}$.}
	    \EndFunction
     \end{algorithmic} 
\end{algorithm}


\begin{algorithm}
    \caption{\textbf{Correction}}\label{alg:k_block_correction}
     \begin{algorithmic}[1]
        \Function{\textbf{Correction}}{A collection of subsets $U_{1}^{\prime}, \cdots, U_{k}^{\prime} \subset Z$ and a graph on $Z$}
	    \State For every $u \in Z$, if $i\in \{1, 2, \cdots, k\}$ is such that $u$ has more neighbors in $U_{i}^{\prime}$ than in any other $U_j^{\prime}$, for $j \neq i$, then add $u$ to $\widehat{U}_{i}$. Break ties arbitrarily.
	    
	    \State \Return{$\widehat{U}_{1}, \cdots, \widehat{U}_{k}$.}
	    \EndFunction
     \end{algorithmic}
\end{algorithm}

\begin{algorithm}
    \caption{\textbf{Merging}}\label{alg:k_block_merging}
     \begin{algorithmic}[1]
        \Function{\textbf{Merging}}{A partition $\widehat{U}_{1}, \cdots, \widehat{U}_{k}$ of $(Z\cap V_{1})\cup (Z\cap V_{2}) \cup \cdots \cup (Z \cap V_{k})$ and a graph between vertices $Y$ and $Z$.}
	    \State For all $u \in Y$, add $u$ to  $\widehat{U_i}$ if the number of hyperedges, which contains $u$ with the remaining vertices located in vertex set $\widehat{U}_i$, is at least $\mu_{\rm{M}}$. Label the conflicts arbitrarily.
	    \State \Return{the label classes as the clusters $\widehat{V}_{1}, \cdots, \widehat{V}_{k}$.}
	    \EndFunction
     \end{algorithmic}
\end{algorithm}


The proof of \Cref{thm:mainresult_k} is structured as follows.
\begin{lemma}\label{lem:spectralclustering_accuracy_k}
  Under the assumptions of \Cref{thm:mainresult_k}, \Cref{alg:k_block_spectral_partition} outputs a $\nu$-correct partition $U_1^{\prime}, \cdots, U_k^{\prime}$ of $Z = (Z\cap V_{1}) \cup \cdots \cup (Z\cap V_{k})$ with probability at least $1 - O(n^{-2})$.
\end{lemma}

\begin{lemma}\label{lem:correction_accuracy_k}
 Under the assumptions of \Cref{thm:mainresult_k}, 
    given any $\nu$-correct partition $U_1^{\prime}, \cdots, U_k^{\prime}$ of $Z = (Z\cap V_{1}) \cup \cdots \cup (Z\cap V_{k})$ and the red hypergraph over $Z$,  \Cref{alg:k_block_correction} computes a $\gamma_{\rm{C}}$-correct partition $\widehat{U}_1, \cdots, \widehat{U}_k$ with probability $1 -O(e^{-n\rho})$, while $\gamma_{\rm{C}} = \max\{\nu,\, 1 - k\rho\}$ with
$
    \rho:= k\exp\left( -C_{\mathcal M} \cdot \mathrm{SNR}_{\mathcal{M}} \right)
$
where $\mathcal{M}$ is obtained from \Cref{alg:parameter_preprocessing}, and $\mathrm{SNR}_{\mathcal{M}}$ and $C_{\mathcal M}$ are defined in \Cref{eqn:SNR_CM}.
\end{lemma}

\begin{lemma}\label{lem:merging_accuracy_k}
    Given any $\nu$-correct partition $\widehat{U}_1, \cdots, \widehat{U}_k$ of $Z = (Z\cap V_{1}) \cup \cdots \cup (Z\cap V_{k})$ and the blue hypergraph between $Y$ and $Z$, with probability $1 -O(e^{-n\rho})$,  \Cref{alg:k_block_merging} outputs a $\gamma$-correct partition $\widehat{V}_1, \cdots, \widehat{V}_{k}$ of $V_1 \cup V_2 \cup \cdots \cup V_k$, while $\gamma =\max\{\nu,\, 1 - k\rho\}$.
\end{lemma}


\subsection{The binary case ($k = 2$)}
The spectral partition step is given in \Cref{alg:2_block_spectral_partition}, and the correction step is given in  \Cref{alg:2_block_correction}.
\begin{algorithm}[ht]
	\caption{\textbf{Spectral Partition}}
	\label{alg:2_block_spectral_partition}
    \begin{algorithmic}[1]
	    \Function{\textbf{Spectral Partition}}{$\bA$, $\bcT$}
	    \State Zero out all the rows and the columns of $\bA$ corresponding to vertices with row sums greater than $20\mathcal{M}_{\max}d$, to obtain the regularized matrix $\bA_{\cI}$.
	    \State Find the subspace $\bU$, which is spanned by the eigenvectors corresponding to the largest two eigenvalues of $\bA_{\cI}$. %\Comment{Here, we assume the $am(\lambda_{1}) = am(\lambda_{2}) = 1$}
	    \State Compute $P_{\bU}\mathbf{1}_n$, the projection of all-ones vector onto $\bU$.
	    \State Let $\bv$ be the unit vector in $\bU$ perpendicular to $P_{\bU}\mathbf{1}_n$.
	    \State Sort the vertices according to their values in $\bv$. Let $V_{1}^{\prime} \subset V$ be the corresponding top $n/2$ vertices, and $V_{2}^{\prime} \subset V$ be the remaining $n/2$ vertices.
	    \State \Return{$V_{1}^{\prime}$, $V_{2}^{\prime}$.}
	    \EndFunction
\end{algorithmic}  
\end{algorithm} 
\begin{lemma}\label{lem:spectralclustering_accuracy_2}
Under the conditions of \Cref{thm:mainresult_2}, the \Cref{alg:2_block_spectral_partition} outputs a $\nu$-correct partition $V_1^{\prime}, V_2^{\prime}$ of $V = V_1 \cup V_2$ with probability at least $1 - O(n^{-2})$.
\end{lemma}

\begin{algorithm}[ht]
	\caption{\textbf{Correction}}
	\label{alg:2_block_correction}
    \begin{algorithmic}[1]   
	    \Function{\textbf{Correction}}{$\bcT$, partition $V_{1}^{\prime}$, $V_{2}^{\prime}$, and a blue graph on $V_{1}^{\prime} \cup V_{2}^{\prime}$}
	    \State For any $v \in V_{1}^{\prime}$, label $v$ ``bad" if the number of blue neighbors of $v$ in $V_{2}^{\prime}$ is at least $\frac{a + b}{4}$ and ``good" otherwise.
	    \State Do the same for $v\in V_{2}^{\prime}$.
	    \State If $v\in V_{i}^{\prime}$ is good, assign it to $\widehat{V}_{i}$, otherwise $\widehat{V}_{3 - i}$.
	    \State \Return{$\widehat{V}_{1}$, $\widehat{V}_{2}$.}
	    \EndFunction
     \end{algorithmic}  
\end{algorithm}
\begin{lemma}\label{lem:correction_accuracy_2}
    Given any $\nu$-correct partition $V_1^{\prime}, V_2^{\prime}$ of $V = V_1 \cup V_2$, with probability at least $1 -O(e^{-n\rho})$, the \Cref{alg:2_block_correction} computes a $\gamma$-correct partition $\widehat{V}_1, \widehat{V}_2$ with $\gamma = \{\nu,\, 1 - 2\rho\}$ and $\rho = 2\exp( -C_{\mathcal M} \cdot \mathrm{SNR}_{\mathcal{M}})$, where $\mathrm{SNR}_{\mathcal{M}}$ and $C_{\mathcal M}$ are defined in \Cref{eqn:SNR_CM_2}.
\end{lemma}

% \section{Numerical Experiments}

% \section{Conclusion}
    
\section{Algorithm's correctness}\label{sec:analysis_k}
In this section, we will show the correctness of \Cref{alg:k_block_partition}. We first introduce some definitions.


\subsubsection*{Vertex set splitting and adjacency matrix}
In \Cref{alg:k_block_partition}, we first randomly partition the vertex set $V$ into two disjoint subsets $Z$ and $Y$ by assigning $+1$ and $-1$ to each vertex independently with equal probability. Let $\bB \in \R^{|Z|\times |Y|}$ denote the submatrix of $\bA$, while $\bA$ was defined in \Cref{eqn:adjacency_matrix_entry}, where rows and columns of $\bB$ correspond to vertices in $Z$ and $Y$ respectively. Let $n_i$ denote the number of vertices in $Z\cap V_i$, where $V_i$ denotes the true partition with $|V_i| = n/k$ for all $i\in [k]$, then $n_i$ can be written as a sum of independent Bernoulli random variables, i.e.,
\begin{align}\label{eqn:dimension_ZcapVi}
    n_i = |Z \cap V_i| = \sum_{v\in V_i} \mathbbm{1}_{ \{v\in Z\} }\, ,
\end{align}
and $|Y\cap V_i| = |V_i| - |Z\cap V_i| = n/k- n_i$ for each $i \in [k]$.
\begin{definition}
The splitting $V = Z \cup Y$ is \textit{perfect} if $|Z\cap V_i| = |Y\cap V_i| = n/(2k)$ for all $i\in [k]$. And the splitting $Y = Y_1 \cup Y_2$ is \textit{perfect} if $|Y_1\cap V_i| = |Y_2\cap V_i|  = n/(4k)$ for all $i\in [k]$.
\end{definition}
However, the splitting will actually be \textit{imperfect} in most cases, since the size of $Z$ and $Y$ would not be exactly the same under the independence assumption. 
The random matrix $\bB$ is parameterized by $\{\bcT^{(m)}\}_{m \in \mathcal{M}}$ and $\{n_i\}_{i=1}^{k}$. If we take expectation over $\{\bcT^{(m)}\}_{m \in \mathcal{M}}$ given the block size information $\{n_i\}_{i=1}^{k}$, then it gives rise to the expectation of the \textit{imperfect} splitting, denoted by $\widetilde{\bB}$,
\begin{equation*}
    \widetilde{\bB} := \begin{bmatrix}
     \alpha \bJ_{n_1 \times ( \frac{n}{k} - n_1)} & \beta \bJ_{n_1 \times (\frac{n}{k} - n_2)}  & \dots & \beta \bJ_{n_1 \times (\frac{n}{k}- n_k)}  \\
    \beta \bJ_{n_2 \times (\frac{n}{k}- n_1)} & \alpha \bJ_{n_2 \times (\frac{n}{k}- n_2)} &\dots & \beta \bJ_{n_2 \times (\frac{n}{k}- n_k)}  \\
    \vdots & \vdots &  \ddots & \vdots \\
    \beta  \bJ_{n_k \times (\frac{n}{k}- n_1)}  & \beta  \bJ_{n_k \times (\frac{n}{k}- n_2)}  & \dots & \alpha \bJ_{n_k \times (\frac{n}{k}- n_k)}
  \end{bmatrix}\,, %=\E_{ \{\bcT^{(m)}\}_{m \in \mathcal{M}} } \left[ \bB \Big| \{n_i\}_{i=1}^{k} \right]
\end{equation*}
where $\alpha$, $\beta$ are defined in \Cref{eqn:alpha_beta}. In the \textit{perfect} splitting case, the dimension of each block is $n/(2k)\times n/(2k)$ since $\E n_i = n/(2k)$ for all $i \in [k]$, and the expectation matrix $\overline{\bB}$ can be written as
\begin{equation*}
    \overline{\bB} := \begin{bmatrix}
     \alpha \bJ_{\frac{n}{2k}} & \beta \bJ_{\frac{n}{2k}} & \dots & \beta \bJ_{\frac{n}{2k}}  \\
    \beta \bJ_{\frac{n}{2k}} & \alpha \bJ_{\frac{n}{2k}} & \dots & \beta \bJ_{\frac{n}{2k}}  \\
    \vdots & \vdots & \ddots & \vdots \\
    \beta \bJ_{\frac{n}{2k}} & \beta \bJ_{\frac{n}{2k}} &\dots & \alpha \bJ_{\frac{n}{2k}}
  \end{bmatrix}\,.
\end{equation*}

In \Cref{alg:k_block_spectral_partition}, $Y_1$ is a random subset of $Y$ obtained by selecting each element with probability $1/2$ independently, and $Y_2 = Y\setminus Y_1$. Let $n^{\prime}_{i}$ denote the number of vertices in $Y_1\cap V_i$, then $n^{\prime}_{i}$ can be written as a sum of independent Bernoulli random variables, 
\begin{align}\label{eqn:dimension_Y1capVi}
    n^{\prime}_{i} = |Y_1 \cap V_i| = \sum_{v\in V_i} \mathbf{1}_{ \{v\in Y_1\} }\,,
\end{align}
and $|Y_2\cap V_i| = |V_i| - |Z\cap V_i| - |Y_1\cap V_i| =  n/k - n_i - n^{\prime}_{i}$ for all $i \in [k]$. 


% Similarly,  $\E n^{\prime}_i = (n/4k)$, which means that this splitting is expected to be perfect. However, we still obtain \textit{imperfect} splitting in most cases due to the independence assumption.

% \begin{remark}
% In \cite{chin2015stochastic}, the authors assumed the \textit{perfect} splitting throughout the proof, which is not necessarily true. Our work provides a rigorous analysis of the algorithm provided in \cite{chin2015stochastic}, without the \textit{perfect} splitting assumption.
% \end{remark}


\subsubsection*{Induced sub-hypergraph}\label{subsubsec:induced_sub_hypergraph}
% For the non-uniform hpyergraph $H = (V, E)$ with $H = \cup_{m \in \mathcal{M}} H_m$ and $E = \cup_{m \in \mathcal{M}}E_m$,  consider the vertex subset $(Y_1 \cup Z) \subset V$, which leads our attention to the corresponding induced sub-hypergraph. 

\begin{definition}[Induced Sub-hypergraph]\label{def:induced sub-hypergraph}
Let $H = (V, E)$ be a non-uniform random hypergraph and $S\subset V$ be any subset of the vertices of $H$. Then the \textit{induced sub-hypergraph} $H[S]$ is the hypergraph whose vertex set is $S$ and whose hyperedge set $E[S]$ consists of all of the edges in $E$ that have all endpoints located in $S$.
\end{definition}

% \begin{remark}
% Note that $H$ can be considered as a collection of $m$-uniform hypergraphs for varying $m$, a.k.a., $H = \cup_{m \in \mathcal{M}} H_m$. The decomposition works for $H[S]$ as well, a.k.a. $H[S] = \cup_{m \in \mathcal{M}} H_m[S]$ and $E[S] = \cup_{m \in \mathcal{M}} E_m[S]$
% \end{remark}

Let $H[Y_1 \cup Z]$(resp. $H[Y_2 \cup Z]$)  denote the induced sub-hypergraph on vertex set $Y_1\cup Z$ (resp. $Y_2 \cup Z$), and $\bB_1 \in \R^{|Z|\times |Y_1|}$ (resp. $\bB_2 \in \R^{|Z|\times |Y_2|}$) denote the adjacency matrices corresponding to the sub-hypergraphs, where rows and columns of $\bB_1$ (resp. $\bB_2$) are corresponding to elements in $Z$ and $Y_1$ (resp., $Z$ and $Y_2$). Therefore,  $\bB_1$ and $\bB_2$ are parameterized by $\{\bcT^{(m)}\}_{m \in \mathcal{M}}$, $\{n_i\}_{i=1}^{k}$ and $\{n^{\prime}_{i}\}_{i=1}^{k}$, and the entries in $\bB_1$ are independent of the entries in $\bB_2$, due to the independence of hyperedges. If we take expectation over $\{\bcT^{(m)}\}_{m \in \mathcal{M}}$ conditioning on $\{n_i\}_{i=1}^{k}$ and $\{n^{\prime}_{i}\}_{i=1}^{k}$, then it gives rise to the expectation of the \textit{imperfect} splitting, denoted by $\widetilde{\bB}_1$,
\begin{equation}\label{eqn:tilde_B1}
    \widetilde{\bB}_1 := \begin{bmatrix}
     \widetilde{\alpha}_{11} \bJ_{n_1 \times n^{\prime}_{1}} & \dots & \widetilde{\beta}_{1k} \bJ_{n_1 \times n^{\prime}_{k}}  \\
    \vdots & \ddots & \vdots \\
    \widetilde{\beta}_{k1}  \bJ_{n_k \times n^{\prime}_{1}}  & \dots & \widetilde{\alpha}_{kk} \bJ_{n_k \times n^{\prime}_{k}} % = \E_{ \{\bcT^{(m)}\}_{m \in \mathcal{M}} } \left[ \bB \Big| \{n_i\}_{i=1}^{k}, \{n^{\prime}_{i}\}_{i=1}^{k} \right]
  \end{bmatrix}\,,
\end{equation}
where
\begin{subequations}
\begin{align}
    \widetilde{\alpha}_{ii} :=&\, \sum_{m \in \mathcal{M}}\left\{ \binom{n_i + n_i^{\prime} - 2}{m-2} \frac{a_m - b_m}{\binom{n}{m-1}} + \binom{\sum_{l=1}^{k}(n_l + n_l^{\prime}) - 2}{m-2} \frac{b_m}{\binom{n}{m-1}} \right\}\,,\label{eqn:tilde_alpha} \\
    \widetilde{\beta}_{ij} :=&\, \sum_{m \in \mathcal{M}} \binom{\sum_{l=1}^{k}(n_l + n_l^{\prime}) - 2}{m-2}\frac{b_m}{\binom{n}{m-1}}\,,\quad i\neq j, \, i, j\in [k]\,. \label{eqn:tilde_beta}
\end{align}
\end{subequations}
The expectation of the \textit{perfect} splitting, denoted by $\overline{\bB}_1$, can be written as
\begin{equation}\label{eqn:bar_B1}
    \overline{\bB}_1 := \begin{bmatrix}
     \overline{\alpha} \bJ_{\frac{n}{2k} \times \frac{n}{4k}} & \overline{\beta} \bJ_{\frac{n}{2k} \times \frac{n}{4k}} & \dots & \overline{\beta} \bJ_{\frac{n}{2k} \times \frac{n}{4k}} \\
    \overline{\beta} \bJ_{\frac{n}{2k} \times \frac{n}{4k}} & \overline{\alpha} \bJ_{\frac{n}{2k} \times \frac{n}{4k}} & \dots & \overline{\beta} \bJ_{\frac{n}{2k} \times \frac{n}{4k}}  \\
    \vdots & \vdots & \ddots & \vdots \\
    \overline{\beta} \bJ_{\frac{n}{2k} \times \frac{n}{4k}} & \overline{\beta} \bJ_{\frac{n}{2k} \times \frac{n}{4k}} &\dots & \overline{\alpha} \bJ_{\frac{n}{2k} \times \frac{n}{4k}}
  \end{bmatrix}\,,
\end{equation}
where
\begin{align}\label{eqn:bar_alphabeta}
    \overline{\alpha} := \sum_{m \in \mathcal{M}}\left\{ \binom{ \frac{3n}{4k} - 2}{m-2} \frac{a_m - b_m}{\binom{n}{m-1}} + \binom{\frac{3n}{4} - 2}{m-2}\frac{b_m}{\binom{n}{m-1}} \right\}\,,\quad  \overline{\beta} := \sum_{m \in \mathcal{M}} \binom{\frac{3n}{4} - 2}{m-2}\frac{b_m}{\binom{n}{m-1}}\,.
\end{align}
The matrices $\widetilde{\bB}_2, \overline{\bB}_2$ can be defined similarly, since the dimensions of $|Y_2\cap V_i|$ are also determined by $n_i$ and $n^{\prime}_{i}$. Obviously, $\overline{\bB}_2 = \overline{\bB}_1$ since $\E n^{\prime}_{i} = \E (n/k- n_{i} - n_i^{\prime}) = n/(4k)$ for all $i\in [k]$.

\subsubsection*{Fixing Dimensions}
The dimensions of $\widetilde{\bB}_1$ and $\widetilde{\bB}_2$, as well as blocks they consist of, are not deterministic---since $n_i$ and $n^{\prime}_i$, defined in \Cref{eqn:dimension_ZcapVi} and \Cref{eqn:dimension_Y1capVi} respectively, are sums of independent random variables. As such we cannot directly compare them. 
In order to overcome this difficulty, we embed $\bB_1$ and $\bB_2$ into the following $n\times n$ matrices:
\begin{equation}\label{eqn:A1A2}
    \bA_1:= \begin{bmatrix}
     \mathbf{0}_{|Z|\times |Z|} & \bB_1 & \mathbf{0}_{|Z|\times |Y_2|} \\
     \mathbf{0}_{|Y|\times |Z|} & \mathbf{0}_{|Y|\times |Y_1|}  & \mathbf{0}_{|Y|\times |Y_2|} 
  \end{bmatrix}\, ,\quad 
  \bA_2:= \begin{bmatrix}
     \mathbf{0}_{|Z|\times |Z|} & \mathbf{0}_{|Z|\times |Y_1|} & \bB_2 \\
     \mathbf{0}_{|Y|\times |Z|} & \mathbf{0}_{|Y|\times |Y_1|}  & \mathbf{0}_{|Y|\times |Y_2|} 
    \end{bmatrix}\,.
\end{equation}
Note that $\bA_1$ and $\bA_2$ have the same size. Also by definition, the entries in $\bA_1$ are independent of the entries in $\bA_2$. If we take expectation over $\{\cT^{(m)}\}_{m\in \mathcal{M}}$ conditioning on $\{n_i\}_{i=1}^{k}$ and $\{n^{\prime}_{i}\}_{i=1}^{k}$, then we obtain the expectation matrices of the \textit{imperfect} splitting, denoted by $\widetilde{\bA}_1$(resp. $\widetilde{\bA}_2$), written as
\begin{equation}\label{eqn:tilde_A1A2}
    \widetilde{\bA}_1:= \begin{bmatrix}
     \mathbf{0}_{|Z|\times |Z|} & \widetilde{\bB}_1 & \mathbf{0}_{|Z|\times |Y_2|} \\
     \mathbf{0}_{|Y|\times |Z|} & \mathbf{0}_{|Y|\times |Y_1|}  & \mathbf{0}_{|Y|\times |Y_2|} 
  \end{bmatrix}\, ,\quad 
  \widetilde{\bA}_2:= \begin{bmatrix}
     \mathbf{0}_{|Z|\times |Z|} & \mathbf{0}_{|Z|\times |Y_1|} & \widetilde{\bB}_2 \\
     \mathbf{0}_{|Y|\times |Z|} & \mathbf{0}_{|Y|\times |Y_1|}  & \mathbf{0}_{|Y|\times |Y_2|} 
    \end{bmatrix}\,.
\end{equation}
The expectation matrix of the \textit{perfect} splitting, denoted by $\overline{\bA}_1$(resp. $\overline{\bA}_2$), can be written as
\begin{equation}\label{eqn:bar_A1A2}
    \overline{\bA}_1:= \begin{bmatrix}
     \mathbf{0}_{ \frac{n}{2} \times \frac{n}{2}} & \overline{\bB}_1 & \mathbf{0}_{\frac{n}{2}\times \frac{n}{4}} \\
     \mathbf{0}_{\frac{n}{2}\times \frac{n}{2}} & \mathbf{0}_{\frac{n}{2}\times \frac{n}{4}}  & \mathbf{0}_{\frac{n}{2}\times \frac{n}{4}} 
  \end{bmatrix}\, ,\quad 
  \overline{\bA}_2:= \begin{bmatrix}
     \mathbf{0}_{ \frac{n}{2} \times \frac{n}{2}} & \mathbf{0}_{\frac{n}{2}\times \frac{n}{4}} & \overline{\bB}_2\\
     \mathbf{0}_{\frac{n}{2}\times \frac{n}{2}} & \mathbf{0}_{\frac{n}{2}\times \frac{n}{4}}  & \mathbf{0}_{\frac{n}{2}\times \frac{n}{4}} 
  \end{bmatrix}\,.
\end{equation}
Obviously, $\widetilde{\bA}_i$ and $\widetilde{\bB}_i$(resp. $\overline{\bA}_i$ and $\overline{\bB}_i$) have the same non-zero singular values for $i = 1, 2$. In the remaining of this section, we will deal with $\widetilde{\bA}_i$ and $\overline{\bA}_i$ instead of $\widetilde{\bB}_i$ and $\overline{\bB}_i$ for $i = 1, 2$.
% since the former have the same size, and can be compared.

\subsection{Spectral Partition: Proof of Lemma \ref{lem:spectralclustering_accuracy_k}}

\subsubsection{Proof Outline}
% In \Cref{alg:k_block_spectral_partition}, the matrix $(\bA_1)_{\cI_1}$ was obtained by zeroing out rows and columns in $\bA_1$ corresponding to vertices with high row sums.
Recall that $\bA_1$ is defined as the adjacency matrix of the induced sub-hypergraph $H[Y_1\cup Z]$ in \Cref{subsubsec:induced_sub_hypergraph}. Consequently, the index set should contain information only from $H[Y_1\cup Z]$. Define the index sets
\begin{align*}
 \cI = \left\{i\in [n]: \mathrm{row}(i)\leq 20\mathcal{M}_{\max}d \right\}, \quad 
    \cI_{1} = \left\{i\in [n]: \mathrm{row}(i)\Big|_{Y_1 \cup Z} \leq 20\mathcal{M}_{\max}d \right\}\,,
\end{align*}
where $d = \sum_{m \in \mathcal{M}}(m-1)a_m$,  and $\mathrm{row}(i)\big|_{Y_1 \cup Z}$ is the row sum of $i$ on $H[Y_1\cup Z]$. We say $\mathrm{row}(i)\big|_{Y_1 \cup Z}=0$ if $i\not\in Y_1\cup Z$, and for vertex $i\in Y_1\cup Z$,
\begin{align*}
    \mathrm{row}(i)\Big|_{Y_1 \cup Z} := \sum_{j=1}^{n} \sum_{m\in \mathcal{M}}\,\,\sum_{ \substack{e\in E_m[Y_1\cup Z]\\ \{i, j\}\subset e} } \,\, \bcT_e^{(m)} = \sum_{m\in \mathcal{M}}(m-1)\sum_{ \substack{e\in E_m[Y_1\cup Z]\\ \{i, j\}\subset e} } \,\, \bcT_e^{(m)}\,.
\end{align*}
As a result, the matrix $(\bA_1)_{\cI_1}$ is obtained by restricting $\bA_1$ on index set $\cI_1$. The next $4$ steps guarantee that \Cref{alg:k_block_spectral_partition} outputs a $\nu$-correct partition.
\begin{enumerate}[label=(\roman*)]
    \item Find the singular subspace $\bU$, which is spanned by the first $k$ left singular vectors of $(\bA_1)_{\cI_1}$.
    \item Randomly pick $s = 2k \log^2 n$ vertices from $Y_2$ and denote the corresponding columns in $\bA_2$ by $\ba_{i_1},\dots, \ba_{i_s}$. Project each vector $\ba_i - \overline{\ba}$ onto the singular subspace $\bU$, with $\overline{\ba}\in \R^{n}$ defined by $\overline{\ba}(j) = \mathbbm{1}_{j\in Z} \cdot (\overline{\alpha} + \overline{\beta})/2$, where $\overline{\alpha}$, $\overline{\beta}$ were defined in  \Cref{eqn:bar_alphabeta}.
    %\begin{equation*}
        %\widetilde{\ba}(j) = \left\{
        %\begin{aligned}
            %&(\overline{\alpha} + \overline{\beta})/2, & \textnormal{if } j\in Z \\
            %&0, & \textnormal{if } j\in Y
        %\end{aligned}
    %\right.\,, \quad \overline{\alpha}, \overline{\beta} \textnormal{ defined in  \Cref{eqn:bar_alphabeta}}\,.
    %\end{equation*}
    \item For each projected vector $P_{\bU}(\ba_i - \overline{\ba})$, identify the top $n/(2k)$ coordinates in value and place the corresponding vertices into a set $U^{\prime}_i$. Discard half of the obtained $s$ subsets, those with the lowest \emph{blue} edge densities. 
    \item Sort the remaining sets according to blue hyperedge density and identify $k$ distinct subsets $U^{\prime}_1, \cdots, U^{\prime}_k$ such that $|U^{\prime}_i \cap U^{\prime}_j| < \lceil(1 - \nu)n/k\rceil$ if $i\neq j$.
\end{enumerate}
Based on the $4$ steps above in \Cref{alg:k_block_spectral_partition}, the proof of Lemma \ref{lem:spectralclustering_accuracy_k} is structured in $4$ parts.
\begin{enumerate}[label=(\roman*)]
    \item Let $\widetilde{\bU}$ denote the subspace spanned by the first $k$ left singular vectors of $\widetilde{\bA}_1$, as defined in \Cref{eqn:tilde_A1A2}. \Cref{subsubsec:bound_subspace_angle_k} shows that the subspace angle between $\bU$ and $\widetilde{\bU}$ is of order $o(1)$.
    
    \item The vector $\widetilde{\bdelta}_i$, defined in \Cref{eqn:tilde_delta},
    reflects the underlying true partition $Z\cap V_{k(i)}$ for each $i\in [s]$, where $k(i)$ denotes the membership of vertex $i$. \Cref{subsubsec:bound_the_projection_error} shows that $\widetilde{\bdelta}_i$ can be approximately recovered by the projected vector $P_{\bU}(\ba_{i} - \overline{\ba})$, since the projection error $\|P_{\bU}(\ba_{i} - \overline{\ba}) - \widetilde{\bdelta}_{i}\|_2$ is of order $O(1/\sqrt{n})$.
    
    \item \Cref{subsubsec:discard_half_blue_accuracy} indicates that the coincidence ratio between the remaining sets and the true partition is at least $\nu$, after discarding half of the sets with lowest blue hyperedge densities.
    
    \item \Cref{lem:sample_k_distinct} proves that we can find $k$ distinct subsets $U^{\prime}_i$ within $k\log^2 n$ trials \textit{w.h.p.}
\end{enumerate}

\subsubsection{Bounding the angle between $\bU$ and $\widetilde{\bU}$.}\label{subsubsec:bound_subspace_angle_k} Recall that the angle between subspaces $\bU$ and $\widetilde{\bU}$ is defined as $\sin \angle (\bU, \widetilde{\bU}) := \|P_{\bU} - P_{\widetilde{\bU}}\|$. A natural idea is to apply Wedin's $\sin \Theta$ Theorem  (\Cref{lem:Wedin_sin}).
% To bound the angle, we should focus on two things here: (i) the spectral gap $[\sigma_k(\widetilde{\bA}_1) - \sigma_{k+1}(\widetilde{\bA}_1)]$ and (ii) the deviation $\|(\bA_1)_{\cI_1} - \widetilde{\bA}_1\|$. 
% We study the spectral gap first. Unfortunately, it is difficult to find the closed-form expression for $\sigma_k(\widetilde{\bA}_1)$, since the dimensions of each block inside $\widetilde{\bA}_1$, as defined in \Cref{eqn:tilde_A1A2}, are not fixed. However, the singular values of $\overline{\bA}_1$, denoted by $\sigma_i(\overline{\bA}_1)$, are easy to compute---see \Cref{eqn:bar_A1A2}.


Lemma \ref{lem:singular_value_approximation} indicates that the difference  of $\sigma_i(\widetilde{\bA}_1)$ and $\sigma_i(\overline{\bA}_1)$ is relatively small, compared to $\sigma_i(\overline{\bA}_1)$. 
\begin{lemma}\label{lem:singular_value_approximation}
Let $\sigma_i(\overline{\bA}_1)$(resp.  $\sigma_i(\widetilde{\bA}_1)$) denote the singular values of $\overline{\bA}_1$ (resp.  $\widetilde{\bA}_1$) for all $i\in [k]$, where the matrices $\overline{\bA}_1$ and $\widetilde{\bA}_1$ are defined in \Cref{eqn:bar_A1A2} and \Cref{eqn:tilde_A1A2} respectively. Then
\begin{align*}
        \sigma_1(\overline{\bA}_1) =&\,\frac{n\left[ \overline{\alpha} + (k-1)\overline{\beta} \right] }{2\sqrt{2}k}= \frac{n}{2\sqrt{2}k}\sum_{m\in \mathcal{M}} \left[ \binom{\frac{3n}{4k} - 2}{m-2} \frac{a_m - b_m}{ \binom{n}{m-1}} + k\binom{\frac{3n}{4} - 2}{m-2}\frac{b_m}{\binom{n}{m-1}} \right] \,,&\\
        \sigma_i(\overline{\bA}_1) =&\, \frac{n( \overline{\alpha} - \overline{\beta})}{2\sqrt{2}k} = \frac{n}{2\sqrt{2}k}\sum_{m\in \mathcal{M}} \binom{\frac{3n}{4k} - 2}{m-2} \frac{a_m - b_m}{ \binom{n}{m-1}}\,, 2 \leq i \leq k\,,&\\
        \sigma_i(\overline{\bA}_1) =&\, 0\,, k+1 \leq i \leq n\,.&
\end{align*}
with $\overline{\alpha}$, $\overline{\beta}$ defined in \Cref{eqn:bar_alphabeta}. Moreover, with probability at least $1 - 2k\exp(-k\log^2(n))$,
\begin{align*}
    \frac{|\sigma_i(\overline{\bA}_1) - \sigma_i(\widetilde{\bA}_1)|}{\sigma_i(\overline{\bA}_1)} = O\left(n^{-1/4}\log^{1/2}(n)\right).
\end{align*}
\end{lemma}
Therefore, with \Cref{lem:singular_value_approximation}, we can write
$
    \sigma_i(\widetilde{\bA}_1) = \sigma_i(\overline{\bA}_1)(1 + o(1)).
$
% We then focus on the deviation $\|(\bA_1)_{\cI_1} - \widetilde{\bA}_1\|$.
Define $\bE_1 := \bA_1 - \widetilde{\bA}_1$ and its restriction on $\cI_1$ as 
\begin{equation}\label{eqn:E1I_k}
    (\bE_1)_{\cI_1} := (\bA_1 - \widetilde{\bA}_1)_{\cI_1} = (\bA_1)_{\cI_1} - (\widetilde{\bA}_1)_{\cI_1}\,,
\end{equation}
as well as $\bDelta_{1}:= (\widetilde{\bA}_1)_{\cI_1} - \widetilde{\bA}_1$. Then $(\bA_1)_{\cI_1} - \widetilde{\bA}_1$ is decomposed as
\begin{align*}
    (\bA_1)_{\cI_1} - \widetilde{\bA}_1 = [(\bA_1)_{\cI_1} - (\widetilde{\bA}_1)_{\cI_1}] + [(\widetilde{\bA}_1)_{\cI_1} -  \widetilde{\bA}_1 ] = (\bE_1)_{\cI_1} + \bDelta_{1}\,.
\end{align*}
\Cref{lem:high_degree_vertices_k} shows that the number of high degree vertices is relatively small with high probability. 
\begin{lemma}\label{lem:high_degree_vertices_k}
    Let $d = \sum_{m\in \mathcal{M}} (m-1)a_m$, where $\mathcal{M}$ is obtained from \Cref{alg:parameter_preprocessing}. There exists a constant $C_1$ such that if $d \geq C_1$, then with probability at least $1 - \exp\left(- d^{-2}n/\mathcal{M}_{\max} \right)$, no more than $d^{-3}n$ vertices have row sums greater than $20\mathcal{M}_{\max}d$.
\end{lemma} 
Consequently, \Cref{cor:norm_of_high_degree_vertices_k_block} indicates $\|\bDelta_1\| \leq \sqrt{d}$ with high probability.
\begin{corollary}\label{cor:norm_of_high_degree_vertices_k_block}
Assume $d \geq \max\{C_1, \sqrt{2}\}$, where $C_1$ is the constant in \Cref{lem:high_degree_vertices_k}, then $\|\bDelta_1\| \leq \sqrt{d}$ \text{w.h.p}.
    \end{corollary}
    \begin{proof}
   Note that $n - |\cI| \leq d^{-3}n$ and $\cI \subset \cI_1$, then $n - |\cI_1| \leq d^{-3}n$.
    From Lemma \ref{lem:high_degree_vertices_k}, there are at most $d^{-3} n$ vertices with row sum greater than $20\mathcal{M}_{\max}d$ in the adjacency matrix $\bA_1$, then the matrix $\bDelta_1 = (\widetilde{\bA}_1)_{\cI_1} - \widetilde{\bA}_1$ has at most $2d^{-3} n^{2}$ non-zero entries. As defined in \Cref{eqn:tilde_A1A2}, every entry in $\widetilde{\bA}_1$ is bounded by $\alpha$, then,
    \begin{subequations}
        \begin{align*}
           \|\bDelta_1\| \leq &\,  \|\bDelta_1 \|_{\frob} = \|(\widetilde{\bA}_1)_{\cI_1} - \widetilde{\bA}_1\|_{\frob} \\
           \leq &\, \sqrt{2 d^{-3}n^{2}} \, \alpha = \sqrt{2 d^{-3}} n \sum_{m\in \mathcal{M}} \left[ \binom{ \frac{n}{k}-2}{m-2} \frac{a_m - b_m}{\binom{n}{m-1}}  + {n - 2 \choose m-2} \frac{b_m}{\binom{n}{m-1}} \right]\\
           \leq &\, \sqrt{2d^{-3} } \sum_{m\in \mathcal{M}} (m-1)a_m \leq \sqrt{2d^{-1}} \leq \sqrt{d}\,.
        \end{align*}
    \end{subequations}
%where the second inequality holds since $\widetilde{\bA}_1$ is obtained by zeroing out elements of $\E\bA$, and the third inequality holds since $a_m\geq b_m$ for all $m\in \mathcal{M}$, while $\mathcal{M}$ is obtained from \Cref{alg:parameter_preprocessing}.
\end{proof}
Moreover, taking $\tau=20\mathcal{M}_{\max}, K=3$ in \Cref{thm:regularization}, with probability at least $1 - n^{-2}$, we have 
\begin{align}\label{eq:E1I1}
    \|(\bE_1)_{\cI_1}\|\leq C_3\sqrt{d}
\end{align}
where $C_3$ is a constant depending only on $\mathcal{M}_{\max}$.
Together with upper bounds for $\|(\bE_1)_{\cI_1}\|$ and $\|\bDelta_1\|$, Lemma \ref{lem:subspace_angle_k}
shows that the angle between $\bU$ and $\widetilde{\bU}$ is relatively small with high probability.

\begin{lemma}\label{lem:subspace_angle_k}
    For any $c\in (0,1)$, there exists some constant $C_2$ such that, if
  $
    \sum_{m\in\mathcal M}(m-1)(a_m-b_m)\geq C_2k^{\mathcal M_{\max} - 1}\sqrt{d}\,,
  $ 
then $\sin \angle (\bU, \widetilde{\bU}) \leq c$ with probability $1 - n^{-2}$. Here $\angle (\bU, \widetilde{\bU})$ is the angle between $\bU$ and $\widetilde{\bU}$. 
\end{lemma}

\begin{proof} 
From \Cref{eq:E1I1} and \Cref{cor:norm_of_high_degree_vertices_k_block}, with high probability,
\begin{align}
     \|(\bA_1)_{\cI_1} - \widetilde{\bA}_1\| \leq  \|(\bE_1)_{\cI_1}\| + \|\bDelta_1 \|\leq (C_3 + 1)\sqrt{d}. \notag 
\end{align}
Since $\sigma_{k+1}(\widetilde{\bA}_1)=0$, using Lemma \ref{lem:singular_value_approximation} to approximate $\sigma_k(\widetilde{\bA}_1)$, we obtain
\begin{subequations}
\begin{align*}
    &\,\sigma_k(\widetilde{\bA}_1) - \sigma_{k+1}(\widetilde{\bA}_1) =\sigma_k(\widetilde{\bA}_1)=(1+o(1))\sigma_k(\overline{\bA}_1) \geq \frac{1}{2}\sigma_k(\overline{\bA}_1)\\
    \geq &\, \frac{n}{4\sqrt{2}k }\sum_{m \in \mathcal{M}} { \frac{3n}{4k} -2 \choose m-2} \frac{a_m - b_m}{{n \choose m-1}} \geq \frac{1}{8k}\sum_{m\in \mathcal M} \bigg(\frac{3}{4k} \bigg)^{m-2}(m-1)(a_m-b_m)\\
    \geq &\, \frac{1}{8k}\left(\frac{1}{2k}\right)^{\mathcal M_{\max} - 2}\sum_{m\in\mathcal M} (m-1)(a_m-b_m)\geq \frac{C_2 \sqrt{d}}{2^{\mathcal{M}_{\max} + 1}}\,.
\end{align*}
\end{subequations}
Then for any $c\in(0, 1)$, we can find $C_2 = [2^{\mathcal{M}_{\max} + 2}(C_3 + 1)/c]$ such that  $\|(\bA_1)_{\cI_1} - \widetilde{\bA}_1\|\leq (1-1/\sqrt{2}) \sigma_k(\widetilde{\bA}_1) 
$, and by Wedin's Theorem (\Cref{lem:Wedin_sin}), the angle between $\bU$ and $\widetilde{\bU}$ is bounded by
    \begin{align}
        \|P_{\bU} - P_{\widetilde{\bU}}\| \leq &\, \frac{\sqrt{2}\|(\bA_1)_{\cI_1} - \widetilde{\bA}_1\|}{ \sigma_k(\widetilde{\bA}_1)} \leq \frac{\sqrt{2}(C_3+1)}{C_2/ 2^{\mathcal M_{\max} + 1}}
        = \frac{\sqrt{2}}{2}c< c\,. \notag 
    \end{align}
\end{proof}

\subsubsection{Bound the projection error.}\label{subsubsec:bound_the_projection_error} Randomly pick $s=2k\log^2 n$ vertices from $Y_2$. Let $\ba_{i_1},\dots, \ba_{i_s}$, $\widetilde{\ba}_{i_1},\dots, \widetilde{\ba}_{i_s}$, $\overline{\ba}_{i_1},\dots, \overline{\ba}_{i_s}$ and $\be_{i_1}, \dots, \be_{i_s}$ be the corresponding columns of $\bA_2$, $\widetilde{\bA}_2$, $\overline{\bA}_2$ and $\bE_2:= \bA_2 - \widetilde{\bA}_2$ respectively, where $\bA_2$, $\widetilde{\bA}_2$ and $\overline{\bA}_2$ were defined in \Cref{eqn:A1A2}, \Cref{eqn:tilde_A1A2} and \Cref{eqn:bar_A1A2}. Let $k(i)$ denote the membership of vertex $i$, hence $i \in V_{k(i)}\cap Y_2$ for any $i \in \{i_1, \cdots, i_s\}$. Define the corresponding vector $\widetilde{\bdelta}_i \in \R^{n}$ with the entries of $\widetilde{\ba}_{i}$ and $\widetilde{\bdelta}_i$ given by
\begin{equation}\label{eqn:tilde_delta}
    \widetilde{\ba}_{i}(j) = \left\{
        \begin{aligned}
            &\widetilde{\alpha}_{ii}, & \textnormal{if } j\in Z\cap V_{k(i)}\\
            &\widetilde{\beta}_{ij}, & \textnormal{if } j\in Z\setminus  V_{k(i)} \\
            &0\,, & \textnormal{if } j\in Y
        \end{aligned}
    \right.\,,\quad \widetilde{\bdelta}_i(j) = \left\{
        \begin{aligned}
            &(\widetilde{\alpha}_{ii} - \widetilde{\beta}_{ij})/2 >0, & \textnormal{if } j\in Z\cap V_{k(i)} \\
            &(\widetilde{\beta}_{ij} - \widetilde{\alpha}_{ii})/2 <0, & \textnormal{if } j\in Z\setminus V_{k(i)} \\
            &0, & \textnormal{if } j\in Y
        \end{aligned}
    \right. \,. 
\end{equation}
where $\widetilde{\alpha}_{ii}$, $\widetilde{\beta}_{ij}$ were defined in \Cref{eqn:tilde_alpha}, \Cref{eqn:tilde_beta} respectively. Consequently, we can identify the block $Z\cap V_{k(i)}$ by recovering $\widetilde{\bdelta}_i$ according to the sign difference between elements in $Z\cap V_{k(i)}$ and $Z\setminus V_{k(i)}$. However, it is hard to handle with $\widetilde{\bdelta}_i$ due to the randomness of $\widetilde{\alpha}_{ii}$, $\widetilde{\beta}_{ij}$. A good approximation of $\widetilde{\bdelta}_i$ was given by $\overline{\bdelta}_i := \overline{\ba}_{i} - \overline{\ba}$, where $\overline{\ba}(j):= \mathbbm{1}_{j\in Z}\cdot (\overline{\alpha} + \overline{\beta})/2$ and entries of $\overline{\bdelta}_i$ given by
\begin{equation*}
    \overline{\ba}_i (j) = \left\{
        \begin{aligned}
            &\overline{\alpha}, & \textnormal{if } j\in Z\cap V_{k(i)} \\
            &\overline{\beta}, & \textnormal{if } j\in Z\setminus V_{k(i)} \\
            &0, & \textnormal{if } j\in Y
        \end{aligned}
        \right. \,,\quad \overline{\bdelta}_i(j) = \left\{
        \begin{aligned}
            &(\overline{\alpha} - \overline{\beta})/2 >0, & \textnormal{if } j\in Z\cap V_{k(i)} \\
            &(\overline{\beta} - \overline{\alpha})/2 <0, & \textnormal{if } j\in Z\setminus V_{k(i)} \\
            &0, & \textnormal{if } j\in Y
        \end{aligned}
    \right. \,.
\end{equation*}
By construction, $\overline{\bdelta}_i$ identifies $Z \cap V_{k(i)}$ in the case of \textit{perfect splitting} for any $i\in\{i_1, \cdots, i_s\}$. However, the fluctuations of $n_i$ and $n^{\prime}_i$ are up to $O(\sqrt{n}\log n)$ \textit{w.h.p.}, as proved in \Cref{lem:singular_value_approximation}. The error in this approximation can be then corrected in \Cref{alg:k_block_correction}. Note that
\begin{align*}
    \ba_i - \overline{\ba} = (\ba_i - \widetilde{\ba}_i) + (\widetilde{\ba}_i - \overline{\ba}_{i}) + (\overline{\ba}_i - \overline{\ba}) = \be_{i} + (\widetilde{\ba}_i - \overline{\ba}_{i}) + \overline{\bdelta}_i\,,\quad \forall i \in V_{k(i)}\cap Y_2 \cap  \{i_1, \cdots, i_s\}\,.
\end{align*}
Since we do not have access to $\widetilde{\bdelta}_i$ and $\overline{\bdelta}_i$, we  use $P_{\bU}(\ba_i - \overline{\ba})$ as a good approximation, and
\begin{align}
    P_{\bU}(\ba_i - \overline{\ba}) -  \widetilde{\bdelta}_i = P_{\bU}\be_{i} + P_{\bU}(\widetilde{\ba}_i - \overline{\ba}_{i}) + P_{\bU}\overline{\bdelta}_i - \widetilde{\bdelta}_i\,. \notag 
\end{align}
By definitions of $\overline{\alpha}$ and $\overline{\beta}$ in \Cref{eqn:bar_alphabeta}, we have
\begin{subequations}
\begin{align*}
 &\, \|P_{\bU}(\widetilde{\ba}_i - \overline{\ba}_{i}) \|_2 \leq \|\widetilde{\ba}_i - \overline{\ba}_{i}\|_2 \leq O\big[ k\cdot \sqrt{n}\log n \cdot (\overline{\alpha} -\overline{\beta})^2 \big]^{1/2} = O[ n^{-3/4}\cdot \log^{1/2}(n) ]\,\\
&\, \|P_{\bU}\overline{\bdelta}_i - \widetilde{\bdelta}_i\|_2 \leq \|\overline{\bdelta}_i\|_{2} + \|\widetilde{\bdelta}_i\|_{2} = 2 \bigg[ \bigg( \sum_{i=1}^{k}n_i \bigg) \cdot (\overline{\alpha} -\overline{\beta})^2 \bigg]^{1/2} = O( n^{-1/2})\,.
\end{align*}
\end{subequations}
% To simplify the rest of the argument, we can obtain accuracy. We defer the proofs of \Cref{lem:projection_deviation_k} 

Lemma \ref{lem:projection_deviation_k} shows that with high probability, at least half of the vectors in $\be_{i_1}, \dots, \be_{i_s}$ has small $l_2$-norm of order $O( n^{-1/2})$, after projection onto $\bU$.
\begin{lemma}\label{lem:projection_deviation_k}
With probability $1 - O(n^{-k\log n})$, at least $s/2$ of the vectors $\be_{i_1}, \dots, \be_{i_s}$ satisfy
    \begin{align}
       \|P_{\bU}\be_i \|_2 \leq 2 \sqrt{k (\alpha + d \mathcal{M}_{\max}/n )} = O(n^{-1/2})\,, \quad \forall i_1, \cdots, i_s \in [n]\,.
    \end{align}
\end{lemma}
To avoid introducing extra notations, let $\ba_{i_1}, \dots, \ba_{i_{s_1}}$ denote those vectors satisfying \[\|P_{\bU}\be_{i_j} \|_2 \leq 2 \sqrt{k (\alpha + d \mathcal{M}_{\max}^2/n )}\] for $j\in [s_1]$, hence referred as \textit{good} vectors. \Cref{lem:projection_deviation_k} indicates that the number of good vectors satisfies $s_1\geq s/2 = k\log^2 n$. Together with \Cref{lem:subspace_angle_k} and discussion in \Cref{subsubsec:bound_the_projection_error}, \Cref{lem:projection_error_k} proves that \textit{good} vectors have vanishing projection errors.

\begin{lemma}[Projection error]\label{lem:projection_error_k}
    If conditions in Lemma \ref{lem:subspace_angle_k} are satisfied, then for all good vectors $\ba_{i_j}$ with $j\in [s_1]$, $\|P_{\bU}(\ba_{i_j} - \overline{\ba}) - \widetilde{\bdelta}_{i_j}\|_2 =O( n^{-1/2})$ holds with probability $1 - O(n^{-k})$.
\end{lemma}
\begin{comment}
\begin{proof}
    \begin{subequations}
        \begin{align}
            &\, \|P_{\bU}(\ba_{i_j} - \widetilde{\ba}) - \widetilde{\bdelta}_{i_j}\|_2 = \|P_{\bU}\be_{i_j} +  P_{\widetilde{\bU}}\widetilde{\bdelta}_{i_j} + (P_{\bU} - P_{\widetilde{\bU}})\widetilde{\bdelta}_{i_j} - \widetilde{\bdelta}_{i_j}\|_2\\
        \leq &\, \|P_{\bU}\be_{i_j}\|_2 +  \| (1 - P_{\widetilde{\bU}})\widetilde{\bdelta}_{i_j}\|_2 + \|(P_{\bU} - P_{\widetilde{\bU}})\widetilde{\bdelta}_{i_j} \|_2\\
        \leq &\, 2 \sqrt{\overline{\alpha}\cdot k(1 + 2k\mathcal{M}_{\max}) } + \|\widetilde{\bdelta}_{i_j}\|_2 + c\|\widetilde{\bdelta}_{i_j}\|_2 \lesssim \frac{1}{\sqrt{n}}\,.
        \end{align}
    \end{subequations}
\end{proof}
\end{comment}

\subsubsection{Accuracy}\label{subsubsec:discard_half_blue_accuracy}
For each projected vector $P_{\bU}(\ba_{i} - \overline{\ba})$, let $U^{\prime}_{i}$ denote the set of its top $n/(2k)$ coordinates in value, where $i\in \{i_1, \cdots, i_s\}$ and $s = 2k \log^2 n$. In \Cref{alg:k_block_spectral_partition}, we discard half of the obtained subsets $U^{\prime}_{i}$ with lowest blue hyperedge densities. \Cref{lem:discard_half_blue} shows that for every good index $i_j$ with vanishing projection errors as referred in \Cref{lem:projection_error_k}, the vertex set $U^{\prime}_{i_j}$ contains at least $\nu$ fraction of the vertices in $V_{k(i_j)} \cap Z$, where $k(i_j)$ denotes the membership of vertex $i_j$.

\begin{lemma}\label{lem:discard_half_blue} 
    Suppose that we are given a set $X \subset Z$ with size $|X| = n/(2k)$. Define
\begin{subequations}
\begin{align*}
    \mu_1:=&\,\frac{1}{2} \sum_{m\in \mathcal{M}}m(m-1) \left\{ \left[ \binom{\frac{\nu n}{2k}}{m} + \binom{\frac{(1 - \nu)n}{2k}}{m} \right] \frac{a_m - b_m }{ \binom{n}{m-1} } + \binom{\frac{n}{2k}}{m} \frac{b_{m}}{ \binom{n}{m-1} } \right\}\,,\\
    \mu_{2} :=&\, \frac{1}{2} \sum_{m\in \mathcal{M}} m(m-1)\left\{ \left[ \binom{\frac{(1 + \nu) n}{4k} }{m} + (k-1)\binom{ \frac{(1 - \nu)n}{4k(k-1)}}{m} \right] \frac{a_m - b_m }{ \binom{n}{m-1} } + \binom{\frac{n}{2k}}{m} \frac{b_{m}}{ \binom{n}{m-1} } \right\}\,,
\end{align*}
\end{subequations}
and $\mu_{\mathrm{T}} := (\mu_1 + \mu_2)/2 \in [\mu_1, \mu_2]$. There is a constant $c>0$ depending on $k, a_m, \nu$ such that for sufficiently large $n$,
    \begin{enumerate}[label=(\roman*)]
        \item If $|X\cap V_i| \leq \nu |X|$ for all $i\in[k]$, then with probability at least $1 - e^{-cn}$, the number of blue hyperedges in the hypergraph induced by $X$ is at most $\mu_{\mathrm{T}}$.

        \item Conversely, if $|X\cap V_i| \geq (1/2)\cdot (1+ \nu)|X|$ for some $i\in\{1, \dots, k\}$, then with probability at least $1 - e^{-cn}$, the number of blue hyperedges in the hypergraph induced by $X$ is at least $\mu_{\mathrm{T}}$.
    \end{enumerate}
\end{lemma}
\begin{remark}
When taking $M=2$, \Cref{lem:discard_half_blue} reduces to \cite[Lemma 31]{chin2015stochastic}. The original proof of \cite[Lemma 31]{chin2015stochastic} contains an error and we are able to fix it in our proof.
\end{remark}

% \begin{remark}
% The vertex set $U^{\prime}_{i_j}$ contains a fraction of the vertices in $V_{k(i_j)} \cap Z$ that is at least $\nu \in (0.5, 1)$ with high probability, where $k(i_j)$ denotes the membership of vertex $i_j$. The key point is to make sure the $\mu_2 - \mu_1 = \Omega(n)$.
% \end{remark}


\begin{comment}
% original statement with numerical numbers
\begin{lemma}
    Suppose that we are given a set $X \subset Z$ with size $|X| = (n/2k)$. Define
\begin{subequations}
\begin{align}
    \mu_1:=&\,\frac{1}{2} \sum_{m\in \mathcal{M}} \left\{ \left[ \binom{\frac{0.9n}{2k}}{m} + \binom{\frac{0.1n}{2k}}{m} \right] \frac{a_m - b_m }{ \binom{n}{m-1} } + \binom{\frac{n}{2k}}{m} \frac{b_{m}}{ \binom{n}{m-1} } \right\}\,,\\
    \mu_{2} :=&\, \frac{1}{2} \sum_{m\in \mathcal{M}} \left\{ \left[ \binom{\frac{0.95n}{2k} }{m} + (k-1)\binom{ \frac{0.05n}{2k(k-1)}}{m} \right] \frac{a_m - b_m }{ \binom{n}{m-1} } + \binom{\frac{n}{2k}}{m} \frac{b_{m}}{ \binom{n}{m-1} } \right\}\,,
\end{align}
\end{subequations}
and $\mu_{\mathrm{T}} := (\mu_1 + \mu_2)/2 \in [\mu_1, \mu_2]$. There is a constant $c>0$ depending on $k, a_m$ such that
    \begin{enumerate}[label=(\roman*)]
        \item If $|X\cap V_i| \leq 0.9|X|$ for all $i\in[k]$, then with probability at least $1 - e^{-cn}$, the number of blue hyperedges in the hypergraph induced by $X$ is at most $\mu_{\mathrm{T}}$.

        \item Conversely, if $|X\cap V_i| \geq 0.95|X|$ for some $i\in\{1, \dots, k\}$, then with probability at least $1 - e^{-cn}$, the number of blue hyperedges in the hypergraph induced by $X$ is at least $\mu_{\mathrm{T}}$.
    \end{enumerate}
\end{lemma}
\end{comment}

Let $U_{i}^{\prime}$ denote the set of $n/(2k)$ largest coordinates of the projected vector $P_{\bU}(\ba_i - \overline{\ba})$ with $i\in V_{k(i)}\cap Y_2$. According to \Cref{lem:discard_half_blue}, if the blue hyperedge density inside $U_{i}^{\prime}$ is at least $\mu_{\mathrm{T}}$, then it guarantees  that $|U_{i}^{\prime}\cap V_{k(i)}| \geq \nu |U_{i}^{\prime}|$. Therefore, it is enough to consider half of the sets with highest blue edge densities in \Cref{alg:k_block_spectral_partition}.

\begin{lemma}\label{lem:sample_k_distinct}
    Through random sampling without replacement in Step 6 of Algorithm \ref{alg:k_block_spectral_partition}, we can find at least $k$ vertices  in $Y_2$ among $k\log^2 n$ samples which belongs to distinct $V_i$, $i\in [k]$ with probability $1 -n^{-\Omega(\log n)}$.
\end{lemma}


\begin{comment}
This is different from Coupon collector's problem where you are sampling with replacement. However, we are expected to have fewer trials when sampling without replacement, since we avoid the chance to draw previous coupons again. In particular, let $T$ be the number of draws needed to collect all $k$ distinct coupons, and let $t_i$ be the number of trials to collect the $i^{\mathrm{th}}$ coupon after $i - 1$ coupons have been collected, then $T = t_1 + \cdots + t_k$. The probability of collecting a new coupon is
\begin{align}
    p_i = \frac{(n/k)[k - (i-1)]}{n - \# \textnormal{ of trials made}} \ge \frac{(n/k)[k - (i-1)])}{n} = \frac{[k - (i-1)]}{k}\,.
\end{align}
Consequently, $t_i$ has geometric distribution with expectation
\begin{align}
    \E t_i = \frac{1}{p_i} \leq \frac{k}{k -(i-1)}\,.
\end{align}
By the linearity and the upper bounds of $\E t_i$, we have
\begin{subequations}
    \begin{align}
        \E T =&\, \sum_{i=1}^{k} \E t_i \leq  k \sum_{i=1}^{k}  \frac{1}{k -(i-1)} = k \sum_{i=1}^{k}  \frac{1}{i} = k \cdot H_k \\
        \approx &\, k\log k + \gamma k + \frac {1}{2} + O \left( \frac{1}{k} \right)\,, \quad \gamma \approx 0.5772156649\,,
    \end{align}
\end{subequations}
where $H_k$ is the $k^{\mathrm{th}}$ harmonic number and the last approximation is obtained by asymptotics of the harmonic numbers. The variance of $T$ can be calculated by independence between $t_i$,
\begin{align}
    \Var(T) = \sum_{i=1}^{k}\Var(t_i) = \sum_{i=1}^{k} \frac{1 - p_i}{p_i^2} <  \sum_{i=1}^{k} \frac{1}{p_i^2} = k^2 \sum_{i=1}^{k}\frac{1}{i^2} < \frac{k^2 \pi^2}{6}\,.
\end{align}
Then by Chebyshev's inequality
\begin{align}
    \P(|\E T - k\cdot H_k| > ck ) \leq \frac{\Var(T)}{c^2 k^2} < \frac{\pi^2}{6c^2}\,.
\end{align}

Let $T_i^r$ denote the event that the $i^{\mathrm{th}}$ coupon was not picked in the first $r$ trials when sampling without replacement, then
\begin{align}
    \P\left(T_{i}^{r}\right)\leq \left(1- \frac{n/k}{n}\right)\cdot \left(1- \frac{n/k}{n-1}\right) \cdots \left(1- \frac{n/k}{n - r}\right) \leq \left(1- \frac{1}{k}\right)^r \leq e^{-r/k}\,.
\end{align}
Let $r = ck\log(k)$ for some $c \geq 5$, then $\P\left(T_{i}^{r}\right) \leq k^{-c}$, and
\begin{align}
    \P\left[T \geq c k \log(k)\right] \leq \P\left( \cup_{i=1}^{k}T_{i}^{ck\log(k)} \right) \leq k\cdot k^{-c} = k^{-c + 1}\,.
\end{align}
As a result, for any $k\geq 3$, we can draw $k$ distinct coupons within $ck \log(k)$ trials. In \Cref{alg:k_block_spectral_partition}, we select $k\log(n)$ different vertices from $Y_2$, which is enough to guarantee $k\log(n) \geq 5k\log(k)$ since $k$ is fixed and $n \gg k$.
\end{comment}

\subsection{Local Correction: Proof of Lemma \ref{lem:correction_accuracy_k}}\label{subsec:local_correction}
For notional convenience, let $U_i:= Z \cap V_i$ denote the intersection of $Z$ and true partition $V_i$ for all $i\in[k]$. In \Cref{alg:k_block_partition}, we first color the hyperedges by red and blue with equal probability. By running \Cref{alg:k_block_spectral_partition} on the red hypergraph, we obtain a $\nu$-correct partition $U_1^{\prime}, \dots, U_k^{\prime}$, i.e.,
\begin{equation}\label{eqn:spectral_clustering_accuracy_k}
    |U_{i}\setminus U_{i}^{\prime}| \leq (1 - \nu)\cdot |U_{i}^{\prime}| = (1 - \nu)\cdot \frac{n}{2k}\,, \quad \forall i\in [k]\,.
\end{equation}
In the rest of this subsection, we condition on the event \eqref{eqn:spectral_clustering_accuracy_k}.
% \begin{comment}
% \begin{lemma}
%     Given any $0.1$-correct partition $U_1^{\prime}, \cdots, U_k^{\prime}$ of $Z = (Z\cap V_{1}) \cup \cdots \cup (Z\cap V_{k})$ and the red hypergraph over $Z$, the \Cref{alg:k_block_correction} computes a $\gamma$-correct partition $\widehat{U}_1, \cdots, \widehat{U}_k$ with probability at least $1 -2k e^{-n}$, while $\gamma$ satisfying
% \begin{align}
%     \gamma = 2k\cdot \exp\left( -C_{\mathcal M} \cdot \mathrm{SNR}_{\mathcal{M}} \right)\,.
% \end{align}
% where $\mathcal{M}\subset \{2, \cdots, M\}$ is obtained from \Cref{alg:parameter_preprocessing} with $\mathcal{M}_{\max}$ denoting its maximum value, and
% \begin{align}
%     \mathrm{SNR}_{\mathcal{M}}:= \frac{ \left[\sum_{m\in \mathcal{M}} (m-1)\left(\frac{a_m - b_m}{k^{m-1}} \right) \right]^2 }{\sum_{m\in \mathcal{M}} (m-1)\left(\frac{a_m - b_m}{k^{m-1}} + b_m \right)}\,, \quad C_{\mathcal M} = \frac{[(0.9)^{\mathcal{M}_{\max}-1} - (0.1)^{\mathcal{M}_{\max}-1}]^2 }{5\cdot 2^{2\mathcal{M}_{\max} }}\,.
% \end{align}
% \end{lemma}
% \end{comment}
Consider a hyperedge $e = \{i_{1}, \cdots, i_{m}\}$ in the underlying $m$-uniform hypergraph. If vertices $i_{1}, \cdots, i_{m}$ are from the same block, then $e$ is a red hyperedge with probability $a_{m}/2\binom{n}{m-1}$;  if vertices $i_{1}, \cdots, i_{m}$ are not from the same block, then $e$ is a red hyperedge with probability $b_{m}/2\binom{n}{m-1}$. The presence of those two types of hyperedges can be denoted by 
% indicator random variables $T_{e}^{(a_m)}$ and $T_{e}^{(b_m)}$ respectively, with definitions 
\begin{equation*}
    T_{e}^{(a_m)}\sim \mathrm{Bernoulli}\left( \frac{a_{m}}{ 2{n \choose m-1} }\right)\,, \quad T_{e}^{(b_m)}\sim \mathrm{Bernoulli}\left( \frac{b_{m}}{ 2{n \choose m-1} }\right)\,,
\end{equation*}
respectively. 
For any finite set $S$, let $[S]^{l}$ denote the family of $l$-subsets of $S$, i.e., $[S]^{l} = \{Z| Z\subseteq S, |Z| = l \}$. 
%Let $E_m([Y]^{l}, [Z]^{m-l})$ denote the set of $m$-hyperedges, where $l$ vertices are from $Y$ and the other $m-l$ vertices from $Z$ within each $m$-hyperedge. 
Consider a vertex $u\in U_1: = Z \cap V_1$. The weighted number of red hyperedges, which contains $u\in U_1$ with the remaining vertices in $U_j^{\prime}$, can be written as
\begin{align}\label{eqn:weighted_neighbor_Z}
    S_{1j}^{\prime}(u):= \sum_{m\in \mathcal{M}}(m-1)\cdot \left\{ \sum_{e\in \, \mathcal{E}^{(a_m)}_{1, j}} T_e^{(a_m)} + \sum_{e\in \mathcal{E}^{(b_m)}_{1, j} } T_e^{(b_m)}\right\}\,,\quad u\in U_1\,,
\end{align}
where $\mathcal{E}^{(a_m)}_{1, j}:= E_m([U_1]^{1}, [ U_1 \cap U_j^{\prime}]^{m-1} )$ denotes the set of $m$-hyperedges with one vertex from $[U_1]^{1}$ and the other $m-1$ from $[U_1 \cap U^{\prime}_{j}]^{m-1}$, while 
 $
 \mathcal{E}^{(b_m)}_{1, j} := E_m \Big([U_1]^{1}, \,\, [U_j^{\prime}]^{m-1} \setminus [U_1\cap U_j^{\prime}]^{m-1} \Big)
 $
denotes the set of $m$-hyperedges with one vertex in $[U_1]^{1}$ while the remaining $m-1$ vertices in $[U_j^{\prime}]^{m-1}\setminus [U_1 \cap U_j^{\prime}]^{m-1}$(not all $m$ vertices are from $V_1$) with their cardinalities
    \begin{align}
       |\mathcal{E}^{(a_m)}_{1, j}| = \binom{|U_1 \cap U^{\prime}_{j}|}{m-1}\,,\quad |\mathcal{E}^{(b_m)}_{1, j}| = \left[ \binom{|U^{\prime}_{j}|}{m-1} - \binom{|U_1 \cap U^{\prime}_{j}|}{m-1} \right]\,. \notag 
    \end{align}
We multiply $(m-1)$ in \Cref{eqn:weighted_neighbor_Z} as weight since the rest $m-1$ vertices are all located in $U^{\prime}_{j}$, which can be regarded as $u$'s neighbors in $U^{\prime}_{j}$. According to the fact $|U^{\prime}_j \cap U_j| \geq (\nu n/2k)$ in \Cref{eqn:spectral_clustering_accuracy_k} and $|U_{j}^{\prime}| = n/(2k)$ for $j\in[k]$, 
\begin{align}
    |\mathcal{E}^{(a_m)}_{1, 1}| \geq \binom{\frac{\nu n}{2k}}{m-1}\,,\quad |\mathcal{E}^{(a_m)}_{1, j}| \leq \binom{\frac{(1 - \nu)n}{2k}}{m-1}\,,\,\, j\neq 1\,. \notag 
\end{align}
To simplify the calculation, we take the lower and upper bound of $|\mathcal{E}^{(a_m)}_{1, 1}|$ and $|\mathcal{E}^{(a_m)}_{1, j}|\,(j\neq 1)$ respectively.
By taking expectation with respect to $T_{e}^{(a_m)}$ and $T_{e}^{(b_m)}$, then for any $u \in U_1$, we have
\begin{subequations}
    \begin{align*}
    \E S_{11}^{\prime}(u) &= \sum_{m\in \mathcal{M}} (m-1)\cdot \left[ \binom{\frac{\nu n}{2k}}{m-1} \frac{a_m - b_m}{2\binom{n}{m-1}} + \binom{\frac{n}{2k}}{m-1} \frac{b_m}{2\binom{n}{m-1}}\right]\,,\\
    \E S_{1j}^{\prime}(u) &= \sum_{m\in \mathcal{M}} (m-1)\cdot \left[ \binom{\frac{(1-\nu)n}{2k}}{m-1} \frac{a_m - b_m}{2\binom{n}{m-1}} +  \binom{\frac{n}{2k}}{m-1} \frac{b_m}{2\binom{n}{m-1}}\right]\,, \quad j\neq 1\,.
\end{align*}
\end{subequations}

By assumptions in \Cref{thm:mainresult_k}, $\E S_{11}^{\prime}(u) - \E S_{1j}^{\prime}(u) = \Omega(1)$. Define
\begin{align}\label{eqn:mu_correction}
    \mu_{\rm{C}} := \frac{1}{2}\sum_{m\in \mathcal{M}} (m-1)\cdot \left\{ \left[ \binom{\frac{\nu n}{2k}}{m-1} + \binom{\frac{(1 - \nu)n}{2k}}{m-1} \right] \frac{a_m - b_m}{2\binom{n}{m-1}} + 2\cdot \binom{\frac{n}{2k}}{m-1}\frac{b_m}{2\binom{n}{m-1}}\right\}\,.
\end{align}
In \Cref{alg:k_block_correction}, vertex $u$ is assigned to $\widehat{U}_{i}$ if it has the maximal number of  neighbors in $U_i^{\prime}$. If $u\in U_1$ is mislabeled, then one of the following events must happen:
\begin{itemize}
    \item $S_{11}^{\prime}(u) \leq \mu_{\rm{C}}$, meaning that $u$ was mislabeled by \Cref{alg:k_block_correction}. 
    \item $S_{1j}^{\prime}(u) \geq \mu_{\rm{C}}$, for some $j\neq 1$, meaning that $u$ survived \Cref{alg:k_block_correction} without being corrected. 
\end{itemize}
\Cref{lem:mislabel_probability_correction} shows that the probabilities of those two events are exponentially small. 
\begin{lemma}\label{lem:mislabel_probability_correction}
For sufficiently large $n$ and  any $u\in U_1 = Z \cap V_1$, we have
    \begin{align}\label{eqn:rhoC_k}
        \rho_1^{\prime} := \P \left( S_{11}^{\prime}(u) \leq \mu_{\rm{C}}\right) \leq \rho \,,\quad \rho_j^{\prime} :=\P \left( S_{1j}^{\prime}(u) \geq \mu_{\rm{C}}\right) \leq \rho\,,\,(j\neq 1),
    \end{align}
where $\rho := \exp\left( -C_{\mathcal M}\cdot \mathrm{SNR}_{\mathcal{M}} \right)$ with $\mathrm{SNR}_{\mathcal{M}}$ and $C_{\mathcal M}$ defined in \Cref{eqn:SNR_CM}.
\end{lemma}
As a result, the probability of that either of those events happened is bounded by $\rho$. The number of mislabeled vertices in $U_1$ after \Cref{alg:k_block_correction} is at most
\begin{align*}
    R_1 = \sum_{t=1}^{|U_1|}\Gamma_{t}\, + \sum_{j=2}^{k}\sum_{t=1}^{|U_1\cap U_j^{\prime}|}\Lambda_{t}\,,
\end{align*}
where $\Gamma_{t}$ (resp. $\Lambda_{t}$) are i.i.d indicator random variables with mean $\rho_1^{\prime}$ (resp. $\rho_j^{\prime}$, $j\neq 1$). Then 
\begin{align*}
   \E R_1 \leq  \frac{n}{2k} \rho_1^{\prime} + \sum_{j=2}^{k}\frac{(1 - \nu)n}{2k} \rho_j^{\prime} \leq \frac{n}{2k} \cdot k\rho = \frac{n\rho}{2}\,. 
\end{align*}
Let $t_1:= n \rho/2$, where $\nu$ is the correctness after \Cref{alg:k_block_spectral_partition}, then by Chernoff bound (\Cref{lem:Chernoff}),
\begin{align}\label{eqn:wrong_prob}
   \P \left( R_1 \geq n\rho \right) = \P \left( R_1 - n \rho/2 \geq t_1 \right) \leq \P \left( R_1 - \E R_1 \geq t_1 \right) \leq e^{-c t_1} = O(e^{-n\rho})\,.
\end{align}
Then with probability $1 - O(e^{-n\rho})$, the fraction of mislabeled vertices in $U_1$ is smaller than $k\rho$, i.e., the correctness of $U_1$ is at least $\gamma_{\rm{C}} := \max\{\nu\,, 1 - k\rho\}$. Therefore, \Cref{alg:k_block_correction} outputs a $\gamma_{\rm{C}}$-correct partition $\widehat{U}_1, \cdots, \widehat{U}_{k}$ with probability  $1- O(e^{-n\rho})$. 
\begin{remark}\label{rem:SNRlogn}
When $\mathrm{SNR}_{\mathcal{M}} = \Omega(\log n)$, we have $\rho = \Omega(n^{-1})$ and $e^{-n\rho} = \Omega(1)$ in \Cref{eqn:wrong_prob}, which may not decrease to $0$ as $n\to\infty$. As a result, we assume $\mathrm{SNR}_{\mathcal{M}} = o(\log n)$ in the statement of Corollary \ref{cor:weak_consistency}.
\end{remark}
\subsection{Merging: Proof of Lemma \ref{lem:merging_accuracy_k}}\label{subsec:merging}
By running \Cref{alg:k_block_correction} on the red graph, we obtain a $\gamma_{\rm{C}}$-correct partition $\widehat{U}_1, \cdots, \widehat{U}_{k}$ where $\gamma_{\rm{C}}:=\max\{\nu\,, 1 - k\rho\} \geq \nu$, i.e.,
\begin{equation}\label{eqn:correction_accuracy_k}
    |U_{j}\cap \widehat{U}_j| \geq \nu\cdot |\widehat{U}_j| =  \frac{\nu n}{2k}\,, \quad \forall j\in [k]\,.
\end{equation}
In the rest of this subsection, we shall condition on this event and abbreviate $Y\cap V_l$ by $W_l:=Y\cap V_l$. The failure probability of \Cref{alg:k_block_merging} is estimated by presence of hyperedges between vertex sets $Y$ and $Z$.

%We also condition on the event that the maximum Red degree of a vertex is at most $\log^{2}(n)$ with probability at least $1 - o(1)$, which can be proved by Bernstein \ref{lem:Bernstein}.

Consider a hyperedge $e = \{i_{1}, \cdots, i_{m}\}$ in the underlying $m$-uniform hypergraph. If vertices $i_{1}, \cdots, i_{m}$ are all from the same cluster $V_{l}$, then the probability that $e$ is an existing blue edge conditioning on the event that $e$ is not a red edge is 
\begin{align}\label{eqn:psi_m}
    \psi_m:=\P\left[e \text{ is a blue edge } \Big| e \text{ is not a red edge } \right] = \frac{\frac{a_{m}}{2{n \choose m-1}}}{1 - \frac{a_{m}}{2{n \choose m-1}}} = \frac{a_{m}}{2{n \choose m-1}}(1+o(1))\,,
\end{align}
and the presence of $e$ can be represented by an indicator random variable $\zeta_e^{(a_m)} \sim \mathrm{Bernoulli}\left(\psi_m\right)$.
Similarly, if vertices $i_{1}, \cdots, i_{m}$ are not all from the same cluster $V_l$, the probability that $e$ is an existing blue edge conditioning on the event that $e$ is not a red edge is
\begin{align}\label{eqn:phi_m}
    \phi_m:=\P\left[e \text{ is a blue edge } \Big| e \text{ is not a red edge } \right] = \frac{\frac{b_{m}}{2{n \choose m-1}}}{1 - \frac{b_{m}}{2{n \choose m-1}}} =\frac{b_{m}}{2{n \choose m-1}}(1+o(1))\,,
\end{align}
and the presence of this hyperedge can be represented by an indicator random variable $\xi_e^{(b_m)} \sim \mathrm{Bernoulli}\left(\phi_m\right)$. 

For any vertex $w\in W_l := Y\cap V_l$ with fixed $l \in [k]$, we want to compute the number of hyperedges containing $w$ with all remaining vertices located in vertex set $\widehat{U}_{j}$ for some fixed $j\in [k]$. Following a similar argument given in \Cref{subsec:local_correction}, this number can be written as
\begin{align}\label{eqn:weighted_neighbor_Y}
    \widehat{S}_{lj}(w):= \sum_{m\in \mathcal{M}}(m-1)\cdot \left\{ \sum_{e\in \, \widehat{\mathcal{E}}^{(a_m)}_{l, j}} \zeta_e^{(a_m)} + \sum_{e\in \widehat{\mathcal{E}}^{(b_m)}_{l, j} } \xi_e^{(b_m)}\right\}\,,\quad w\in W_l\,,
\end{align}
 where $\widehat{\mathcal{E}}^{(a_m)}_{l, j}:= E_m([W_l]^{1}, [U_l\cap \widehat{U}_j]^{m-1})$ denotes the set of $m$-hyperedges with $1$ vertex from $[W_l]^{1}$ and the other $m-1$ vertices from $[U_l\cap \widehat{U}_j]^{m-1}$, while
    $
        \widehat{\mathcal{E}}^{(b_m)}_{l, j} := E_m ([W_l]^{1}, \,\, [\widehat{U}_j]^{m-1} \setminus [U_l\cap \widehat{U}_j]^{m-1})
    $
    denotes the set of $m$-hyperedges with $1$ vertex in $[W_l]^{1}$ while the remaining $m-1$ vertices are in $[\widehat{U}_j]^{m-1}\setminus [U_l\cap \widehat{U}_j]^{m-1}$, with their cardinalities %\todo{d(u) is the red degree}
    \begin{align*}
       |\widehat{\mathcal{E}}^{(a_m)}_{l, j}| = \binom{|U_l \cap \widehat{U}_{j}|}{m-1}\,,\quad |\widehat{\mathcal{E}}^{(b_m)}_{l, j}| = \left[ \binom{|\widehat{U}_{j}|}{m-1} - \binom{|U_l \cap \widehat{U}_{j}|}{m-1} \right]\,. \notag 
    \end{align*}
Similarly, we multiply $(m-1)$ in \Cref{eqn:weighted_neighbor_Y} as weight since the rest $m-1$ vertices can be regarded as $u$'s neighbors in $\widehat{U}_{j}$. By accuracy of \Cref{alg:k_block_correction} in \Cref{eqn:correction_accuracy_k}, $|\widehat{U}_{j} \cap U_j| \geq \nu n/(2k)$, then
\begin{align*}
    |\widehat{\mathcal{E}}^{(a_m)}_{l, l}| \geq \binom{\frac{\nu n}{2k}}{m-1}\,,\quad |\widehat{\mathcal{E}}^{(a_m)}_{l, j}| \leq \binom{\frac{(1 - \nu)n}{2k}}{m-1}\,,\,\, j\neq l\,. \notag 
\end{align*}
% To simplify the calculation, we take the lower and upper bounds for $|\widehat{\mathcal{E}}^{(a_m)}_{l, l}|$ and $|\widehat{\mathcal{E}}^{(a_m)}_{l, j}|$, $(j\neq l)$, respectively. 
Taking expectation with respect to $\zeta_e^{(a_m)}$ and $\xi_e^{(b_m)}$,  for any $w \in W_l$, we have
\begin{subequations}
    \begin{align*}
    \E \widehat{S}_{ll}(w) &= \sum_{m\in \mathcal{M}} (m-1)\cdot\left[ \binom{\frac{\nu n}{2k}}{m-1} (\psi_{m} - \phi_{m}) + \binom{\frac{n}{2k}}{m-1} \phi_{m}\right],\\
    \E \widehat{S}_{lj}(w) &= \sum_{m\in \mathcal{M}} (m-1)\cdot\left[ \binom{\frac{(1 - \nu) n}{2k}}{m-1} (\psi_{m} - \phi_{m}) +  \binom{\frac{n}{2k}}{m-1} \phi_{m} \right],\,\, j\neq l\,.
\end{align*}
\end{subequations}
By assumptions in \Cref{thm:mainresult_k}, $\E \widehat{S}_{ll}(w) - \E \widehat{S}_{lj}(w) = \Omega(1)$. We define
\begin{align}
    \mu_{\rm{M}} := \frac{1}{2}\sum_{m\in \mathcal{M}} (m-1)\cdot\left\{ \left[ \binom{\frac{\nu n}{2k}}{m-1} + \binom{\frac{(1 - \nu )n}{2k}}{m-1} \right](\psi_{m} - \phi_{m}) + 2\binom{\frac{n}{2k}}{m-1} \phi_{m}\right\}\,. \notag 
\end{align}
After \Cref{alg:k_block_merging}, if a vertex $w\in W_l$ is mislabelled, one of the following events must  happen
\begin{itemize}
    \item $\widehat{S}_{ll}(w) \leq \mu_{\rm{M}}$, which implies that $u$ was mislabelled by \Cref{alg:k_block_merging}.
    \item $\widehat{S}_{lj}(w) \geq \mu_{\rm{M}}$ if $j\neq l$, which implies that $u$ survived \Cref{alg:k_block_merging} without being corrected.
\end{itemize}
By an argument similar to \Cref{lem:mislabel_probability_correction}, we can prove that for any $w\in W_l$,
\begin{align}
    \hat{\rho}_{l} := \P( \widehat{S}_{ll}(w)\leq \mu_{\rm{M}} ) \leq \rho\,,\quad \hat{\rho}_{j} := \P ( \widehat{S}_{lj}(w) \geq \mu_{\rm{M}} ) \leq \rho\,, \,(j\neq l), \notag
\end{align}
where $\rho := \exp\left( -C_{\mathcal M}\cdot \mathrm{SNR}_{\mathcal{M}} \right)$. The mis-classified probability for $w\in W_l$ is upper bounded by $\sum_{j=1}^{k}\hat{\rho}_{j} \leq k\rho$. The number of mislabelled vertices in $W_l$ is at most
$
    R_l = \sum_{t=1}^{|W_l|}\Gamma_{t} \,,
$
    where $\Gamma_{t}$ are i.i.d indicator random variables with mean $k \rho$ and $\E R_l \leq n/(2k)\cdot k\rho = n\rho/2$. Let $t_l:= n \rho/2$, by Chernoff bound (\Cref{lem:Chernoff}),
\begin{align*}
   \P \left( R_l \geq n\rho \right) = \P \left( R_l - n\rho/2\geq t_l \right) \leq \P \left( R_l - \E R_l \geq t_l \right) \leq e^{-c t_l} = O(e^{-n\rho})\,.
\end{align*}
Hence with probability $1 - O(e^{-n\rho})$, the fraction of mislabeled vertices in $W_l$ is smaller than $k\rho$, i.e., the correctness in $W_l$ is at least $\gamma_{\rm{M}} :=\max\{ \nu,\, 1 - k\rho \}$.

\subsection{Proof of  \Cref{thm:mainresult_k}}
Now we are ready to prove our main result (\Cref{thm:mainresult_k}). Note that the correctness of \Cref{alg:k_block_correction} and \Cref{alg:k_block_merging} are $\gamma_{\rm{C}}$ and $\gamma_{\rm{M}}$ respectively, then with probability at least $1 - o(1)$, the correctness $\gamma$ of \Cref{alg:k_block_partition} is
$
    \gamma := \min\{\gamma_{\rm{C}}, \gamma_{\rm{M}}\} = \max\{\nu,\, 1 - k\rho \}.
$
%A simple analysis shows that the value of correctness $\gamma$ depends on the relationship between $\text{SNR}_{\mathcal M}$ and $\nu$. 
We will have $\gamma = 1 - k\rho$ if $\nu \leq 1 - k\rho$, equivalently,
\begin{align}\label{eqn:gammacondition}
    \text{SNR}_{\mathcal M} \geq \frac{1}{C_{\mathcal M}} \log \frac{k}{1 - \nu}\,,
\end{align}
otherwise $\gamma = \nu$. The inequality \eqref{eqn:gammacondition} holds since
\begin{subequations}
    \begin{align*}
        \text{SNR}_{\mathcal M}
    =& \frac{ \left[\sum_{m\in \mathcal{M}} (m-1)\left(\frac{a_m - b_m}{k^{m-1}} \right) \right]^2 }{\sum_{m\in \mathcal{M}} (m-1)\left(\frac{a_m - b_m}{k^{m-1}} + b_m \right)}\\
    \geq & \frac{[\sum_{m\in\mathcal M} (a_m-b_m)]^2}{k^{2\mathcal M_{\max}-2}(\mathcal M_{\max}-1)d}\geq \frac{(C_{\nu})^2}{\mathcal M_{\max}-1} \log \frac{k}{1 - \nu}\geq \frac{1}{C_{\mathcal M}}\log \frac{k}{1-\nu}\,.
    \end{align*}
\end{subequations}
where the first two inequalities holds according to $d := \sum_{m\in \mathcal{M}}(m-1)a_m$ and Condition \eqref{eqn:SNRcondition}, while the last inequality holds when taking $C_{\nu}$ in \eqref{eqn:SNRcondition} satisfying  $C_{\nu}\geq \sqrt{(\mathcal M_{\max}-1)/C_{\mathcal M} }$, which finishes the proof.

\begin{remark}
Condition \eqref{eqn:gammacondition} indicates that the improvement of accuracy from local refinement (\Cref{alg:k_block_correction} and \Cref{alg:k_block_merging})  will be guaranteed when $\text{SNR}_{\mathcal M}$ is large enough. 
% However, the correctness lower bound $1 - \rho_{\rm{C}}$ is not tight when $\text{SNR}_{\mathcal M}$ is small, even though the actual correctness was improved by local refinement \todo{it wasn't proved}. 
If $\text{SNR}_{\mathcal M}$ is small,  we use correctness of \Cref{alg:k_block_spectral_partition} instead, i.e., $\gamma = \nu$, to represent the correctness of \Cref{alg:k_block_partition}. 
% The result is consistent with our \Cref{alg:parameter_preprocessing}, since we want to obtain the maximum signal-to-noise ratio to get a better performance guarantee. 
\end{remark}

\subsection{Proof of \Cref{cor:weak_consistency}}
    For any fixed $\nu \in (0.5, 1)$, $\mathrm{SNR}_{\mathcal{M}} \to \infty$ implies $\rho \to 0$ and $d =\sum_{m\in \mathcal{M}}(m-1)a_m \to \infty$. Since
    \begin{align}
    \frac{\sum_{m\in\mathcal M} (m-1)(a_m-b_m)^2}{\sum_{m}(m-1)a_m} \geq \frac{\sum_{m\in\mathcal M} (m-1)(a_m-b_m)^2}{\sum_{m}(m-1)k^{m-1}(a_m+ k^{m-1}b_m)} = \mathrm{SNR}\,, \notag 
\end{align}
    Condition \eqref{eqn:SNRcondition} is satisfied. Applying \Cref{thm:mainresult_k}, we find $\gamma = 1 - o(1)$, which implies weak consistency. The constraint of $\mathrm{SNR}_{\mathcal{M}} = o(\log n)$ is used in the proof of \Cref{lem:correction_accuracy_k}, see Remark \ref{rem:SNRlogn}.

\subsection*{Acknowledgements}
This work is partially supported by  NSF DMS-1949617.  I.D. and Y.Z. acknowledge support from NSF  DMS-1928930 during their participation in the program “Universality and Integrability in Random Matrix
Theory and Interacting Particle Systems” hosted by the Mathematical Sciences Research Institute
in Berkeley, California during the Fall semester of 2021. Y.Z. thanks Zhixin Zhou for helpful comments.
\section{Proof of Theorem \ref{thm:concentration} and  Theorem \ref{thm:regularization}}\label{sec:proof_concentration}

\subsection{Discretization}
To prove Theorem \ref{thm:concentration}, we start with a standard $\epsilon$-net argument.
% An $\epsilon$-net for a compact metric space $(\mathcal{X}, d)$ is a finite subset $\mathcal{N}$ of $\mathcal{X}$ such that for each point $\bx \in \mathcal{X}$, there is a point $\by \in \mathcal{N} \subset \mathcal{X}$ with $d(\bx, \by) \leq \epsilon$.


\begin{lemma}[Lemma 4.4.1 in \cite{vershynin_2018}]
Let $\bW$ be any Hermitian $n\times n$ matrix and let $\mathcal{N}_{\epsilon}$ be an $\epsilon$-net on the unit sphere $\mathbb{S}^{n-1}$ with $\epsilon\in (0,1)$, then 
$
\|\bW\| \leq \frac{1}{1-\epsilon}\sup_{\bx\in \mathcal{N}_{\epsilon}} |\langle \bW \bx, \bx\rangle |.
$
\end{lemma}
\begin{comment}
\begin{proof}
The spectral norm of $\bM$ is defined as
\begin{equation*}
    \|\bM\| := \sup_{\bx \in \mathbb{S}^{n-1}} \|\bM \bx\|_2 = \sup_{\bx \in \mathbb{S}^{n-1}} |\langle \bM \bx, \bx \rangle|\, . 
\end{equation*}
 Notice that $\mathcal{N}_{\epsilon}$ is an $\epsilon$-net on the unit sphere $\mathbb{S}^{n-1}$, then for any $\bx \in \mathbb{S}^{n-1}$, there exists some $\by \in \mathcal{N}_{\epsilon}$ such that $\|\bx - \by\|_2 \leq \epsilon$, and we have
\begin{equation*}
    \|\bM \bx\|_2 - \|\bM \by\|_2 \leq \|\bM \bx - \bM \by\|_2 \leq \|\bM\|_2 \cdot \|\bx - \by\|_2 \leq \epsilon \|\bM\|\, ,
\end{equation*}
leading to the inequality $\|\bM \bx\|_2 - \epsilon \|\bM\| \leq \|\bM \by\|_2$. For any $\by \in \mathcal{N}_{\epsilon}$, if we take the supremum over $\bx$ on the left hand side, then 
\begin{equation*}
    \sup_{\bx \in \mathbb{S}^{n-1}}\|\bM \bx\|_2 - \epsilon \|\bM\| = (1 - \epsilon) \|\bM\| \leq \|\bM \by\|_2 \leq \sup_{\bz \in \mathcal{N}_{\epsilon}} \|\bM \bz\|_2. 
\end{equation*}
Therefore,
\begin{equation*}
    \|\bM\| \leq \frac{1}{1 - \epsilon} \sup_{\bx \in \mathcal{N}_{\epsilon}} \|\bM \bx\|_2 = \frac{1}{1 - \epsilon} \sup_{\bx\in \mathcal{N}_{\epsilon}} |\langle \bM \bx, \bx \rangle|\, . 
\end{equation*}
\end{proof}
\end{comment}
By \cite[Corollary 4.2.13]{vershynin_2018}, the size of $\mathcal{N}_{\epsilon}$ is bounded by $|\mathcal{N}_{\epsilon}|\leq (1+2/\epsilon)^n$. We would have $\log|\mathcal{N}| \leq n \log(5)$ when $\mathcal{N}$ is taken as an $(1/2)$-net of $\mathbb S^n$. Define $\bW := \bA - \E \bA$, then $\bW_{ii}=0$ for each $i\in [n]$ by the definition of adjacency matrix in \Cref{eqn:expected_adjacency_k}, and we obtain 
\begin{align}\label{eqn:epsilon_net}
    \|\bA - \E \bA\| := \|\bW\| \leq 2\sup_{\bx \in \mathcal{N} }| \langle \bW \bx, \bx\rangle|\,.
\end{align}
 For any fixed $\bx\in \mathbb{S}^{n-1}$, consider the \textit{light} and \textit{heavy} pairs as follows.
\begin{align}\label{eqn:lightheavypair}
    &\mathcal{L}(\bx)=\Bigg\{ (i,j): |\bx_i \bx_j|\leq {\frac{\sqrt d}{n}}\Bigg\}, \quad \mathcal{H}(\bx)=\Bigg\{ (i,j): |\bx_i \bx_j|> {\frac{\sqrt d}{n}}\Bigg\}\,,
\end{align}
where $d = \sum_{m=2}^{M}(m-1)d_m$. Thus by the triangle inequality, 
\[
    | \langle \bx, \bW \bx\rangle | \leq \left|\sum_{(i,j)\in \mathcal{L}(\bx)} \bW_{ij}\bx_i \bx_j \right| + \left|\sum_{(i,j)\in \mathcal{H}(\bx)} \bW_{ij}\bx_i \bx_j \right|
,\] 
and by \Cref{eqn:epsilon_net},
\begin{align}\label{eq:contribution}
    \|\bA - \E \bA\| \leq 2\sup_{\bx \in \mathcal{N} }\Bigg|\sum_{(i,j)\in \mathcal{L}(\bx)} \bW_{ij} \bx_i \bx_j\Bigg| + 2\sup_{\bx\in \mathcal{N} }\Bigg|\sum_{(i,j)\in \mathcal{H}(\bx)} \bW_{ij} \bx_i \bx_j \Bigg|\,.
\end{align}
\subsection{Contribution from light pairs}\label{sec:light}
For each $m$-hyperedge $e\in E_m$, we define $\bcW^{(m)}_e := \bcT^{(m)}_e - \E\bcT^{(m)}_e.$ Then for any fixed $\bx\in \mathbb{S}^{n-1}$, the contribution from light couples can be written as

\begin{align} 
       &\,\sum_{(i,j)\in \mathcal{L}(\bx)} \bW_{ij} \bx_i \bx_j = \sum_{(i,j)\in \mathcal{L}(\bx)} \left( \sum_{m=2}^{M}\,\, \sum_{\substack{e\in E_m\\ \{i,j\}\subset e} } \bcW^{(m)}_e \right) \bx_i\bx_j  \notag \\
       =&\, \sum_{m=2}^M \sum_{e\in E_m} \bcW^{(m)}_e \left( \sum_{\substack{(i,j)\in \mathcal{L}(\bx)\\ i\neq j,\, \{i,j\}\subset e} }\bx_i \bx_j \right) = \sum_{m=2}^M \sum_{e\in  E_m} \bcY_e^{(m)} \,,\label{eq:lightW}
\end{align}
where the constraint $i\neq j$ comes from the fact $\bW_{ii} =0$ and we denote
\begin{align*}
    \bcY_e^{(m)}:= \bcW^{(m)}_e \left( \sum_{\substack{(i,j)\in \mathcal{L}(\bx)\\i\neq j, \, \{i,j\}\subset e}}\,\, \bx_i \bx_j\right)\,.
\end{align*}
Note that $\E \bcY_e^{(m)} = 0$, and by the  definition of light pair (\ref{eqn:lightheavypair}),
\begin{align*}
|\bcY_e^{(m)}|\leq m(m-1)\sqrt{d}/n \leq M(M-1)\sqrt{d}/n\,, \quad \forall m\in \{2, \cdots, M\}\,.
\end{align*}
Moreover, \Cref{eq:lightW} is a sum of independent, mean-zero random variables, and
\begin{subequations}
\begin{align*}
    &\,\sum_{m=2}^M \sum_{e\in E_m} \E [(\bcY_e^{(m)})^2] := \sum_{m=2}^M \sum_{e\in E_m} \Bigg[ \E[(\bcW_{e}^{(m)})^2] \Bigg( \sum_{\substack{(i,j)\in \mathcal{L}(\bx)\\ i\neq j, \{i,j\}\subset  e}}\bx_i \bx_j \Bigg)^2 \Bigg] \\
    \leq &\, \sum_{m=2}^M \sum_{e\in E_m} \Bigg[ \E[ \bcT^{(m)}_e ] \cdot m(m-1) \Bigg(\sum_{\substack{(i,j)\in \mathcal{L}(x)\\ i\neq j, \{i,j\}\subset e}} \bx_i^2 \bx_j^2 \Bigg) \Bigg] 
    \leq \sum_{m=2}^M \frac{d_m  \cdot  m(m-1)}{\binom{n}{m-1}}\binom{n}{m-2} \sum_{(i,j)\in [n]^2}\bx_i^2 \bx_j^2 \\
    \leq &\, \sum_{m=2}^M \frac{d_m m(m-1)^2}{n-m+2} \leq \frac{2}{n}\sum_{m=2}^M d_m (m-1)^3\leq \frac{2d(M-1)^2}{n}\,,
\end{align*}
\end{subequations}
when $n\geq 2m-2$,
where $d_m=\max d_{i_1,\dots, i_m}$ and $d = \sum_{m=2}^{M}(m-1)d_m$. Then Bernstein's inequality (Lemma \ref{lem:Bernstein}) implies that for any $\alpha>0$,
\begin{subequations}
\begin{align*}
    &\P \Bigg( \Bigg| \sum_{(i,j)\in\mathcal{L}(\bx)} \bW_{ij}\bx_i \bx_j\Bigg| \geq \alpha\sqrt{d}\Bigg) = \P \Bigg( \Bigg| \sum_{m=2}^M \sum_{e\in  E_m} \bcY_e^{(m)} \Bigg| \geq \alpha\sqrt{d}\Bigg) \\
    \leq & 2\exp \Bigg( -\frac{\frac{1}{2}\alpha^2d}{\frac{2d}{n}(M-1)^2+ \frac{1}{3} (M-1)M \frac{\sqrt{d}}{n} \alpha\sqrt{d}}\Bigg)
    \leq 2\exp \Bigg( -\frac{\alpha^2 n}{4(M-1)^2+\frac{2\alpha(M-1)M}{3}}\Bigg).
\end{align*}
\end{subequations}
Therefore by taking a union bound,
    \begin{align}\label{eq:lightbound}
        &\,\P \Bigg(  \sup_{\bx\in \mathcal{N} }\Bigg| \sum_{(i,j)\in\mathcal{L}(\bx)} \bW_{ij} \bx_i \bx_j \Bigg|\geq \alpha\sqrt{d}\Bigg)\leq |\mathcal{N}| \cdot \P \Bigg( \Bigg| \sum_{(i,j)\in\mathcal{L}(\bx)} \bW_{ij}\bx_i \bx_j\Bigg| \geq \alpha\sqrt{d}\Bigg) \notag\\
        \leq&\, 2\exp \Bigg( \log(5) \cdot n-\frac{\alpha^2 n}{4(M-1)^2+\frac{2\alpha(M-1)M}{3}}\Bigg)\leq 2e^{-n}\, ,
    \end{align}
where we choose $\alpha=5M(M-1)$ in the last line. 

\subsection{Contribution from heavy pairs}
Note that for any $i\not=j$,
\begin{align}\label{eq:EA}
    \E \bA_{ij}\leq \sum_{m=2}^M \binom{n-2}{m-2}\frac{d_m}{\binom{n}{m-1}}= \sum_{m=2}^M\frac{(m-1)d_m}{n}=\frac{d}{n}.
\end{align}
and
\begin{align}\label{eq:heavyexp}
   & \Bigg|\sum_{(i,j)\in \mathcal{H}(\bx)} \E \bA_{ij} \bx_i \bx_j\Bigg| =\Bigg|\sum_{(i,j)\in \mathcal{H}(\bx)} \E \bA_{ij} \frac{\bx_i^2 \bx_j^2}{\bx_i \bx_j}\Bigg| 
    \leq \sum_{(i,j)\in \mathcal{H}(\bx)} \frac{d}{n}\frac{\bx_i^2 \bx_j^2}{|\bx_i\bx_j|}\leq \sqrt{d}\sum_{(i,j)\in \mathcal{H}(\bx)} \bx_i^2 \bx_j^2\leq \sqrt{d}.
\end{align}
Therefore it suffices to show that, with high probability, 
\begin{align}\label{eq:Osqrtd}
    \sum_{(i,j)\in \mathcal{H}(\bx)}  \bA_{ij} \bx_i \bx_j=O \Big(\sqrt{d} \Big).
\end{align}
Here we use the discrepancy analysis from \cite{feige2005spectral,cook2018size}. We consider the weighted graph associated with the adjacency matrix $\bA$.

\begin{definition}[Uniform upper tail property, \textbf{UUTP}]
Let $\bM$ be an $n\times n$ random symmetric matrix with non-negative entries and define
\begin{align*}
    \mu := \sum_{i,j=1}^n \bQ_{ij}\E \bM_{ij}, \quad \tilde{\sigma}^2 :=\sum_{i,j=1}^n \bQ_{ij}^2 \E \bM_{ij}.
\end{align*}
We say that $\bM$ satisfies the uniform upper tail property $\textnormal{\textbf{UUTP}}(c_0,\gamma_0)$ with $c_0>0,\gamma_0\geq 0$, if for any $a,t>0$ and any $n\times n$ symmetric matrix $\bQ$ with entries $\bQ_{ij}\in [0,a]$ for all $i,j\in [n]$,
\begin{align*}
    \P \Bigg(  f_{\bQ}(\bM)\geq (1 + \gamma_0)\mu +t\Bigg)\leq \exp \Bigg( -c_0 \frac{\tilde{\sigma}^2}{a^2} h\Bigg( \frac{at}{\tilde{\sigma}^2}\Bigg)\Bigg).
\end{align*}
where function $f_{\bQ}(\bM): \R^{n \times n} \mapsto \R$ is defined by $f_{\bQ}(\bM):= \sum_{i,j=1}^n \bQ_{ij} \bM_{ij}$ for $\bM \in \R^{n \times n}$, and function $h(x):=(1 + x)\log(1+x) - x$ for all $x>-1$.
\end{definition}

\begin{lemma}\label{lem:UUTPforA}
Let $\bA$ be the adjacency matrix of the non-uniform hypergraph $H = \cup_{m=2}^{M}H_m$, then $\bA$ satisfies $\textnormal{\bf{UUTP}}(c_0,\gamma_0)$ with $c_0 = [M(M-1)]^{-1}, \gamma_0 = 0$.
\end{lemma}
\begin{proof}
Note that
\begin{subequations}
\begin{align*}
    &\,f_{\bQ}(\bA) - \mu = \sum_{i,j=1}^n \bQ_{ij}(\bA_{ij} - \E \bA_{ij}) = \sum_{i,j=1}^n \bQ_{ij}\bW_{ij}\\
    = &\, \sum_{i,j=1}^n \bQ_{ij} \Bigg( \sum_{m=2}^{M}\,\, \sum_{\substack{e\in E_m\\ i\neq j,\, \{i,j\}\subset e} } \bcW^{(m)}_e  \Bigg) =\sum_{m=2}^M \,\, \sum_{e\in E_m} \bcW^{(m)}_e \Bigg( \sum_{\{i,j\}\subset e, \, i\neq j}\bQ_{ij}\Bigg) =\sum_{m=2}^{M} \,\,\sum_{e\in E_m}\bcZ^{(m)}_e\,,
\end{align*}
\end{subequations}
where $\bcZ^{(m)}_e = \bcW^{(m)}_e \big( \sum_{\{i,j\}\subset e, i\neq j} \bQ_{ij}\big)$ are independent centered random variables upper bounded by
$
    |\bcZ^{(m)}_e| \leq \sum_{ \{i,j\} \subset e,\, i\neq j}\bQ_{ij}\leq M(M-1)a
$
for each $ m\in \{2, \dots, M\}$ since  $\bQ_{ij} \in [0, a]$.
Moreover, the variance of the sum can be written as
\begin{subequations}
\begin{align*}
    &\, \sum_{m=2}^M \sum_{e\in E_m} \E(\bcZ^{(m)}_e)^2:=\sum_{m=2}^M \,\, \sum_{e\in E_m} \E(\bcW^{(m)}_e)^2 \Bigg( \sum_{\{i,j\}\subset e, \, i\neq j} \bQ_{ij}\Bigg)^2\\
    \leq &\,\sum_{m=2}^M \,\,\sum_{e\in E_m}\E [\bcT^{(m)}_{e}] \cdot m(m-1) \sum_{ \{i,j \}\subset e, i\neq j }\bQ_{ij}^2
    \leq M(M-1)\sum_{i,j=1}^n \bQ_{ij}^2 \E \bA_{ij} = M(M-1)\tilde{\sigma}^2.
\end{align*}
\end{subequations}
where the last inequality holds since by definition $\E \bA_{ij} = \sum_{m=2}^M \,\,\sum_{\substack{e\in E_m\\ \{i,j\}\subset e}}\E [\bcT^{(m)}_{e}]$. Then by Bennett's inequality \ref{lem:Bennett}, we obtain
\begin{align*}
    \P (f_{\bQ}(\bA)- \mu\geq t) \leq \exp\Bigg( -\frac{\tilde{\sigma}^2}{M(M-1)a^2} h \Bigg( \frac{at}{\tilde{\sigma}^2}\Bigg) \Bigg)\,
\end{align*}
where the inequality holds since the function $x\cdot h(1/x) = (1+x)\log(1 + 1/x) - 1$ is decreasing with respect to $x$.
\end{proof}

\begin{definition}[Discrepancy property, \bf{DP}]
Let $\bM$ be an $n\times n$ matrix with non-negative entries. For $S,T\subset [n]$, define
$
    e_{\bM}(S,T)=\sum_{i\in S, j\in T}\bM_{ij}.
$
We say $\bM$ has the discrepancy property with parameter $\delta>0$, $\kappa_1>1, \kappa_2\geq 0$, denoted by $\textnormal{\bf{DP}}(\delta,\kappa_1,\kappa_2)$, if for all non-empty $S,T\subset [n]$, at least one of the following hold:
\begin{enumerate}
    \item $e_{\bM}(S,T)\leq \kappa_1 \delta |S| |T|;$
    \item  $e_{\bM}(S,T) \cdot \log \Bigg(\frac{e_{\bM}(S,T)}{\delta |S|\cdot |T|}\Bigg)\leq \kappa_2 (|S|\vee |T|)\cdot \log \Bigg(\frac{en}{|S|\vee |T|}\Bigg) $.
\end{enumerate}
\end{definition}

The following shows that if a symmetric random matrix $\bA$ satisfies the upper tail property $\textnormal{\textbf{UUTP}}(c_0,\gamma_0)$ with parameter $c_0>0,\gamma_0\geq 0$, then the discrepancy property holds with high probability.

\begin{lemma}[Lemma 6.4 in \cite{cook2018size}]\label{lem:DPforM}
Let $\bM$ be an $n\times n$ symmetric random matrix with non-negative entries. Assume that for some $\delta>0$, $\E \bM_{ij}\leq \delta$ for all $i,j\in [n]$ and $\bM$ has  $\textnormal{\textbf{UUTP}}(c_0,\gamma_0)$ with parameter $c_0,\gamma_0 >0$. Then for any $K>0$,  the  discrepancy  property $\textnormal{\bf{DP}}(\delta,\kappa_1,\kappa_2)$ holds for $\bM$ with probability at least $1 - n^{-K}$ with 
$
    \kappa_1 = e^2(1+\gamma_0)^2,  \kappa_2 = \frac{2}{c_0}(1+\gamma_0)(K+4).
$
\end{lemma}

When the discrepancy property holds, then deterministically the  contribution from heavy pairs is $O(\sqrt{d})$, as shown in the following lemma.
\begin{lemma}[Lemma 6.6 in \cite{cook2018size}]\label{lem:heavybound}
Let $\bM$ be a non-negative symmetric $n\times n$ matrix with all row sums bounded by $d$. Suppose $\bM$ has $\textnormal{\bf{DP}}(\delta,\kappa_1,\kappa_2)$ with $\delta=Cd/n$ for some $C>0,\kappa_1>1,\kappa_2\geq 0$. Then for any $x\in \mathbb{S}^{n-1}$,
\begin{align*}
    \Bigg| \sum_{(i,j)\in \mathcal{H}(x)} \bM_{ij} \bx_i \bx_j \Bigg|\leq \alpha_0\sqrt{d},
\end{align*}
where
$
    \alpha_0=16+32C(1+\kappa_1)+64\kappa_2  (1+\frac{2}{\kappa_1\log\kappa_1}).
$
\end{lemma}

The next lemma shows $\bA$ has bounded row and column sums with high probability.
\begin{lemma}\label{lem:maxdeg}
For any $K>0$, there is a constant $\alpha_1>0$ such that with probability at least $1-n^{-K}$,
\begin{align}\label{eq:bounded_degree}
    \max_{1\leq i\leq n} \,\sum_{j=1}^n \bA_{ij} \leq \alpha_1 d
\end{align}
with $\alpha_1= 4 + \frac{2(M-1)(1+K)}{3c}$ and $d\geq c\log n$.
\end{lemma}
\begin{proof}
For a fixed $i\in [n]$,
\begin{subequations}
\begin{align*}
   &\sum_{j=1}^n \bA_{ij} =\sum_{m=2}^M\sum_{e\in E_m: i\in e} (m-1)\bcT^{(m)}_e\,, \quad \, \sum_{j=1}^n (\bA_{ij} - \E \bA_{ij}) =\sum_{m=2}^M\sum_{e\in E_m: i\in e} (m-1)\bcW^{(m)}_e\,, \\
    &\sum_{j=1}^n\E \bA_{ij} \leq \sum_{m=2}^M \binom{n}{m-1} \frac{(m-1)d_m}{\binom{n}{m-1}}=d,\\ 
   & \sum_{m=2}^M(m-1)^2\sum_{e\in E_m: i\in e} \E[(\bcW^{(m)}_e)^2] \leq \sum_{m=2}^M (m-1)^2\sum_{e\in E_m: i\in e} \E[\bcT^{(m)}_e]\leq (M-1)d\,. 
\end{align*}
\end{subequations}
Then for $\alpha_1= 4 + \frac{2(M-1)(1+K)}{3c}$, by Bernstein's inequality, with the assumption that $d\geq c\log n$,
\begin{align} 
    \P\Bigg (\sum_{j=1}^n \bA_{ij}\geq \alpha_1 d \Bigg) &\leq \P \Bigg(\sum_{j=1}^n \bA_{ij} - \E \bA_{ij}\geq (\alpha_1-1) d\Bigg) \notag\\
    &\leq \exp \Bigg( -\frac{\frac{1}{2}(\alpha_1-1)^2 d^2}{(M-1)d+\frac{1}{3}(M-1)(\alpha_1-1)d}\Bigg) 
    \leq n^{-\frac{3c(\alpha_1-1)^2}{(M-1)(2\alpha_1+4)}}\leq n^{-1-K.}\label{eq:K}
\end{align}
Taking a union bound over $i\in [n]$, then \Cref{eq:bounded_degree} holds with probability $1-n^{-K}$.
\end{proof}

Now we are ready to obtain \Cref{eq:Osqrtd}.
\begin{lemma}
For any $K>0$,  there is a constant $\beta$ depending on $K, c, M$ such that with probability at least $1-2n^{-K}$,
 \begin{align}\label{eq:betasqrtd}
  \Bigg|\sum_{(i,j)\in \mathcal H(x)}  \bA_{ij}\bx_i \bx_j\Bigg|\leq \beta \sqrt{d}.   
 \end{align}
\end{lemma}
\begin{proof}
By Lemma \ref{lem:UUTPforA}, $\bA$ satisfies $\textnormal{\bf{UUTP}}(\frac{1}{M(M-1)},0)$. From \Cref{eq:EA} and Lemma \ref{lem:DPforM}, the property $\textnormal{\bf{DP}}(\delta, \kappa_1,\kappa_2)$ holds for $\bA$ with probability at least $1-n^{-K}$ with 
\begin{align*}
    \delta = \frac{d}{n}, \quad \kappa_1=e^2,\quad  \kappa_2=2M(M-1)(K+4).
\end{align*}
Let $\mathcal{E}_1$ be the event that $\textnormal{\textbf{DP}}(\delta, \kappa_1,\kappa_2)$ holds for $\bA$. Let $\mathcal{E}_2$ be the event that all row sums of $\bA$ are bounded by $\alpha_1 d$. Then
$
    \P (\mathcal{E}_1 \cap \mathcal{E}_2)\geq 1-2n^{-K}.
$
On the event $\mathcal{E}_1 \cap \mathcal{E}_2$, by Lemma \ref{lem:heavybound}, \Cref{eq:betasqrtd} holds with $
    \beta= \alpha_0\alpha_1,
$
where 
 \begin{align*}
    &\alpha_0= 16+32(1+e^2)+128M(M-1)(K+4) (1+e^{-2}), 
    \quad \alpha_1= 4+\frac{2(M-1)(1+K)}{3c}.
\end{align*} 

\end{proof}

\subsection{Proof of Theorem \ref{thm:concentration}}
\begin{proof}
From \eqref{eq:lightbound}, with probability at least $1-2e^{-n}$, the contribution from light pairs in \eqref{eq:contribution} is bounded by  $2\alpha \sqrt{d}$ with $\alpha=5M(M-1)$. From \eqref{eq:heavyexp} and  \eqref{eq:betasqrtd}, with probability at least $1-2n^{-K}$, the contribution from heavy pairs in \eqref{eq:contribution} is bounded by $2\sqrt{d}+2\beta\sqrt{d}$.  Therefore with probability at least $1-2e^{-n}-2n^{-K}$, 
\[
    \|\bA - \E \bA \|\leq C_M \sqrt{d}, 
\]
where $C_M$ is a constant depending only on $c, K, M$ such that
$
    C_M= 2(\alpha+1+\beta).
$ 
In particular, we can take $\alpha=5M(M-1),$ $\beta=512M(M-1)(K+5)\left( 2+\frac{(M-1)(1+K)}{c}\right)$, and
$
    C_M=512M(M-1)(K+6)\left( 2+\frac{(M-1)(1+K)}{c}\right).
$
This finishes the proof of Theorem \ref{thm:concentration}.
\end{proof}

\subsection{Proof of Theorem \ref{thm:regularization}}

Let $\cS\subset [n]$ be any given subset.  From \eqref{eq:lightbound}, with probability at least $1-2e^{-n}$,
\begin{align}\label{eq:S}
    \sup_{\bx\in \mathcal{N} }  \Bigg| \sum_{(i,j)\in \mathcal{L}(\bx)}{\Bigg(\bA_{\cS} - \E \bA_{\cS}\Bigg)}_{ij} \bx_i \bx_j\Bigg|\leq 5M(M-1)\sqrt{d}.
\end{align}
Since there are at most $2^n$ many choices for $\cS$, by taking a union bound, with probability at least $1-2(e/2)^{-n}$, we have for all $\cS\subset [n]$, \Cref{eq:S} holds. In particular, by taking $\cS = \cI = \{i\in [n]: \textnormal{row}(i)\leq \tau d\}$, with probability at least $1-2(e/2)^{-n}$, we have
\begin{align}\label{eq:reg1}
    \sup_{\bx\in \mathcal{N}}  \Bigg| \sum_{(i,j)\in \mathcal{L}(\bx)}[(\bA - \E \bA)_{\cI}]_{ij} \bx_i \bx_j\Bigg|\leq 5M(M-1)\sqrt{d}.
\end{align}
Similar to \Cref{eq:heavyexp}, deterministically,
\begin{align}\label{eq:reg2}
     \Bigg|\sum_{(i,j)\in \mathcal{H}(\bx)} [(\E \bA)_{\cI}]_{ij}\bx_i \bx_j\Bigg|\leq (M-1)\sqrt{d}.
\end{align}
Next we show the contribution from heavy pairs for $\bA_{\cI}$ is bounded.
\begin{lemma}
For any $K>0$,  there is a constant $\beta_{\tau}$ depending on $K, c, M,\tau$ such that with probability at least $1-n^{-K}$,
 \begin{align}\label{eq:betasqrtd_tau}
  \Bigg|\sum_{(i,j)\in \mathcal{H}(\bx)}  [(\bA)_{\cI}]_{ij}\bx_i \bx_j\Bigg|\leq \beta_{\tau} \sqrt{d}.   
 \end{align}
\end{lemma}
\begin{proof}
Note that $\bA$ satisfies $\textnormal{\bf{UUTP}}\left(\frac{1}{M(M-1)},0 \right)$ from Lemma \ref{lem:UUTPforA}. Then from Lemma \ref{lem:DPforM}, with probability at least $1-n^{-K}$, $\textnormal{\bf{DP}}(\delta,\kappa_1,\kappa_2)$ holds for $\bA$ with
\begin{align*}
    \delta =\frac{d}{n}, \quad \kappa_1=e^2,\quad  \kappa_2=2M(M-1)(K+4).
\end{align*}
The $\textnormal{\bf{DP}}(\delta,\kappa_1,\kappa_2)$ property holds for $\bA_{\cI}$ as well, since $\bA_{\cI}$ is obtained from $\bA$ by restricting to $\cI$. Note that all row sums in $\bA_{\cI}$ are bounded by $\tau d$. By \Cref{lem:heavybound},
\begin{align}\label{eq:reg3}
    \Bigg| \sum_{(i,j)\in \mathcal{H}(\bx)} [\bA_{\cI}]_{ij}\bx_i \bx_j \Bigg|\leq \alpha_0\sqrt{\tau d},
\end{align}
where we can take
$
    \alpha_0=16+\frac{32}{\tau}(1+e^2)+128M(M-1)(K+4)\left (1+\frac{1}{e^2}\right).
$
\end{proof}
We can then take $\beta_{\tau}=\alpha_0\sqrt{\tau}$ in \Cref{eq:betasqrtd_tau}. Therefore, combining \Cref{eq:reg1}, \Cref{eq:reg2}, \Cref{eq:reg3}, with probability at least $1-2(e/2)^{-n}-n^{-K}$, there exists a constant $C_{\tau}$ depending only on $\tau, M, K$ such that 
$
    \|(\bA-\E \bA)_{\cI}\|\leq C_{\tau} \sqrt{d},
$
where $C_{\tau}=2( (5M+1)(M-1)+\alpha_0\sqrt{\tau}).$
This finishes the proof of Theorem \ref{thm:regularization}.
\section{Technical Lemmas}

\subsection{Proof of \Cref{lem:general_eigen_value_approximation}}
\begin{proof}
By Weyl's inequality \ref{lem:Weyl_inequality}, the difference between eigenvalues of $\widetilde{\E \bA}$ and $\E \bA$ can be upper bounded by
\begin{subequations}
\begin{align*}
    |\lambda_i( \widetilde{\E \bA} ) - \lambda_i(\E \bA)| & \leq \|\widetilde{\E \bA} - \E \bA \|_2 \leq \|\widetilde{\E \bA} - \E \bA \|_{\frob}\\
    &\leq \left[2k \cdot \frac{n}{k} \cdot \sqrt{n}\log(n) \cdot (\alpha - \beta)^2 \right]^{1/2}= O \Big(n^{3/4}\log^{1/2}(n)(\alpha - \beta)\Big).
\end{align*}
\end{subequations}
The lemma follows, as $\lambda_i(\E \bA) = \Omega \left ( n (\alpha - \beta) \right ) $ for all $1 \leq i \leq k$. 
\end{proof}

\subsection{Proof of \Cref{lem:singular_value_approximation}}\label{subsec:prf_singular_value_approximation}
\begin{proof}
We first compute the singular values of $\overline{\bB}_1$. From \Cref{eqn:bar_B1}, the rank of matrix $\overline{\bB}_1$ is $k$, and 
% the matrix product $\overline{\bB}_1^{\sT}\overline{\bB}_1$ is written as
% \begin{align*}
%     \overline{\bB}_1^{\sT}\overline{\bB}_1 &= \begin{bmatrix}
%       \frac{n}{2k}[\overline{\alpha}^2 + (k-1)\overline{\beta}^2]\bJ_{\frac{n}{4k}}& \cdots & \frac{n}{2k}[2\overline{\alpha}\overline{\beta} + (k-2)\overline{\beta}^2]\bJ_{\frac{n}{4k}} \\
%       \vdots & \ddots & \vdots\\
%         \frac{n}{2k}[2\overline{\alpha} \overline{\beta} + (k-2)\overline{\beta}^2]\bJ_{\frac{n}{4k}} & \cdots &\frac{n}{2k}[\overline{\alpha}^2 + (k-1)\overline{\beta}^2]\bJ_{\frac{n}{4k}}
%     \end{bmatrix}\,,
% \end{align*}
% where $\overline{\alpha}$, $\overline{\beta}$ are defined in \Cref{eqn:bar_alphabeta} respectively, and $\bJ_{\frac{n}{4k}}$ denotes the $n/(4k)\times n/(4k)$ matrix with all ones. The eigenvalues of $\overline{\bB}_1^{\sT}\overline{\bB}_1$ are
% \begin{subequations}
% \begin{align*}
%     \lambda_1(\overline{\bB}_1^{\sT}\overline{\bB}_1) =&\,\frac{n}{4k}\left\{ \frac{n[\overline{\alpha}^2 + (k-1)\overline{\beta}^2]}{2k} + \frac{(k-1)n[2\overline{\alpha}\overline{\beta} + (k-2)\overline{\beta}^2]}{2k}\right\} = \frac{n^2[\overline{\alpha} + (k-1)\overline{\beta}]^2 }{8k^2}\,,\\
%     \lambda_i(\overline{\bB}_1^{\sT}\overline{\bB}_1) = &\,\frac{n}{4k}\left\{ \frac{n[ \overline{\alpha}^2 + (k-1)\overline{\beta}^2]}{2k} -\frac{n[2\overline{\alpha}\overline{\beta} + (k-2)\overline{\beta}^2]}{2k} \right\} = \frac{n^2(\overline{\alpha} -\overline{\beta})^2}{8k^2}\,, \quad 2 \leq i \leq k\,,\\
%     \lambda_i(\overline{\bB}_1^{\sT}\overline{\bB}_1) = &\, 0\,\quad \quad \quad \quad \quad k+1 \leq i \leq n\,.
% \end{align*}
% \end{subequations}
% associated with eigenvectors
% \newcommand\MyLBrace[2]{\left.\rule{0pt}{#1}\right\}\text{#2}}
% \begin{align*}
%     \bv_1 = \begin{bmatrix}
%         \left(\begin{smallmatrix}1\\ \vdots\\ \\ 1\end{smallmatrix}\right)\\\\
%         \left(\begin{smallmatrix}1\\ \vdots\\ \\1\end{smallmatrix}\right)\\
%         \vdots\\
%         \left(\begin{smallmatrix}1\\ \vdots\\ \\1\end{smallmatrix}\right)\\
%       \end{bmatrix}\!\!\!\!\!\!
%       \begin{array}{l}
%           \MyLBrace{4ex}{(n/4k)}  \\
%           \\
%           \MyLBrace{4ex}{(n/4k)}  \\
%           \\
%           \\
%           \MyLBrace{4ex}{(n/4k)} 
%       \end{array}\,, \quad
%       \bv_2 = \begin{bmatrix}
%         \left(\begin{smallmatrix}1\\ \vdots\\ \\ 1\end{smallmatrix}\right)\\\\
%         \left(\begin{smallmatrix}-1\\ \vdots\\ \\-1\end{smallmatrix}\right)\\
%         \vdots\\
%         \left(\begin{smallmatrix}0\\ \vdots\\ \\0\end{smallmatrix}\right)\\
%       \end{bmatrix}\!\!\!\!\!\!
%       \begin{array}{l}
%           \MyLBrace{4ex}{(n/4k)}  \\
%           \\
%           \MyLBrace{4ex}{(n/4k)}  \\
%           \\
%           \\
%           \MyLBrace{4ex}{(n/4k)} 
%       \end{array}\,, \cdots, \bv_k = \begin{bmatrix}
%         \left(\begin{smallmatrix}1\\ \vdots\\ \\ 1\end{smallmatrix}\right)\\\\
%         \left(\begin{smallmatrix}0\\ \vdots\\ \\0\end{smallmatrix}\right)\\
%         \vdots\\
%         \left(\begin{smallmatrix}-1\\ \vdots\\ \\-1\end{smallmatrix}\right)\\
%       \end{bmatrix}\!\!\!\!\!\!
%       \begin{array}{l}
%           \MyLBrace{4ex}{(n/4k)}\\
%           \\
%           \MyLBrace{4ex}{(n/4k)}\\
%           \\
%           \\
%           \MyLBrace{4ex}{(n/4k)} 
%       \end{array}\,.
% \end{align*}
 the least non-trivial singular value  of $\overline{\bB}_1$ is
\begin{align*}
    \sigma_k (\overline{\bB}_1) = \frac{n}{2\sqrt{2}k}(\overline{\alpha} - \overline{\beta}) = \frac{n}{2\sqrt{2}k}\sum_{m\in \mathcal{M}} \binom{\frac{3n}{4k} - 2}{m-2} \frac{a_m - b_m}{ \binom{n}{m-1}}\,,
\end{align*}
where $\mathcal{M}$ is obtained from \Cref{alg:parameter_preprocessing}. By the definition of $\overline{\bA}_1$ in \Cref{eqn:bar_A1A2}, the least non-trivial singular value  of $\overline{\bA}_1$ is
\begin{align*}
    \sigma_k (\overline{\bA}_1) = \sigma_k (\overline{\bB}_1) = \frac{n}{2\sqrt{2}k}(\overline{\alpha} - \overline{\beta}) = \frac{n}{2\sqrt{2}k}\sum_{m\in \mathcal{M}} \binom{\frac{3n}{4k} - 2}{m-2} \frac{a_m - b_m}{ \binom{n}{m-1}}\,.
\end{align*}
Recall that $n_i$, defined in \Cref{eqn:dimension_ZcapVi}, denotes the number of vertices in $Z\cap V_i$, which can be written as  $n_i = \sum_{v\in V_i} \mathbf{1}_{\{ v \in Z\} }.$ 
% The expectation of $n_i$ is $\E n_i = n/(2k)$.
By Hoeffding's Lemma \ref{lem:Hoeffding}, 
\begin{align*}
    \P\left( \bigg|n_i - \frac{n}{2k}\bigg| \geq \sqrt{n}\log(n) \right) \leq 2\exp \left( - k\log^2 (n)\right)\,. 
\end{align*}
Similarly, $n^{\prime}_{i}$, defined in \Cref{eqn:dimension_Y1capVi}, satisfies
% denotes the number of vertices in $Y_1\cap V_i$, which can be written as $n^{\prime}_{i} = |Y_1 \cap V_i| = \sum_{v\in V_i} \mathbf{1}_{ \{v\in Y_1\} }$ and 
% with $\E n^{\prime}_{i} = n/(4k)$. By Hoeffding's \Cref{lem:Hoeffding}, 
\begin{align}
    \P\left( \bigg|n^{\prime}_{i} - \frac{n}{4k}\bigg| \geq \sqrt{n}\log(n) \right) \leq 2\exp \left( - k\log^2 (n)\right)\,. \notag
\end{align}
As defined in \Cref{eqn:tilde_B1} and \Cref{eqn:bar_B1}, both $\widetilde{\bB}_1$ and $\overline{\bB}_1$ are deterministic block matrices. Then with probability at least $1 - 2k\exp \left( - k\log^2 (n)\right)$, the dimensions of each block inside $\widetilde{\bB}_1$ and $\overline{\bB}_1$ are approximately the same, with deviations up to $\sqrt{n}\log(n)$.
% Consequently, the matrix $\widetilde{\bA}_1$, which was defined in \Cref{eqn:tilde_A1A2}, can be treated as a perturbed version of $\overline{\bA}_1$. 
By Weyl's Lemma \ref{lem:Weyl_inequality}, for any $i\in [k]$,
\begin{subequations}
\begin{align*}
    |\sigma_i(\overline{\bB}_1) - \sigma_i(\widetilde{\bB}_1)| &=  |\sigma_i(\overline{\bA}_1) - \sigma_i(\widetilde{\bA}_1)|\leq \|\overline{\bA}_1 - \widetilde{\bA}_1\|_2 \leq \|\overline{\bA}_1 - \widetilde{\bA}_1\|_{\frob}\\
    &\leq \left[2k \cdot \frac{n}{k} \cdot \sqrt{n}\log(n) \cdot (\overline{\alpha} - \overline{\beta})^2 \right]^{1/2}= O \left(n^{3/4}\log^{1/2}(n) \cdot (\overline{\alpha} - \overline{\beta})\right).
\end{align*}
\end{subequations}
As a result, with probability at least $1 - 2k\exp( -k \log^2(n))$, we have
\begin{align*}
    \frac{|\sigma_k(\overline{\bA}_1) - \sigma_k(\widetilde{\bA}_1)|}{\sigma_k(\overline{\bA}_1)} = \frac{|\sigma_k(\overline{\bB}_1) - \sigma_k(\widetilde{\bB}_1)|}{\sigma_k(\overline{\bB}_1)} = O\left(n^{-1/4}\log^{1/2}(n)\right).
\end{align*}
% Therefore, singular values of $\widetilde{\bA}_1$ can be approximated by singular values of $\overline{\bA}_1$ with vanishing errors.
\end{proof}



\subsection{Proof of Lemma \ref{lem:high_degree_vertices_k}}
\begin{proof}\label{prf:high_degree_vertices_k}


Without loss of generality, we can assume $\mathcal M=\{2,\dots, M\}$. If $\mathcal M$ is a subset of $\{2,\dots, M\}$, we can take $a_m=b_m=0$ for $m\not\in \mathcal M$. Note that in fact, if the best SNR is obtained when $\mathcal{M}$ is a strict subset, we can substitute $\mathcal{M}_{\max}$ for $M$. 




Let $X\subset V$ be a subset of vertices in hypergraph $H = (V, E)$ with size $|X| = cn$ for some $c\in (0,1)$ to be decided later.
Suppose $X$ is a set of vertices with high degrees that we want to zero out. 
% Our goal is to prove the following statement:
% \textit{The probability that all vertices in $X$ has degree higher than $20 d$ is small.}
 We first count the $m$-uniform hyperedges on $X$ separately, then weight them by $(m-1)$, and finally sum over $m$ to compute the row sums in $\bA$  corresponding to each vertex in $X$.  Let $E_m(X)$ denote the set of $m$-uniform hyperedges with all vertices located in $X$, and $E_m (X^{c})$ denote the set of $m$-uniform hyperedges with all vertices in $X^{c}=V\setminus X$, respectively. Let $E_m(X, X^{c})$ denote the set of $m$-uniform hyperedges with at least $1$ endpoint in $X$ and $1$ endpoint in $X^{c}$. 
% The relationship between the number of total degrees and the number $m$-uniform hyperedges in $X$ satisfies
%     \begin{equation}\label{eqn:m_uniform_hyperedges_in_X}
%         m|E_m(X)| + |E_m(X, X^{c})| \leq \sum_{v \in X} \mathrm{deg}(v) \leq m|E_m(X)| + (m-1)|E_m(X, X^{c})|\,.
%     \end{equation}
% The upper bound is achieved when there are exactly $m - 1$ endpoints in $X$ and $1$ endpoint in $V \setminus X$ for each hyperedge in $E(X_{m}, X_{m}^{c})$, while the lower bound is achieved when there are exactly $m-1$ endpoints in $V \setminus X$ and $1$ endpoint in $X$, respectively. 
% By symmetry, we have
%     \begin{equation}\label{eqn:m_uniform_hyperedges_in_Xc}
%         m|E_m(X^{c})| + |E_m(X, X^{c})| \leq \sum_{v \in V \setminus X} \mathrm{deg}(v) \leq m|E_m(X^{c})| + (m-1)|E_m(X, X^{c})|\,.
%     \end{equation}
% Obviously, the upper bounds in \Cref{eqn:m_uniform_hyperedges_in_X} and \Cref{eqn:m_uniform_hyperedges_in_Xc} can't be achieved simultaneously. If the upper bound in \Cref{eqn:m_uniform_hyperedges_in_X} is achieved, then the lower bound in \Cref{eqn:m_uniform_hyperedges_in_Xc} must be taken. 
% By summing over $m$, 
The relationship between total row sums and the number of non-uniform hyperedges in vertex set $X$ can be expressed as
    \begin{align}\label{eqn:non_uniform_hyperedges_in_X}
        %\sum_{m = 2}^{M} (m-1)\Big( m|E_m(X)| + |E_m(X, X^{c})| \Big) \leq 
        \sum_{v \in X} \mathrm{row}(v) \leq &\, \sum_{m = 2}^{M} (m-1)\Big( m|E_m(X)| + (m-1)|E_m(X, X^{c})| \Big) %\leq &\, \sum_{m = 2}^{M} m\Big( |E_m(X)| + |E_m(X, X^{c})| \Big).
    \end{align} 
If the row sum of each vertex $v \in X$ is at least $20Md$, where $d = \sum_{m=2}^{M} (m-1)a_m$, it follows that
\begin{equation}\label{eq:lowerboundEM}
\sum_{m = 2}^{M} (m-1)\Big(m |E_m(X)| + (m-1)|E_m(X, X^{c})| \Big) \geq cn\cdot(20Md)\,.
\end{equation}
Then either 
\begin{align*}
    \sum_{m = 2}^{M} m(m-1) |E_m(X)| \geq 4Mcnd, \quad \text{or} \quad  \sum_{m = 2}^{M} (m-1)^2|E_m(X, X^{c})| \geq 16Mcnd.  \notag 
\end{align*}

\subsubsection{Concentration of $\sum_{m=2}^{M} m(m-1)|E_m(X)|$.} Recall that $\big|E_m(X)\big|$ denotes the number of $m$-uniform hyperedges with all vertices located in $X$, which can be viewed as the sum of independent Bernoulli random variables $T_{e}^{(a_m)}$ and $T_{e}^{(b_m)}$ given by
    \begin{equation}\label{eq:EmXBernoulli}
        T_{e}^{(a_m)}\sim \mathrm{Bernoulli}\left( \frac{a_{m}}{ {n \choose m-1} }\right), \quad T_{e}^{(b_m)}\sim \mathrm{Bernoulli}\left( \frac{b_{m}}{ {n \choose m-1} }\right)\,.
    \end{equation}
% Here $ T_{e}^{(a_m)}$ denotes the presence of the hyperedge $e$ within same cluster and  $ T_{e}^{(b_m)}$ denotes the presence of the hyperedge $e$ between different clusters.
Let \{$V_{1}, \dots, V_{k}$\} be the true partition of $V$. Suppose that there are $\eta_{i}cn$ vertices in block $V_{i}\cap X$ for each $i\in [k]$ with restriction $\sum_{i=1}^{k}\eta_{i} = 1$, then  $\big|E_m(X)\big|$ can be written as
\begin{align}
   \big|E_m(X)\big| = \sum_{e\in E_{m}(X, a_m)} T_e^{(a_m)} \, + \sum_{e\in E_m(X, b_m)} T_e^{(b_m)}\,, \notag 
\end{align}
where $E_{m}(X, a_m) := \cup_{i=1}^{k}E_{m}(V_i\cap X)$ denotes the union for sets of hyperedges with all vertices in the same block $V_i\cap X$ for some $i\in [k]$, and
\begin{align}
    E_m(X, b_m) :=\, E_{m}(X) \setminus E_{m}(X, a_m) = E_{m}(X) \setminus \Big( \cup_i^{k} E_{m}(V_i\cap X) \Big) \notag 
\end{align}
denotes the set of hyperedges with vertices crossing different $V_i\cap X$. %Since 
%\begin{align}
    %|E_m(X, a_m)| = \sum_{i=1}^{k} |E_{m}(V_i\cap X)|= \sum_{i=1}^{k} {\eta_{i}cn \choose m}, \quad |E_m(X, b_m)|= { cn \choose m} - \sum_{i=1}^{k} {\eta_{i}cn \choose m}\,,
%\end{align}
We can compute the expectation of $\big|E_m(X)\big|$ as 
    \begin{equation}\label{eq:ExpE_m}
        \E |E_m(X)| = \sum_{i=1}^{k} {\eta_{i}cn \choose m} \frac{a_m - b_m }{ {n \choose m-1} } + { cn \choose m} \frac{b_{m}}{ {n \choose m-1} }.
    \end{equation}
Then
\begin{align}\label{eqn:expectation_EX}
    \sum_{m=2}^{M} m(m-1) \cdot \E |E_m(X)| & = \sum_{m=2}^{M} m(m-1) \Bigg[ \sum_{i=1}^{k} {\eta_{i}cn \choose m} \frac{a_m - b_m }{ {n \choose m-1} } + { cn \choose m} \frac{b_{m}}{ {n \choose m-1} } \Bigg].
\end{align}
As $\sum\limits_{i=1}^k \eta_i=1$, it follows that 
$
\sum_{i=1}^k \binom{\eta_i c n}{m} \leq \binom{cn}{m}
$ by induction, thus
\[
 \frac{a_m - b_m }{ {n \choose m-1} } \sum_{i=1}^{k} {\eta_{i}cn \choose m}  +  \frac{b_{m}}{ {n \choose m-1} } { cn \choose m} =  \frac{a_m}{ {n \choose m-1} } \sum_{i=1}^{k} {\eta_{i}cn \choose m}  + \frac{b_{m}}{ {n \choose m-1} } \left ( { cn \choose m} - \sum_{i=1}^k {\eta_{i}cn \choose m} \right)  
\]
where both terms on the right are  positive numbers. Using this and taking $b_m = a_m$, we obtain the following upper bound for all $n$,
\begin{align*}
 \sum_{m=2}^{M} m(m-1) \E |E_m(X)|\leq \sum_{m=2}^M m(m-1)\binom{cn}{m} \frac{a_m}{\binom{n}{m-1}} \leq cn\sum_{m=2}^M (m-1)a_m =cnd \,.
\end{align*}
% By taking $b_m=a_m$ and $b_m=0$, respectively, we obtain the upper and lower bounds of \eqref{eqn:expectation_EX_Xc} as 
% \begin{align}
%   \sum_{m=2}^M m \sum_{i=1}^k \binom{\eta_icn}{m} \frac{a_m}{\binom{n}{m-1}}  \leq \sum_{m=2}^{M} m \mu_{|E_m(X)|}\leq \sum_{m=2}^M m\binom{cn}{m} \frac{a_m}{\binom{n}{m-1}}.
% \end{align}
% We can further simplify the lower bound by taking $\eta_i=1/k$. Then
% \begin{align}
%   \sum_{m=2}^M m k \binom{cn/k}{m} \frac{a_m}{\binom{n}{m-1}}  \leq \sum_{m=2}^{M} m \mu_{|E_m(X)|}\leq \sum_{m=2}^M m\binom{cn}{m} \frac{a_m}{\binom{n}{m-1}}.
% \end{align}
%And for sufficiently large $n$,
% \begin{align}
%  \frac{1}{2}kc^M nd \leq \frac{1}{2}n\sum_{m=2}^M k^{m-1} c^m a_m \leq \sum_{m=2}^{M} m \mu_{|E_m(X)|}\leq 2n\sum_{m=2}^M c^m a_m \leq 2c^2 nd.
% \end{align}
Note that $\sum_{m = 2}^{M} m(m-1) |E_m(X)|$ is a weighted sum of independent Bernoulli random variables (corresponding to hyperedges), each upper bounded by $M^2$. Also, its variance is bounded by
\begin{align}
    \sigma^2:= \,& \Var\left(\sum_{m = 2}^{M} m(m-1) |E_m(X)|\right)=\sum_{m=2}^M m^2(m-1)^2 \Var\left( |E_m(X)|\right) \notag \\
    \leq\,& \sum_{m=2}^M m^2(m-1)^2\E |E_m(X)|\leq M^2cnd. \notag 
\end{align}
We can apply Bernstein's Lemma \ref{lem:Bernstein} and obtain
\begin{align}
  \mathbb P\left( \sum_{m = 2}^{M} m(m-1) |E_m(X)| \geq 4Mcnd \right)\notag \leq \,& \mathbb P\left( \sum_{m = 2}^{M} m(m-1)( |E_m(X)| -\E|E_m(X)| )\geq 3Mcnd \right)\notag\\
\leq \,& \exp \left(-\frac{(3Mcnd)^2}{M^2cnd + M^2cnd/3} \right) \leq \exp(-6cnd)\,.\label{eq:chernoff111}
\end{align}

%we can apply Bennett's inequality \ref{lem:Bennett}. Take $K=M$ in \ref{lem:Bennett}. And note that
%\begin{align}
    %\sigma^2:= \,& \Var\left(\sum_{m = 2}^{M} m |E_m(X)|\right)=\sum_{m=2}^M m^2 \Var\left( |E_m(X)|\right)\\
    %\leq\,& \sum_{m=2}^M m^2\E |E_m(X)|\leq 2Mcnd.
%\end{align}
%which implies for sufficiently large $n$,
%\begin{align}
 % \sigma^2 &\leq \sum_{m=2}^M m^2\mu_{|E_m(X)|}\leq 2Mcnd.
%\end{align}

%We obtain
%\begin{align}
  %& \mathbb P\left( \sum_{m = 2}^{M} m |E_m(X)| \geq 4cnd \right)
%\leq  \mathbb P\left( \sum_{m = 2}^{M} m( |E_m(X)| -\mu_{|E_m(X)|})\geq 3cnd \right)\\
%\leq & \exp(-\frac{\sigma^2}{M^2}h(3M^2cnd/\sigma^2)
%\leq \exp(-cnd h(3))\leq \exp(-2cnd).\label{eq:chernoff111}
%\end{align}
%where we use the fact that $\frac{1}{x} h(x)=\frac{1 + x}{x}\log(1 + x) - 1$ is increasing when $x>0$ in the last line.

\subsubsection{Concentration of  $\sum_{m=2}^{M} (m-1)^2|E_m(X, X^{c})|$.}

For any finite set $S$, let $[S]^{j}$ denote the family of $j$-subsets of $S$, i.e., $[S]^{j} = \{Z| Z\subseteq S, |Z| = j \}$. Let $E_m([Y]^{j}, [Z]^{m-j})$ denote the set of $m$-hyperedges, where $j$ vertices are from $Y$ and $m-j$ vertices are from $Z$ within each $m$-hyperedge. We want to count the number of $m$-hyperedges between $X$ and $X^c$, according to the number of vertices located in $X^c$ within each $m$-hyperedge. Suppose that there are $j$ vertices from $X^c$ within each $m$-hyperedge for some $1\leq j \leq m-1$.
\begin{enumerate}[label=(\roman*)]
     \item Assume that all those $j$ vertices are in the same $[V_{i}\setminus X]^{j}$. If the remaining $m-j$ vertices are from $[V_{i}\cap X]^{m-j}$, then this $m$-hyperedge is connected with probability $a_m / \binom{n}{m-1}$, otherwise $b_m / \binom{n}{m-1}$. The number of this type $m$-hyperedges can be written as
    \begin{align}
        \sum_{i=1}^{k} \left[ \sum_{e\in \mathcal{E}^{(a_m)}_{j, i}} T_{e}^{(a_m)} + \sum_{e\in \mathcal{E}^{(b_m)}_{j, i}} T_{e}^{(b_m)} \right]\,, \notag 
    \end{align} 
    where $\mathcal{E}^{(a_m)}_{j, i}:= E_m([V_i \cap X^c]^{j}, [V_i\cap X]^{m-j})$, and
        \begin{align}
            \mathcal{E}^{(b_m)}_{j, i} := E_m \Big([V_i\cap X^c]^{j}, \,\, [X]^{m-j} \setminus [V_i\cap X]^{m-j} \Big) \notag 
        \end{align}
    denotes the set $m$-hyperedges with $j$ vertices in $[V_i\cap X^c]^{j}$ and the remaining $m-j$ vertices  in $[X]^{j}\setminus [V_i\cap X]^{j}$. We compute all possible choices and upper bound the  cardinality of $\mathcal{E}^{(a_m)}_{j, i}$ and $\mathcal{E}^{(b_m)}_{j, i}$ by
    \begin{align}
            \big|\mathcal{E}^{(a_m)}_{j, i} \big| \leq &\, \binom{ (\frac{1}{k} - \eta_{i}c)n}{j} \binom{\eta_{i}cn}{m - j}\,,\quad \big|\mathcal{E}^{(b_m)}_{j, i} \big| \leq \binom{ (\frac{1}{k} - \eta_{i}c)n}{j} \left[ \binom{cn}{m - j} - \binom{\eta_{i}cn}{m - j} \right]\,. \notag 
        \end{align}
        
    \item If those $j$ vertices in $[V\setminus X]^{j}$ are not in the same $[V_i \cap X]^{j}$ (which only happens $j \geq 2$), then the number of this type hyperedges can be written as $\sum_{e\in \mathcal{E}^{(b_m)}_{j}} T_{e}^{(b_m)}$, where 
    \begin{subequations}
        \begin{align*}
            \mathcal{E}^{(b_m)}_{j} :=&\, E_m \Big([V\setminus X]^{j}\setminus \big( \cup_{i=1}^{k}[V_i\setminus X]^{j} \big), \,\,\, [X]^{m-j}\Big)\,,\\
            \big| \mathcal{E}^{(b_m)}_{j} \big| \leq &\, \left[ \binom{(1-c)n}{j} -\sum_{i=1}^{k} \binom{( \frac{1}{k} - \eta_{i}c) n}{j} \right] \binom{cn}{m-j}\,.
        \end{align*}
    \end{subequations}
\end{enumerate}
Therefore, $|E_m(X, X^c)|$ can be written as a sum of independent Bernoulli random variables,
\begin{align}\label{eq:EXmcBernoulli}
    |E_m(X, X^c)| = \sum_{j=1}^{m-1}\sum_{i=1}^{k} \left[ \sum_{e\in \mathcal{E}^{(a_m)}_{j, i}} T_{e}^{(a_m)} + \sum_{e\in \mathcal{E}^{(b_m)}_{j, i}} T_{e}^{(b_m)} \right] + \sum_{j=2}^{m-1}\sum_{e\in \mathcal{E}^{(b_m)}_{j}} T_{e}^{(b_m)}\,.
\end{align}
%Note that $\binom{(1-c)n}{1} = \sum_{i=1}^{k}\binom{(1/k - \eta_i c)n}{1}$, and so $\big| \mathcal{E}^{(b_m)}_{1} \big| = 0$. Thus we write $\E (|E_m(X, X^c)|)$ as
%\begin{align}
    %\E(|E_m(X, X^c)|) =\E \left ( \sum_{j=1}^{m-1} \left\{ \sum_{i=1}^{k} \left[ \sum_{e\in \mathcal{E}^{(a_m)}_{j, i}} T_{e}^{(a_m)} + \sum_{e\in \mathcal{E}^{(b_m)}_{j, i}} T_{e}^{(b_m)} \right] + \sum_{e\in \mathcal{E}^{(b_m)}_{j}} T_{e}^{(b_m)} \right\}\,\right ).
%\end{align}
Then the expectation can be rewritten as
\begin{align}
    &\, \E \left(|E_m(X, X^c)| \right) \label{eqn:regrouping_EXXc}\\
    =&\, \sum_{j=1}^{m-1} \sum_{i=1}^{k} \binom{( \frac{1}{k} - \eta_{i}c) n}{j} \Bigg\{ \binom{\eta_{i}cn}{m-j} \frac{a_m}{ \binom{n}{m-1} } + \bigg[ \binom{cn}{m-j} - \binom{\eta_{i}cn}{m-j} \bigg] \frac{b_m}{ \binom{n}{m-1} } \Bigg\} \notag \\
    \,& + \sum_{j=1}^{m-1} \Bigg[ {(1 - c) n \choose j} - \sum_{i=1}^{k}{( \frac{1}{k} - \eta_{i}c) n \choose j} \Bigg] {cn \choose m - j} \frac{b_m}{ \binom{n}{m-1} } \notag \\ 
    =\,& \sum_{j=1}^{m-1} \sum_{i=1}^{k}\binom{( \frac{1}{k} - \eta_{i}c) n}{j} \binom{\eta_{i}cn}{m-j} \frac{a_m - b_m}{ \binom{n}{m-1} } + \sum_{j=1}^{m-1} \binom{(1-c)n}{j} \binom{cn}{m-j} \frac{b_m}{ \binom{n}{m-1} }\notag\\
    =&\, \sum_{i=1}^{k} \Bigg[ { \frac{n}{k} \choose m} -  {\eta_{i}cn \choose m} - {(\frac{1}{k} - \eta_{i}c) n \choose m} \Bigg]\frac{a_m - b_m}{{n \choose m-1} } + \Bigg[ {n \choose m} - {cn \choose m} - {(1 - c) n \choose m} \Bigg] \frac{b_m}{{n \choose m-1} }\notag\,,
\end{align}
where we used the fact $\binom{(1-c)n}{1} = \sum_{i=1}^{k}\binom{(1/k - \eta_i c)n}{1}$ in the first equality and Vandermonde's identity ${n_{1} + n_{2} \choose m} = \sum_{j=0}^{m}{n_{1} \choose j}{n_{2} \choose m-j}$ in last equality. Note that 
\[
f_c := {n \choose m} - {cn \choose m} - {(1 - c) n \choose m}
\]
counts the number of subsets of $V$ with $m$ elements such that at least one element belongs to $X$ and at least one element belongs to $X^c$. On the other hand,
\[
g_c = \sum_{i=1}^{k} \left[  { \frac{n}{k} \choose m} -  {\eta_{i}cn \choose m} - {(\frac{1}{k} - \eta_{i}c) n \choose m} \right]\,.
\]
counts the number of subsets of $V$ with $m$ elements such that all elements belong to a single $V_i$, and given such an $i$, that at least one element belongs to $X \cap V_i$ and at least one belongs to $X^c \cap V_i$. 
\begin{figure}[ht]
    \centering
    \includegraphics[width= 0.4\linewidth]{fig/fig1.png}
    \caption{Comparison of $f_c$ and $g_c$}
    \label{fig:comparison_fcgc}
\end{figure}
As \Cref{fig:comparison_fcgc} shows, $g_{c}$ only counts the blue pairs while $f_c$ counts red pairs in addition. By virtue of the fact that there are fewer conditions imposed on the sets included in the count for $f_c$, we must have $f_c \geq g_c$. Thus, rewriting \Cref{eqn:regrouping_EXXc}, we obtain
\begin{align*}
\E (|E_m(X, X^c)|) &\,= g_c \frac{a_m}{{n \choose m-1} } + (f_c-g_c) \frac{b_m}{{n \choose m-1} }\,~.
\end{align*}
Since both terms in the above sum are positive, we can upper bound by taking $a_m=b_m$ to obtain
\begin{align}
    \E \left( |E_m(X, X^c)| \right) \leq f_c  \frac{a_m}{{n \choose m-1} } = \Bigg[ {n \choose m} - {cn \choose m} - {(1 - c) n \choose m} \Bigg] \frac{a_m}{{n \choose m-1} } ~. \notag 
\end{align}
By summing over $m$, the expectation of $\sum_{m=2}^{M}(m-1)^2|E_m(X, X^{c})|$ satisfies
\begin{subequations}
\begin{align}
  \sum_{m=2}^M (m-1)^2\cdot \E \left(|E_m(X, X^c)| \right)&\leq  \sum_{m=2}^M (m-1)^2\Bigg[ {n \choose m} - {cn \choose m} - {(1 - c) n \choose m} \Bigg] \frac{a_m}{{n \choose m-1} }~, \notag \\
  &\leq 2n \sum_{m=2}^M(1-c^m-(1-c)^m)(m-1)a_m\leq 8Mcnd\,, \notag 
\end{align}
\end{subequations}
where the last upper inequality holds when $c \leq c_0 = c_0(m)$ is sufficiently small, since
\begin{align}
    &\, [(1 - c) + c]^m - c^m -(1-c)^m \notag\\
    = &\, \binom{m}{1}(1 - c)^{m-1}c + \binom{m}{2}(1 - c)^{m-2}c^2 + \cdots + \binom{m}{m-1}(1 - c)^{1}c^{m-1}\notag\\
    \leq &\, \binom{m}{1}c + \binom{m}{2}c^2 + \cdots + \binom{m}{m-1}c^{m-1}\leq (1 + c)^{m} - 1 \leq 2(1 + mc - 1) \leq 2mc\,. \label{eqn:czero}
\end{align}
Similarly, we apply Bernstein \Cref{lem:Bernstein} again with $K=M^2$, $\sigma^2\leq 8M^3cnd$ and obtain
\begin{align}
  & \, \mathbb P\left( \sum_{m = 2}^{M} (m-1)^2 |E_m(X,X^c)| \geq 16Mcnd \right) \notag\\
  \leq \,&\, \mathbb P\left( \sum_{m = 2}^{M} (m-1)^2( |E_m(X,X^c)| -\E|E_m(X,X^c)|)\geq 8Mcnd \right) \leq \exp(-6cnd/M)\,. \label{eq:chernoff222}
\end{align}
    % \begin{align}\label{eqn:expectation_EX_Xc}
    %     \sum_{m=2}^{M}m \mu_{|E_m(X, X^{c})|} \= &\, \sum_{m=2}^{M}m \sum_{i=1}^{k} \Bigg[ {\frac{n}{k} \choose m} -  {\eta_{i}cn \choose m} - {(\frac{1}{k} - \eta_{i}c) n \choose m} \Bigg]\frac{a_m - b_m}{{n \choose m-1} } \\
    %     &\, + \sum_{m=2}^{M}m \Bigg[ {n \choose m} - {cn \choose m} - {(1 - c) n \choose m} \Bigg] \frac{b_m}{{n \choose m-1} }\,.\notag 
    %     &\leq \sum_{m=2}^{M}m\Bigg[ {n \choose m} - {cn \choose m} - {(1 - c) n \choose m} \Bigg] \frac{a_m}{{n \choose m-1} }\\
    %   \leq & 2\sum_{m=2}^M n (1-c^m-(1-c)^m) a_m 
    % \end{align}
% According to the results in \Cref{eqn:expectation_EX} and \Cref{eqn:expectation_EX_Xc}, we have
%     \begin{align}
%   \mu:=      \,& \sum_{m=2}^{M} \Big( m \mu_{|E_m(X)|} + (m-1)\mu_{|E_m(X, X^{c})|} \Big) \notag\\
%         =\,& \sum_{m=2}^{M} \Bigg\{ k(m-1) {\frac{n}{k} \choose m} +  \sum_{i=1}^{k} \Big[ {\eta_{i}cn \choose m} - (m-1){( \frac{1}{k} - \eta_{i}c) n \choose m}  \Big] \Bigg\} \frac{a_m - b_m }{ {n \choose m-1} } \\
%         &\, + \sum_{m=2}^{M}\Bigg[ (m-1){n \choose m} + {cn \choose m} - (m-1){(1 - c) n \choose m} \Bigg] \frac{b_m}{{n \choose m-1} } \notag\\
%         \leq &\sum_{m=2}^{M}\Bigg[ (m-1){n \choose m} + {cn \choose m} - (m-1){(1 - c) n \choose m} \Bigg] \frac{a_m}{{n \choose m-1} } \\
%         &\leq 2n \sum_{m=2}^M (1+c^m-(1-c)^m) a_m
%     \end{align}
%  \todo{We need an upper bound and lower bound by optimizing over $\{\eta_{i}\}_{i=1}^{k}$. I believe the upper bound is achieve at corner and lower bound is achieved when $\eta_{i} = \frac{1}{k}$. }  We need this upper bound of $\mu$ to ensure $\nu cnd - \mu > 0$ to apply Bernstein.


% Since $\sum_{m=2}^{M} \Big( m |E_m(X)| + (m-1)|E_m(X, X^{c})| \Big)$ is a weighted sum of independent Bernoulli random variables (see \eqref{eq:EmXBernoulli} and \eqref{eq:EXmcBernoulli}), 

% \begin{align}
%     &\, \sum_{m=2}^{M} \Big( m |E_m(X)| + (m-1)|E_m(X, X^{c})| \Big)\\
%     = &\, \sum_{m=2}^{M} \left\{ \sum_{e\in E_{m}(X, a_m)} mT_e^{(a_m)} \, + \sum_{e\in E_m(X, b_m)} mT_e^{(b_m)} \right\} \\
%     + &\,\sum_{m=2}^{M} \left\{ \sum_{j=1}^{m-1} \left[ \sum_{i=1}^{k} \left( \sum_{e\in \mathcal{E}^{(a_m)}_{j, i}} (m-1)T_{e}^{(a_m)} + \sum_{e\in \mathcal{E}^{(b_m)}_{j, i}} (m-1)T_{e}^{(b_m)} \right) + \sum_{e\in \mathcal{E}^{(b_m)}_{j}} (m-1)T_{e}^{(b_m)} \right] \right\}\notag
% \end{align}

% by Bennett's inequality (\Cref{thm:Bennett}),
% \begin{align}
%     & \mathbb{P}\Bigg[\sum_{m=2}^{M} \Big( m |E(X_{m})| + (m-1)|E(X_{m}, X_{m}^{c})| \Big) -\mu \geq \nu cnd -\mu \Bigg]\\
%     \leq &
% \end{align}
% Here we should ensure $\nu cn d - \mu > 0$ to apply Bernstein
% \newpage
% Let
% \begin{align*}
%     \delta:= \nu cn d \Bigg[\sum_{m=2}^{M} \Big( m \mu_{|E_m(X)|} + (m-1)\mu_{|E_m(X, X^{c})|} \Big) \Bigg]^{-1},
% \end{align*}
% where the second inequality is from \eqref{eq:lowerboundEM}.\todo{Try Bernstein??}
% Then Chernoff inequality  gives
%     \begin{align*}
%         \, & \mathbb{P}\Bigg[\sum_{m=2}^{M} \Big( m |E(X_{m})| + (m-1)|E(X_{m}, X_{m}^{c})| \Big) \geq \nu cnd \Bigg] \\
%         = \, & \mathbb{P}\Bigg[\sum_{m=2}^{M} \Big( m |E(X_{m})| + (m-1)|E(X_{m}, X_{m}^{c})| \geq \delta \sum_{m=2}^{M} \Big( m \mu_{|E(X_{m})|} + (m-1)\mu_{|E(X_{m}, X_{m}^{c})|} \Big) \Bigg] \\
%         \leq  \,& \Bigg( \frac{\exp(\delta - 1)}{\delta^{\delta} } \Bigg)^{\sum_{m=2}^{M} \Big( m \mu_{|E(X_{m})|} + (m-1)\mu_{|E(X_{m}, X_{m}^{c})|} \Big)}\\
%         =\, & \exp \Bigg[\Big( \delta - 1 - \delta \log \delta \Big) \sum_{m=2}^{M} \Big( m \mu_{|E(X_{m})|} + (m-1)\mu_{|E(X_{m}, X_{m}^{c})|} \Big) \Bigg]
%     \end{align*}  
By the binomial coefficient upper bound $\binom{n}{k}\leq (\frac{en}{k})^k$ for $1\leq k\leq n$, there are at most 
    \begin{align}\label{eq:binomialupperbound}
        \binom{n}{cn}\leq \left(\frac{e}{c}\right)^{cn}= \exp(-c(\log c-1)n)
    \end{align}
    % \begin{equation}
    %     {n \choose cn} \leq \exp\left[-\frac{1}{2}\log n - \Big[ c\log c + (1 - c)\log(1- c)  \Big]n - \frac{1}{2}\log c -\frac{1}{2}\log(1 - c)\right]
    % \end{equation}
many subsets $X$ of size $|X| = cn$. Let  $d$ be sufficiently large so that $d^{-3} \leq c_0$. Substituting $c = d^{-3}$ in \eqref{eq:binomialupperbound}, we have
        \begin{equation}
            {n \choose d^{-3}n} \leq \exp \left[ 3d^{-3} \log(d)n \right]. \notag 
        \end{equation}
Taking $c=d^{-3}$ in \eqref{eq:chernoff111} and \eqref{eq:chernoff222}, we obtain
\begin{align}
     & \mathbb P\left( \sum_{m = 2}^{M} (m-1) (m|E_m(X)| + (m-1)|E_m(X,X^c)|)\geq 20Md^{-2}n \right)\leq 2\exp(-2d^{-2}n/M)\,. \notag 
\end{align}
Taking a union bound over all possible $X$ with $|X|=d^{-3}n$, we obtain with probability at least $1-2\exp(3d^{-3}\log d n-2d^{-2}n/M)\leq 1-2\exp(-d^{-2}n/M)$, no more than $d^{-3}n$ many vertices have total row sum greater than $20Md$. Note that we have imposed the condition that $d^{-3} \leq c_0$ (as in \ref{eqn:czero}), so $d$ needs to be sufficiently large. 
\end{proof} 

\subsection{Proof of \Cref{lem:projection_deviation_k}}
\begin{proof}\label{prf:projection_deviation_k}
Note that $\bU$ is spanned by the first $k$ singular vectors of $(\bA_1)_{\cI_1}$. Let $\bu_1, \dots, \bu_k$ be an orthonormal basis of the subspace $\bU$, then the projection $P_{\bU}$ is given by $P_{\bU} := \sum_{l=1}^{k}\< \bu_{l},\cdot\, \> \bu_{l}$. For some fixed vertex $i$ with $k(i)$ indexing its membership, we have for $i\in V_{k(i)} \cap Y_2\cap \{i_1, \cdots, i_s\}$,
\begin{align*}
    P_{\bU}\be_i =&\, \sum_{l=1}^{k}\< \bu_{l}, \be_i\> \bu_{l}, \quad 
    \|P_{\bU}\be_i\|_2^2 =\, \sum_{l=1}^k \< \bu_{l}, \be_i\>^2.
\end{align*}

By definition of $\bA_1$ and $\bA_2$ in \Cref{eqn:A1A2}, and the independence between entries in $\bA_1$ and entries in $\bA_2$, we know that $\{\bu_l\}_{l=1}^{k}$ and $e_i$ are independent of each other, since $\{\bu_l\}_{l=1}^{k}$ is constructed from $\bA_1$ and $e_i$ is chosen from columns of $\bE_2 := \bA_2 - \widetilde{\bA}_2$. If we take expectation over $\{\bcT^{(m)}\}_{m \in \mathcal{M}}$ conditioning on $\{\bu_l\}_{l=1}^{k}$, then by direct calculation,
    \begin{align*}
        \E_{\{\bcT^{(m)}\}_{m \in \mathcal{M}}} \left[ \< \bu_{l}, \be_i\> \Big\vert \{\bu_l\}_{l=1}^{k} \right]=&\, \sum_{j=1}^{n}\bu_{l}(j) \cdot \E \left( \Big[ (\bA_2)_{ji} - (\E\bA_2)_{ji} \Big] \right) = 0\,,\\
        \E_{\{\bcT^{(m)}\}_{m \in \mathcal{M}}} \left[ \|P_{\bU}\be_i\|_2^2 \Big\vert \{\bu_l\}_{l=1}^{k} \right] =&\, \sum_{l=1}^{k} \E_{\{\bcT^{(m)}\}_{m \in \mathcal{M}}} \left[ \< \bu_{l}, \be_i\>^2 \Big\vert \{\bu_l\}_{l=1}^{k} \right]\,,
    \end{align*}
where $\mathcal{M}$ is obtained from \Cref{alg:parameter_preprocessing}. We expand each $\< \bu_{l}, \be_i\>^2$ and rewrite it into $2$ parts,

    \begin{align}
        \< \bu_{l}, \be_i\>^2 =&\, \sum_{j_1 = 1}^{k} \sum_{j_2 = 1}^{k} \bu_{l}(j_1)\be_i(j_1)\bu_{l}(j_2)\be_i(j_2) \notag \\
    =&\, \underbrace{\sum_{j = 1}^{k}[\bu_{l}(j)]^2[\be_i(j)]^2}_{(a)} + \underbrace{\sum_{j_1 \neq j_2} \bu_{l}(j_1)\be_i(j_1)\bu_{l}(j_2)\be_i(j_2)}_{(b)}\,,\quad l\in [k]\,.\label{eqn:contribution_part_a_b}
    \end{align}

% \begin{remark}
%     Part $(a)$ is the contribution from graph, i.e., $2$-uniform hypergraph. Part $(b)$ is the contribution from $m$-uniform hypergraph with $m\geq 3$, which only occurs in hypergraph clustering. 
% \end{remark}
The expectation of part $(a)$ in \Cref{eqn:contribution_part_a_b} is upper bounded by $\alpha$ as defined in \Cref{eqn:alpha_beta}, since
\begin{subequations}
    \begin{align*}
        &\, \E_{\{\bcT^{(m)}\}_{m \in \mathcal{M}}} \left[ \sum_{j = 1}^{k}[\bu_{l}(j)]^2[\be_i(j)]^2 \Bigg\vert \{\bu_l\}_{l=1}^{k} \right]\\
        =&\, \sum_{j=1}^{n}[\bu_{l}(j)]^2 \cdot \Var \Big((\bA_{2})_{ji} \Big) \leq \sum_{j=1}^{n}[\bu_{l}(j)]^2 \cdot (\E\bA_{2})_{ji} \leq \alpha \,, \quad \forall l\in [k]\,,
    \end{align*}
\end{subequations}
where $\|\bu_l\|_2 = 1$ and $\bA_2$ is defined in \Cref{eqn:A1A2}. For part $(b)$,  
\begin{subequations}
    \begin{align*}
        &\, \E_{\{\bcT^{(m)}\}_{m \in \mathcal{M}}} \left[ \sum_{j_1 \neq j_2} \bu_{l}(j_1)\be_i(j_1)\bu_{l}(j_2)\be_i(j_2) \Bigg\vert \{\bu_l\}_{l=1}^{k} \right]\\
        =&\, \sum_{j_1 \neq j_2} \bu_{l}(j_1)\bu_{l}(j_2) \E \left[ \Big( (\bA_2)_{j_{1}i} - (\E\bA_2)_{j_{1}i} \Big) \Big( (\bA_2)_{j_{2}i} - (\E\bA_2)_{j_{2}i} \Big) \right]\\
        =&\, \sum_{j_1 \neq j_2} \bu_{l}(j_1)\bu_{l}(j_2) \E   \Bigg(\sum_{m\in \mathcal{M}}\,\,\sum_{\substack{e\in E_m[Y_2\cup Z]\\ \{i,j_1\}\subset e} }\, (\bcT^{(m)}_{e} - \E\bcT^{(m)}_{e} ) \Bigg) \Bigg(\sum_{m\in \mathcal{M}}\,\,\sum_{\substack{e\in E_m[Y_2\cup Z]\\ \{i,j_2\}\subset e} }\, (\bcT^{(m)}_{e} - \E\bcT^{(m)}_{e} ) \Bigg).
    \end{align*}
\end{subequations}
According to \Cref{def:adjacency_tensor} of the adjacency tensor, $\bcT^{(m)}_{e_1}$ and $\bcT^{(m)}_{e_2}$ are independent if hyperedge $e_1 \neq e_2$, then only the terms with hyperedge $e\supset \{i, j_1, j_2\}$ have nonzero expectation. Then the expectation of part $(b)$ can be rewritten as
    \begin{align}
        &\, \E_{\{\bcT^{(m)}\}_{m \in \mathcal{M}}} \left[ \sum_{j_1 \neq j_2} \bu_{l}(j_1)\be_i(j_1)\bu_{l}(j_2)\be_i(j_2) \Bigg\vert \{\bu_l\}_{l=1}^{k} \right]\notag \\
        = &\, \sum_{j_1 \neq j_2} \bu_{l}(j_1)\bu_{l}(j_2) \sum_{m\in \mathcal{M}}\,\, \sum_{\substack{e\in E_m[Y_2\cup Z]\\ \{i,j_1, j_2\}\subset e} } \E \big(\bcT^{(m)}_{e} - \E\bcT^{(m)}_{e} \big)^2 \notag\\
        \leq &\, \sum_{j_1 \neq j_2} \bu_{l}(j_1)\bu_{l}(j_2) \sum_{m\in \mathcal{M}} \,\,\sum_{\substack{e\in E_m[Y_2\cup Z]\\ \{i,j_1, j_2\}\subset e} } \E\bcT^{(m)}_{e}\notag\\
        = &\sum_{j_1 \neq j_2} \bu_{l}(j_1)\bu_{l}(j_2) \sum_{m\in \mathcal{M}}\sum_{\substack{e\in E_m[Y_2\cup Z]\\ \{i,j_1, j_2\}\subset e} } \frac{a_m}{\binom{n}{m-1}}\,.\label{eqn:expectation_Y2cupZ}
    \end{align}
Note that $|Y_2 \cup Z| \leq n$, then the number of possible hyperedges $e$, while $e\in E_m[Y_2\cup Z]$ and $e \supset \{i,j_1, j_2\}$, is at most $\binom{n}{m-3}$. Thus \Cref{eqn:expectation_Y2cupZ} is upper bounded by
\begin{subequations}
    \begin{align*}
        &\, \sum_{j_1 \neq j_2} \bu_{l}(j_1)\bu_{l}(j_2) \sum_{m\in \mathcal{M}} \binom{n}{m-3} \frac{a_m}{\binom{n}{m-1}} \\
        \leq &\,\sum_{j_1 \neq j_2} \bu_{l}(j_1)\bu_{l}(j_2) \sum_{m\in \mathcal{M}} \frac{(m-1)(m-2)}{(n-m)^2}a_m
        \leq \frac{d \mathcal{M}_{\max} }{n^2} \sum_{j_1 \neq j_2} \bu_{l}(j_1)\bu_{l}(j_2) \,\\
        \leq &\, \frac{d \mathcal{M}_{\max} }{2n^2} \sum_{j_1 \neq j_2}\Big( [\bu_{l}(j_1)]^2 + [\bu_{l}(j_2)]^2 \Big)
        \leq \frac{d \mathcal{M}_{\max}(n-1) }{2n^2}\Bigg( \sum_{j_1 = 1}^{n} [\bu_{l}(j_1)]^2 + \sum_{j_2 = 1}^{n}[\bu_{l}(j_2)]^2 \Bigg) \leq \frac{d \mathcal{M}_{\max} }{n}\,,
    \end{align*}
\end{subequations}
where $\|\bu_l\|_2 = 1$, $d = \sum_{m\in \mathcal{M}}(m-1)a_m$. With the upper bounds for part $(a)$ and $(b)$ in \Cref{eqn:contribution_part_a_b}, the conditional expectation of $\|P_{\bU}\be_i\|_2$ is bounded by
\begin{comment}
We now compute the sum of expectation in \Cref{eqn:expectation_Y2cupZ}. Note that $n_{k(i)}$ and $n^{\prime}_{k(i)}$ denote the cardinality of $Z\cap V_{k(i)}$ and $Y_1 \cap V_{k(i)}$, as defined in \Cref{eqn:dimension_ZcapVi} and \Cref{eqn:dimension_Y1capVi} respectively, while the size of true partition $V_{k(i)}$ is $(n/k)$ for all $k(i)\in [k]$. Thus the cardinality of $Y_2 \cup Z$ and $V_{k(i)} \cap Y_2$ can be written as
\begin{align}
    |Y_2 \cup Z| =&\, n - \sum_{t=1}^{k} n_t^{\prime}\,, \quad |V_{k(i)} \cap Y_2| = \frac{n}{k} - n_{k(i)} - n^{\prime}_{k(i)}\,,\quad k(i) \in [k]\,.
\end{align}
If vertices $i,j_1, j_2$ are all from the true block $V_{k(i)}$, then the $m$-hyperedge $e \supset \{i,j_1, j_2\}$ can be sampled with probability $a_m/\binom{n}{m-1}$ when the remaining $m-3$ vertices are located in $V_{k(i)}$ as well, otherwise with probability $b_m/\binom{n}{m-1}$. In this scenario, its expectation can be written as
\begin{subequations}
    \begin{align*}
        &\,\sum_{\substack{e\in E_m[Y_2\cup Z]\\ \{i,j_1, j_2\}\subset e} } \E\bcT^{(m)}_{e} = \binom{\frac{n}{k} - n^{\prime}_{k(i)} - 3}{m-3} \frac{a_m}{\binom{n}{m-1}} + \left[\binom{n - \sum_{t=1}^{k} n_t^{\prime} - 3}{m-3} - \binom{\frac{n}{k} - n^{\prime}_{k(i)} - 3}{m-3} \right]\frac{b_m}{\binom{n}{m-1}}\\
        \leq &\, \frac{m-2}{\frac{n}{k} - n^{\prime}_{k(i)} - 2}\left\{\binom{\frac{n}{k} - n^{\prime}_{k(i)} - 2}{m-2} \frac{a_m}{\binom{n}{m-1}} + \left[\binom{n - \sum_{t=1}^{k} n_t^{\prime} - 2}{m-2} - \binom{\frac{n}{k} - n^{\prime}_{k(i)} - 2}{m-2} \right]\frac{b_m}{\binom{n}{m-1}} \right\}\\
        = &\, \frac{m-2}{\frac{3n}{4k} - 2}\left\{\binom{ \frac{3n}{4k} - 2}{m-2} \frac{a_m}{\binom{n}{m-1}} + \left[\binom{\frac{3n}{4} - 2}{m-2} - \binom{\frac{3n}{4k} - 2}{m-2} \right]\frac{b_m}{\binom{n}{m-1}} \right\}(1 + o(1))\,,\quad (3\leq m \ll \frac{n}{k})\,,
    \end{align*}
\end{subequations}
where the last equality holds with high probability, since $\E n_{k(i)}^{\prime} = (n/4k)$ and the probability
\begin{align}
    \P\left( \bigg|n^{\prime}_{k(i)} - \frac{n}{4k}\bigg| \geq \sqrt{n}\log(n) \right) \leq 2\exp \left( - k\log^2 (n)\right)\,,\quad k(i)\in [k]\,,
\end{align}
which was shown by Hoeffding \Cref{thm:Hoeffding} in the proof of Lemma \ref{lem:singular_value_approximation}.

On the other hand, if one of the vertices $j_1, j_2$ are not from $V_{k(i)}$, then the hyperedge $e \supset \{i,j_1, j_2\}$ can only be sampled with probability $b_m/\binom{n}{m-1}$, with its expectation
\begin{subequations}
\begin{align}
    \sum_{\substack{e\in E_m[Y_2\cup Z]\\ \{i,j_1, j_2\}\subset e} } \E\bcT^{(m)}_{e} = &\, \binom{n - \sum_{t=1}^{k} n_t^{\prime} - 3}{m-3} \frac{b_m}{\binom{n}{m-1}}\\
    \leq &\, \frac{m-2}{\frac{3n}{4k} - 2}\left\{\binom{ \frac{3n}{4k} - 2}{m-2} \frac{a_m}{\binom{n}{m-1}} + \left[\binom{\frac{3n}{4} - 2}{m-2} - \binom{\frac{3n}{4k} - 2}{m-2} \right]\frac{b_m}{\binom{n}{m-1}} \right\}\,,
\end{align}
\end{subequations}
where the last inequality holds by assumption $a_m \geq b_m$. Returning to \Cref{eqn:expectation_Y2cupZ}, it is upper bounded by
\begin{subequations}
    \begin{align}
        &\, \E_{\{\bcT^{(m)}\}_{m \in \mathcal{M}}} \left[ \sum_{j_1 \neq j_2} \bu_{l}(j_1)\be_i(j_1)\bu_{l}(j_2)\be_i(j_2) \Bigg\vert \{\bu_l\}_{l=1}^{k} \right]\\
        \leq &\, \sum_{j_1 \neq j_2} \bu_{l}(j_1)\bu_{l}(j_2) \sum_{m\in \mathcal{M}} \frac{m-2}{\frac{3n}{4k} - 2}\left\{\binom{ \frac{3n}{4k} - 2}{m-2} \frac{a_m}{\binom{n}{m-1}} + \left[\binom{\frac{3n}{4} - 2}{m-2} - \binom{\frac{3n}{4k} - 2}{m-2} \right]\frac{b_m}{\binom{n}{m-1}} \right\}\\
        \leq &\, \frac{\mathcal{M}_{\max} -2}{\frac{3n}{4k} - 2} \sum_{j_1 \neq j_2} \bu_{l}(j_1)\bu_{l}(j_2) \overline{\alpha} \leq  \frac{\mathcal{M}_{\max} -2}{\frac{3n}{4k} - 2} \cdot  \frac{\overline{\alpha}}{2}\sum_{j_1 \neq j_2}\Big( [\bu_{l}(j_1)]^2 + [\bu_{l}(j_2)]^2 \Big)\\
        \leq &\,\frac{\overline{\alpha}(\mathcal{M}_{\max} -2)}{2} \cdot \frac{n-1}{\frac{3n}{4k} - 2}\Bigg( \sum_{j_1 = 1}^{n} [\bu_{l}(j_1)]^2 + \sum_{j_2 = 1}^{n}[\bu_{l}(j_2)]^2 \Bigg) \leq  2k\overline{\alpha}\mathcal{M}_{\max}\,,
    \end{align}
\end{subequations}
where the last inequality holds since $n\gg k$ and $\|\bu_l\|_2 = 1$. 
\end{comment}
\begin{align*}
     \E_{\{\bcT^{(m)}\}_{m \in \mathcal{M}}} \left[ \|P_{\bU}\be_i\|_2^2 \Big\vert \{\bu_l\}_{l=1}^{k} \right] =&\, \sum_{l=1}^{k} \E_{\{\bcT^{(m)}\}_{m \in \mathcal{M}}} \left[ \< \bu_{l}, \be_i\>^2 \Big\vert \{\bu_l\}_{l=1}^{k} \right] \leq k \Big(\alpha + \frac{d \mathcal{M}_{\max}}{n} \Big)\,,
\end{align*}
By Markov's inequality,
\begin{align*}
    &\,\P\left( \|P_{\bU}\be_i \|_2 > 2 \sqrt{k \Big(\alpha + \frac{d \mathcal{M}_{\max}}{n} \Big) } \right) \leq \frac{ \E_{\{\bcT^{(m)}\}_{m \in \mathcal{M}}} \left[ \|P_{\bU}\be_i\|_2^2 \Big\vert \{\bu_l\}_{l=1}^{k} \right] }{4k(\alpha + d \mathcal{M}_{\max}/n )} \leq \frac{1}{4}\,.
\end{align*}
Let $X_i$ be the Bernoulli random variable defined by
\begin{equation*}
    X_i = \mathbbm{1}\{\|P_{\bU}\be_{i} \|_2 > 2 \sqrt{k (\alpha + d \mathcal{M}_{\max}/n ) } \} \,, \quad i\in \{i_1, \cdots, i_s\}\,.
\end{equation*}
Obviously, $\E X_i = \P(X_i = 1)\leq 1/4$. Define
\begin{align*}
    \delta := \frac{s}{2 \sum_{j=1}^{s}\E X_{i_j}} - 1 \geq 1\,.
\end{align*}
Since $s = 2k\log^2 n$, then by Chernoff's \Cref{lem:Chernoff},
\begin{align*}
    \P \Bigg( \sum_{j=1}^{s}X_{i_j} \geq \frac{s}{2}\Bigg) = \P \Bigg( \sum_{j=1}^{s}X_{i_j} \geq (1 + \delta) \sum_{j=1}^{s}\E X_{i_j}\Bigg) \leq \exp \Bigg( -c\delta^2\sum_{j=1}^{s}\E X_{i_j}\Bigg) =O(e^{-k\log^2 n})\,,
\end{align*}
for some constant $c$ depending on $k$. Therefore, with probability $1 - O(n^{-k\log n})$, at least $s/2$ of the vectors $\be_{i_1}, \dots, \be_{i_s}$ satisfy $\|P_{\bU}\be_i \|_2 \leq 2 \sqrt{k (\alpha + d \mathcal{M}_{\max}/n )}$ for each $i\in \{i_1, i_2,\cdots, i_s\}$.
\end{proof}

\subsection{Proof of \Cref{lem:discard_half_blue}}
\begin{proof}
We start with the following simple claim: for any $m \geq 2$ and any $\nu \in [1/2, 1)$, 
\begin{eqnarray} \label{eq:asympt}
\nu^m + (1-\nu)^m < \left ( \frac{1+\nu}{2} \right)^m~.
\end{eqnarray}
Indeed, one quick way to see this is by induction on $m$; we will induct from $m$ to $m+2$. Assume the inequality is true for $m$; then  
\begin{eqnarray*}
\nu^{m+2} +(1-\nu)^{m+2} & = & \nu^2 \nu^{m} + (1-\nu)^2 (1-\nu)^m \\
& = &  \nu^2 \nu^{m} + (1-2\nu+\nu^2)(1-\nu)^m ~\leq ~\nu^2 \nu^{m}  
+ \nu^2 (1-\nu)^m \\
& = & \nu^2 (\nu^m + (1-\nu)^m ) ~<~ (\nu^2 + (1-\nu)^2) (\nu^m + (1-\nu)^m ) \\
& < & \left (\frac{1+\nu}{2} \right)^2 \left ( \frac{1+\nu}{2} \right)^m = \left ( \frac{1+\nu}{2} \right)^{m+2}~,
\end{eqnarray*}
where we have used the induction hypothesis together with $1-2\nu \leq 0$ and $(1-\nu)^2>0$. After easily checking that the inequality works for $m=2,3$, the induction is complete. 
We shall now check that the quantities defined in   \Cref{lem:discard_half_blue} obey the relationship $\mu_2 \geq \mu_1$ and $\mu_2 - \mu_1 = \Omega(n)$, for $n$ large enough.  First, note that the only thing we need to check is that, for sufficiently large $n$, 
\[
\binom{\frac{\nu n}{2k}}{m} + \binom{\frac{(1-\nu)n}{2k}}{m} \leq \binom{\frac{(1+\nu)n}{4k}}{m} + (k-1) \binom{\frac{(1-\nu)n}{4k(k-1)}}{m}~;
\]
in fact, we will show the stronger statement that for any $m \geq 2$ and $n$ large enough, 
\begin{align} \label{eq:strict_binom}
\binom{\frac{\nu n}{2k}}{m} + \binom{\frac{(1-\nu)n}{2k}}{m} <  \binom{\frac{(1+\nu)n}{4k}}{m} ~,
\end{align}
and this will suffice to see that the second part of the assertion, $\mu_2-\mu_1 = \Omega(n)$, is also true. 
Asymptotically, $\binom{\frac{\nu n}{2k}}{m} \sim \frac{\nu^m}{m!} \left ( \frac{n}{2k} \right)^m $, $\binom{\frac{(1-\nu)n}{2k}}{m} \sim \frac{(1-\nu)^m}{m!} \left ( \frac{n}{2k} \right)^m$, and $\binom{\frac{(1+\nu)n}{4k}}{m} \sim \frac{\left ( \frac{1+\nu}{2} \right)^m}{m!} \left ( \frac{n}{2k} \right)^m$. Note then that \eqref{eq:strict_binom} follows from \eqref{eq:asympt}. 

\vspace{.2cm}

Let \{$V_{1}, \dots, V_{k}$\} be the true partition of $V$. Recall that hyperedges in $H = \cup_{m\in \mathcal{M}}H_m$ are colored red and blue with equal probability in \Cref{alg:k_block_partition}.  Let $E_m(X)$ denote the set of blue $m$-uniform hyperedges with all vertices located in vertex set $X$. 
%and $E(X)$ denote the set of all \textcolor{blue}{blue} hyperedges for varying $m\in \mathcal{M}$, then $E(X) = \cup_{m\in\mathcal{M}}E_{m}(X)$.
Assume $|X \cap V_i| = \eta_i |X|$ with $\sum_{i=1}^{k}\eta_{i} = 1$. For each $m\in\mathcal{M}$, the presence of hyperedge $e\in E_{m}(X)$ can be represented by independent Bernoulli random variables
\begin{equation*}
    T_{e}^{(a_m)}\sim \mathrm{Bernoulli}\left( \frac{a_{m}}{ 2 \binom{n}{m-1} }\right), \quad T_{e}^{(b_m)} \sim \mathrm{Bernoulli}\left( \frac{b_{m}}{ 2\binom{n}{m-1} }\right)\,,
\end{equation*}
depending on whether $e$ is a hyperedge with all vertices in the same block. Denote by \[E_{m}(X, a_m) := \cup_{i=1}^{k}E_{m}(V_i\cap X)\] the union of all $m$-uniform sets of hyperedges with all vertices in the same $V_i\cap X$ for some $i\in [k]$, and by
\begin{align*}
    E_m(X, b_m) :=\, E_{m}(X) \setminus E_{m}(X, a_m) = E_{m}(X) \setminus \Big( \cup_i^{k} E_{m}(V_i\cap X) \Big)
\end{align*}
the set of $m$-uniform hyperedges with vertices across different blocks $V_i\cap X$. Then the cardinality $|E_m(X)|$ can be written as the
\begin{align*}
   |E_m(X)| = \sum_{e\in E_{m}(X, a_m)} T_e^{(a_m)} \, + \sum_{e\in E_m(X, b_m)} T_e^{(b_m)}\,,
\end{align*}
and by summing over $m$, the weighted cardinality $|E(X)|$ is written as
\begin{align*}
    |E(X)| := \sum_{m\in \mathcal{M}}m(m-1)|E_m(X)| = \sum_{m\in \mathcal{M}}m(m-1)\left\{ \sum_{e\in E_{m}(X, a_m)} T_e^{(a_m)} \, + \sum_{e\in E_m(X, b_m)} T_e^{(b_m)} \right\}\,,
\end{align*}
with its expectation
\begin{align} \label{eq:true_expect}
    \E |E(X)| = \sum_{m\in \mathcal{M}}m(m-1)\left\{ \sum_{i=1}^{k} \binom{ \eta_{i} \frac{n}{2k}}{m} \frac{a_m - b_m }{ 2\binom{n}{m-1} } + \binom{\frac{n}{2k}}{m} \frac{b_{m}}{ 2 \binom{n}{m-1} } \right\}\,,
\end{align}
since 
\begin{align*}
    |E_m(X, a_m)| = \sum_{i=1}^{k} |E_{m}(V_i\cap X)|= \sum_{i=1}^{k} \binom{\eta_{i} \frac{n}{2k}}{m}\,, \quad |E_m(X, b_m)|= \binom{ \frac{n}{2k}}{m} - \sum_{i=1}^{k} \binom{\eta_{i} \frac{n}{2k}}{m}\,,
\end{align*}

Next, we prove the two statements in \Cref{lem:discard_half_blue} separately. First, assume that  $|X\cap V_i| \leq \nu |X|$ ( i.e., $\eta_i \leq \nu $) for all $i\in [k]$. Then
        \begin{align*}
            \E |E(X)| \leq &\, \frac{1}{2} \sum_{m\in \mathcal{M}}m(m-1) \left\{ \left[ \binom{\frac{\nu n}{2k}}{m} + \binom{\frac{(1 - \nu)n}{2k}}{m} \right] \frac{a_m - b_m }{ \binom{n}{m-1} } + \binom{\frac{n}{2k}}{m} \frac{b_{m}}{ \binom{n}{m-1} } \right\} =:\mu_1\,.
        \end{align*}
To justify the above inequality, note that since $\sum\limits_{i=1}^k \eta_i = 1$, the sum 
$\sum_{i=1}^k \binom{\eta_i \frac{n}{2k}}{m}
$ is maximized when all but $2$ of the $\eta_i$ are $0$, and since all $\eta_i \leq \nu$, this means that 
\[
\sum_{i=1}^k \binom{\eta_i \frac{n}{2k}}{m} \leq \binom{\frac{\nu n}{2k}}{m} + \binom{\frac{(1-\nu) n}{2k}}{m}\,.
\] 
Note that $m(m-1)(T_{e}^{(a_m)} - \E T_{e}^{(a_m)})$ and $m(m-1)(T_{e}^{(b_m)} - \E T_{e}^{(b_m)})$ are independent mean-zero random variables bounded by $M(M-1)$ for all $m\in \mathcal{M}$, and $\Var(|E(X)|) \leq M^2(M - 1)^2\E |E(X)| = \Omega(n)$. Recall that $\mu_{\mathrm{T}}:= (\mu_1 + \mu_2)/2$. Define $t = \mu_{\mathrm{T}} - \E|E(X)|$, then $ 0 < (\mu_2 - \mu_1)/2 \leq t \leq \mu_{\mathrm{T}}$, hence $t = \Omega(n)$. By Bernstein's Lemma \ref{lem:Bernstein}, we have
\begin{eqnarray*}
       \P \Big( |E(X)| \geq \mu_{\mathrm{T}} \Big) = \P \Big( |E(X)| - \E|E(X)|
       \geq  t \Big) 
       \leq & \exp\left(-\frac{t^2/2}{\Var(|E(X)|) + M(M-1)t/3} \right) = O(e^{-cn})\,,
\end{eqnarray*}
where $c>0$ is some constant.
On the other hand, if $|X\cap V_i| \geq \frac{1+\nu}{2}|X|$ for some $i\in [k]$,  then
    \begin{align*}
       &\,  \E|E(X)|\geq \frac{1}{2} \sum_{m\in \mathcal{M}}m(m-1) \left\{ \left[ \binom{\frac{(1 + \nu)n}{4k} }{m} + (k-1)\binom{ \frac{(1 - \nu)n}{4k(k-1)}}{m} \right] \frac{a_m - b_m }{ \binom{n}{m-1} } + \binom{\frac{n}{2k}}{m} \frac{b_{m}}{ \binom{n}{m-1} } \right\} =:\mu_2\,.
    \end{align*}
    The above can be justified by noting that at least one $|X\cap V_i| \geq \frac{1+\nu}{2} |X|$, and that the rest of the vertices will yield a minimal binomial sum when they are evenly split between the remaining $V_j$. Similarly, define $t = \mu_{\mathrm{T}} - \E|E(X)|$, then $0 < (\mu_2 - \mu_1)/2 \leq -t = \Omega(n)$, and Bernstein's Lemma \ref{lem:Bernstein} gives
    \begin{align*}
       \P \Big( |E(X)| \leq \mu_{\mathrm{T}} \Big) = \, \P \Big( |E(X)| - \E|E(X)| \leq -t \Big) 
       \leq  \exp\left(-\frac{t^2/2}{\Var(|E(X)|) + M(M-1)(-t)/3} \right) = O(e^{-c^{\prime} n})\,,
\end{align*}
where $c^{\prime}>0$ is some other constant.
\end{proof}

\subsection{Proof of \Cref{lem:sample_k_distinct}}

\begin{proof}
    If vertex $i$ is uniformly chosen from $Y_2$, the probability that $i \notin V_{l}$ for some $l \in [k]$ is 
\begin{align*}
    \P(i\notin V_{l} | i\in Y_2) = \frac{\P(i\notin V_{l}, i\in Y_2) }{\P(i \in Y_2)} = 1 - \frac{|V_l \cap Y_2|}{|Y_2|} = 1 - \frac{\frac{n}{k} - n_l - n^{\prime}_l}{n - \sum_{t=1}^{k}(n_t + n_t') }\,,\quad l\in [k]\,,
\end{align*}
where $n_t$ and $n_t^{\prime}$, defined in \Cref{eqn:dimension_ZcapVi} and \Cref{eqn:dimension_Y1capVi}, denote the cardinality of $Z\cap V_t$ and $Y_1\cap V_t$ respectively. As proved in \Cref{subsec:prf_singular_value_approximation}, with probability at least $1 - 2\exp( -k \log^2(n))$, we have $|n_t - n/(2k)|\leq \sqrt{n}\log(n)$ and $|n_t^{\prime} - n/(4k)|\leq \sqrt{n}\log(n)$, then
$
    \P(i\notin V_{l} | i\in Y_2) = 1-\frac{1}{k} \Big(1 + o(1) \Big)\,.
$
After $k\log^2 n$ samples from $Y_2$, the probability that there exists at least one node which belongs to $V_l$ is at least
\[
    1 - \bigg(1 - \frac{1+o(1)}{k} \bigg)^{k\log^2 n}=1-n^{-(1+o(1))k\log(\frac{k}{k-1})\log n}.
\]
The proof is completed by a union bound over $l\in [k]$.

% Note that we could reduce the probability of failure to $n^{-C}$ for any $C>0$, by taking more samples. 
\end{proof}
\subsection{Proof of \Cref{lem:mislabel_probability_correction}}


\begin{proof}
    We calculate $\P(S_{11}^{\prime}(u) \leq \mu_{\rm{C}})$ first. Define $t_{\rm{1C}} := \mu_{\rm{C}} - \E S_{11}^{\prime}(u)$, then by Bernstein's inequality (\Cref{lem:Bernstein}) and taking $K = \mathcal{M}_{\max} - 1$,
\begin{subequations}
    \begin{align*}
    &\, \P \left( S_{11}^{\prime}(u) \leq \mu_{\rm{C}}\right) = \P \left( S^{\prime}_{11}(u) - \E S^{\prime}_{11}(u) \leq t_{\rm{1C}}\right)\\
    \leq &\, \exp\left( - \frac{t_{\rm{1C}}^2/2}{\Var[S_{11}^{\prime}(u)] +(\mathcal{M}_{\max} - 1)\cdot  t_{\rm{1C}}/3} \right) \leq \exp\left( -\frac{3t_{\rm{1C}}^2/(\mathcal{M}_{\max} - 1)}{6(\mathcal{M}_{\max} - 1)\cdot \E S_{11}^{\prime}(u) +  2 t_{\rm{1C}}} \right)\\
    \leq&\, \exp\left( -\frac{[(\nu)^{\mathcal{M}_{\max}-1} - (1 - \nu)^{\mathcal{M}_{\max}-1}]^2 }{(\mathcal{M}_{\max} - 1)^2\cdot 2^{2\mathcal{M}_{\max} + 3}} \cdot \frac{ \left[\sum_{m\in \mathcal{M}} (m-1)\left(\frac{a_m - b_m}{k^{m-1}} \right) \right]^2 }{\sum_{m\in \mathcal{M}} (m-1)\left(\frac{a_m - b_m}{k^{m-1}} + b_m \right)}\right)\,,
\end{align*}
\end{subequations}
where $\mathcal{M}$ is obtained from \Cref{alg:parameter_preprocessing} with $\mathcal{M}_{\max}$ denoting the maximum value in $\mathcal{M}$, and the last two inequalities hold since $\Var[S_{11}^{\prime}(u)] \leq (\mathcal{M}_{\max} - 1)^2\E S_{11}^{\prime}(u)$, and for sufficiently large $n$,
\begin{subequations}
    \begin{align*}
     t_{\rm{1C}} :=&\, \mu_{\rm{C}} - \E S_{11}^{\prime}(u) = -\frac{1}{2} \sum_{m\in \mathcal{M}} (m-1)\cdot\left[ \binom{\frac{\nu n}{2k}}{m-1} - \binom{\frac{(1 - \nu)n}{2k}}{m-1} \right]\frac{a_m - b_m}{2\binom{n}{m-1}}\\
     \leq &\, - \frac{1}{2}\sum_{m\in \mathcal{M}} \frac{(\nu)^{m-1} - (1 - \nu )^{m-1}}{2^{m}}\cdot(m-1) \frac{a_m - b_m}{k^{m-1}}(1 + o(1))\\
    \leq &\, -\frac{(\nu)^{\mathcal{M}_{\max} - 1} - (1 - \nu)^{\mathcal{M}_{\max} - 1} }{2^{\mathcal{M}_{\max} + 2}} \sum_{m\in \mathcal{M}}(m-1) \cdot\frac{a_m - b_m}{k^{m-1}} \,,
    \end{align*}
\end{subequations}

%\todo{change constant C(M)}.

\begin{subequations}
    \begin{align*}
    &\, 6(\mathcal{M}_{\max} - 1)\E S_{11}^{\prime}(u) +  2t_{\rm{1C}} = 2\mu_{\rm{C}} + (6\mathcal{M}_{\max} - 8)\E S^{\prime}_{11}(u)\\
         = &\, \sum_{m\in \mathcal{M}} (m-1)\left\{  \left[ (6\mathcal{M}_{\max} - 7)\binom{\frac{\nu n}{2k}}{m-1} + \binom{\frac{(1 - \nu) n}{2k}}{m-1} \right] \frac{a_m - b_m}{2\binom{n}{m-1}} +  6(\mathcal{M}_{\max} - 1)\binom{\frac{n}{2k}}{m-1}\frac{b_m}{2\binom{n}{m-1}}\right\}\\
         = &\, \sum_{m\in \mathcal{M}} (m-1)\left[ \frac{(6\mathcal{M}_{\max} - 7)\cdot (\nu)^{m-1} + (1 - \nu )^{m-1} }{2^{m}}\cdot \frac{a_m - b_m}{k^{m-1}} + \frac{(6\mathcal{M}_{\max} - 6)b_m}{2^{m}k^{m-1}}
         \right](1+o(1))\\
         \leq &\, \sum_{m\in \mathcal{M}} \frac{(6\mathcal{M}_{\max} - 7)\cdot(\nu)^{m-1} + (1 - \nu )^{m-1} }{2^{m}}\cdot (m-1)\left(\frac{a_m - b_m}{k^{m-1}} + b_m\right)(1+o(1))\\
         \leq &\, \frac{3(\mathcal{M}_{\max} - 1)}{2}\sum_{m\in \mathcal{M}} (m-1)\left(\frac{a_m - b_m}{k^{m-1}} + b_m\right)\,.
\end{align*}
\end{subequations}
%for sufficiently large $n$.
Similarly, for $\P(S_{1j}^{\prime}(u) \geq \mu_{\rm{C}})$, define $t_{\rm{jC}} := \mu_{\rm{C}} - \E S_{1j}^{\prime}(u)$ for $j\neq 1$, by Bernstein's Lemma \ref{lem:Bernstein}, 
\begin{subequations}
\begin{align*}
    &\, \P \left( S_{1j}^{\prime}(u) \geq \mu_{\rm{C}}\right) = \P \left( S^{\prime}_{1j}(u) - \E S^{\prime}_{1j}(u) \geq t_{\rm{jC}}\right)\\
    \leq &\, \exp\left( - \frac{t_{\rm{jC}}^2/2}{\Var[S_{1j}^{\prime}(u)] + (\mathcal{M}_{\max} - 1)\cdot t_{\rm{jC}}/3} \right) \leq \exp\left( -\frac{3t_{\rm{jC}}^2/(\mathcal{M}_{\max} - 1)}{6(\mathcal{M}_{\max} - 1)\cdot\E S_{1j}^{\prime}(u) +  2 t_{\rm{jC}}} \right)\\
    \leq&\, \exp\left( -\frac{[(\nu)^{\mathcal{M}_{\max}-1} - (1 - \nu )^{\mathcal{M}_{\max}-1}]^2 }{(\mathcal{M}_{\max} - 1)^2\cdot 2^{2\mathcal{M}_{\max} + 3}} \cdot \frac{ \left[\sum_{m\in \mathcal{M}} (m-1)\left(\frac{a_m - b_m}{k^{m-1}} \right) \right]^2 }{\sum_{m\in \mathcal{M}} (m-1)\left(\frac{a_m - b_m}{k^{m-1}} + b_m \right)} \right)\,.
\end{align*}
\end{subequations}
The last two inequalities holds since $\Var[S_{1j}^{\prime}(u)] \leq (\mathcal{M}_{\max} - 1)^2\E S_{1j}^{\prime}(u)$, and for sufficiently large $n$, 
\begin{subequations}
    \begin{align*}
     t_{\rm{jC}} :=&\, \mu_{\rm{C}} - \E S_{1j}^{\prime}(u) = \frac{1}{2} \sum_{m\in \mathcal{M}} (m-1)\cdot \left[ \binom{\frac{\nu n}{2k}}{m-1} - \binom{\frac{(1 - \nu)n}{2k}}{m-1} \right]\frac{a_m - b_m}{2\binom{n}{m-1}}\\
    \geq &\, \frac{(\nu)^{\mathcal{M}_{\max} - 1} - (1 - \nu )^{\mathcal{M}_{\max} - 1} }{2^{\mathcal{M}_{\max} + 2}} \sum_{m\in \mathcal{M}}(m-1) \cdot\frac{a_m - b_m}{k^{m-1}} \,\,,
    \end{align*}
\end{subequations}

\begin{subequations}
    \begin{align*}
    &\,6\E S_{1j}^{\prime}(u) +  2t_{\rm{jC}} = 2\mu_{\rm{C}} + (6\mathcal{M}_{\max} - 8)\E S^{\prime}_{1j}(u)\\
         = &\, \sum_{m\in \mathcal{M}} (m-1)\left\{ \left[ \binom{\frac{\nu n}{2k}}{m-1} + (6\mathcal{M}_{\max} - 7)\binom{\frac{(1 - \nu)n}{2k}}{m-1} \right] \frac{a_m - b_m}{2\binom{n}{m-1}} +  6(\mathcal{M}_{\max} -1)\binom{\frac{n}{2k}}{m-1}\frac{b_m}{2\binom{n}{m-1}}\right\}\\
        = &\, \sum_{m\in \mathcal{M}} (m-1)\cdot\left( \frac{(\nu)^{m-1} + (6\mathcal{M}_{\max} - 7)\cdot (1 - \nu )^{m-1}}{2^{m}}\cdot \frac{a_m - b_m}{k^{m-1}} + \frac{(6\mathcal{M}_{\max} - 6)\cdot b_m}{2^{m}k^{m-1}}\right)(1+o(1))\\
         \leq &\, \sum_{m\in \mathcal{M}} \frac{ (\nu)^{m-1} + (6\mathcal{M}_{\max} - 7)\cdot(1 - \nu )^{m-1}}{2^{m}}\cdot (m-1)\left(\frac{a_m - b_m}{k^{m-1}} + b_m\right)(1+o(1))\\
         \leq &\, \frac{3(\mathcal{M}_{\max} - 1)}{2}\sum_{m\in \mathcal{M}} (m-1)\left(\frac{a_m - b_m}{k^{m-1}} + b_m\right).
\end{align*}
\end{subequations}
\end{proof}
\section{Algorithm’s correctness for the binary case}\label{sec:analysis_2}

We will show the correctness of \Cref{alg:2_block_partition} and prove \Cref{thm:mainresult_2} in this section. The analysis will mainly follow from the analysis in \Cref{sec:analysis_k}. We only detail the differences. 
% First, \Cref{eqn:expected_adjacency_k} becomes 
% \begin{align}
%   \overline{\bA}: = \E \bA = \begin{bmatrix}
%      \alpha \bJ_{\frac{n}{2}}  &  \beta \bJ_{\frac{n}{2}}\\
%     \beta  \bJ_{\frac{n}{2}} & \alpha \bJ_{\frac{n}{2}}
%   \end{bmatrix} - \alpha \bI_{n}\,,\notag
% \end{align}
% with $\alpha$ and $\beta$ defined in \Cref{eqn:alpha_beta}. 


Without loss of generality, we assume $n$ is even to guarantee the existence of a binary partition of size $n/2$. The method to deal with the odd $n$ case was discussed in \Cref{lem:general_eigen_value_approximation}. 
Then, let the index set be \(\cI = \{i\in [n]: \mathrm{row}(i) \leq 20\mathcal{M}_{\max} d\}\), as shown in \Cref{eqn:restricted_adjacency_matrix}.  Let $\bu_i$ (resp. $\bar{\bu}_i$) denote the eigenvector associated to $\lambda_i(\bA_{\cI})$ (resp. $\lambda_i(\overline{\bA})$) for $i =1, 2$. Define two linear subspaces $\bU:= \mathrm{Span}\{\bu_{1}, \bu_{2}\}$ and $\overline{\bU}:= \mathrm{Span}\{\bar{\bu}_{1}, \bar{\bu}_{2}\}$, then the angle between $\bU$ and $\overline{\bU}$ is defined as
 $
      \sin \angle(\bU, \overline{\bU}) := \|P_{\bU} - P_{\overline{\bU}}\|
  $,
where $P_{\bU}$ and $P_{\overline{\bU}}$ are the orthogonal projections onto $\bU$ and $\overline{\bU}$, respectively.

\subsection{Proof of \Cref{lem:spectralclustering_accuracy_2}}

\subsubsection{Bound the angle between $\bU$ and $\overline{\bU}$.}\label{subsubsec:bound_subspace_angle_2} The strategy to bound the angle is similar to \Cref{subsubsec:bound_subspace_angle_k}, except that we apply Davis-Kahan Theorem (\Cref{lem:Davis_Kahan_sin}) here. 
% We should pay attention to two things: (i) the spectral gap $|\lambda_2(\overline{\bA})| - |\lambda_{3}(\overline{\bA})|$ and (ii) the deviation $\|\bA_{\cI} - \overline{\bA}\|$.
% While the spectrum of $\overline{\bA}$ was discussed in \Cref{lem:eigenvalues_HSBM}, we focus on the deviation $\|\bA_{\cI} - \overline{\bA}\|$. 

Define $\bE := \bA - \overline{\bA}$ and its restriction on $\cI$, namely $\bE_{\cI} := (\bA - \overline{\bA})_{\cI} = \bA_{\cI} - \overline{\bA}_{\cI}$, as well as $\bDelta:= \overline{\bA}_{\cI} - \overline{\bA}$. Then the deviation $\bA_{\cI} - \overline{\bA}$ is decomposed as
\begin{align}
    \bA_{\cI} - \overline{\bA} = (\bA_{\cI} - \overline{\bA}_{\cI}) + (\overline{\bA}_{\cI} - \overline{\bA}) = \bE_{\cI} + \bDelta\,.\notag
\end{align}

\Cref{thm:regularization} indicates $\|\bE_{\cI}\|\leq C_3\sqrt{d}$ with probability at least $1 - n^{-2}$ when taking $\tau=20\mathcal{M}_{\max}, K=3$, where $C_3$ is a constant depending only on $\mathcal{M}_{\max}$. Moreover, \Cref{lem:high_degree_vertices_k} shows that the number of vertices with high degrees is relatively small. Consequently, an argument similar to \Cref{cor:norm_of_high_degree_vertices_k_block} leads to the conclusion $\|\bDelta\| \leq \sqrt{d}$ \textit{w.h.p.} Together with upper bounds for $\|\bE_{\cI}\|$ and $\|\bDelta\|$, \Cref{lem:subspace_angle_2} shows that the angle between $\bU$ and $\overline{\bU}$ is relatively small \textit{w.h.p.}

\begin{lemma}\label{lem:subspace_angle_2}
    For any $c\in (0,1)$, there exists a constant $C_2$ depending on $\mathcal M_{\max}$ and $c$ such that if
  \begin{align}
     \sum_{m\in\mathcal M}(m-1)(a_m-b_m)\geq C_2 \cdot 2^{\mathcal M_{\max}+2}\sqrt{d},\notag
  \end{align}
   then $\sin \angle (\bU, \overline{\bU}) \leq c$ with probability $1 - n^{-2}$. 
\end{lemma}
\begin{proof}
    First, with probability $1 - n^{-2}$, we have
\begin{align}
     \|\bA_{\cI} - \overline{\bA}\| \leq  \|\bE_{\cI}\| + \|\bDelta \|\leq (C_3+1)\sqrt{d}. \notag 
\end{align}
    According to the definitions in \Cref{eqn:alpha_beta}, $\alpha \geq \beta$ and $\alpha  = O(1/n)$, $\beta = O(1/n)$. Meanwhile, \Cref{lem:eigenvalues_HSBM} shows that $|\lambda_2(\overline{\bA})| = [- \alpha + (\alpha - \beta)n/2]$ and $|\lambda_{3}(\overline{\bA})| = \alpha$. Then
    \begin{subequations}
        \begin{align*}
        &\,|\lambda_2(\overline{\bA})| - |\lambda_{3}(\overline{\bA})| = \frac{n}{2}(\alpha - \beta) - 2\alpha \geq \frac{3}{4}\cdot \frac{n}{2}(\alpha - \beta) = \frac{3n}{8} \sum_{m\in \mathcal{M}} \binom{\frac{n}{2} - 2}{m-2}  \frac{(a_m - b_m)}{\binom{n}{m-1}}\\ \geq &\, \frac{1}{4}\sum_{m\in \mathcal{M}} \frac{(m-1)(a_m - b_m)}{2^{m-2}} \geq \frac{1}{2^{\mathcal{M}_{\max}}}\sum_{m\in \mathcal{M}}(m-1)(a_m - b_m) \geq 4C_2\sqrt{d}\,.
        \end{align*}
    \end{subequations}
    Then for some large enough $C_2$, the following condition for Davis-Kahan Theorem (\Cref{lem:Davis_Kahan_sin}) is satisfied
\begin{align}
      \|\bA_{\cI} - \overline{\bA}\|\leq (1-1/\sqrt{2}) \left(|\lambda_2(\overline{\bA})| - |\lambda_{3}(\overline{\bA})|\right). \notag 
\end{align}
Then for any $c\in(0, 1)$, we can choose $C_2 = (C_3 + 1)/c$ such that 
    \begin{align}
        \|P_{\bU} - P_{\overline{\bU}}\| \leq \frac{2\|\bA_{\cI} - \overline{\bA}\|}{ |\lambda_2(\overline{\bA})| - |\lambda_{3}(\overline{\bA})|} \leq \frac{2(C_3+1)\sqrt{d}}{4C_2\sqrt{d}} = \frac{c}{2} \leq c\,. \notag 
    \end{align}
\end{proof}

Now, we focus on the accuracy of \Cref{alg:2_block_spectral_partition}, once the conditions in \Cref{lem:subspace_angle_2} are satisfied.

\begin{lemma}[Lemma 23 in \cite{chin2015stochastic}]\label{lem:vector_angle_2block}
If $\sin \angle(\overline{\bU}, \bU) \leq c \leq \frac{1}{4}$,  there exists a unit vector $\bv\in \bU$ such that $\sin \angle(\bar{\bu}_{2}, \bv) \leq 2 \sqrt{c}$.
\end{lemma}
The desired vector $\bv$, as constructed in \Cref{alg:2_block_spectral_partition}, is the unit vector perpendicular to $P_{\bU}\mathbf{1}_n$, where $P_{\bU}\mathbf{1}_n$ is the projection of all-ones vector onto $\bU$. \Cref{lem:subspace_angle_2} and \Cref{lem:vector_angle_2block} together give the following corollary.
\begin{corollary}
    For any  $c \in (0, 1)$, there exists a unit vector $\bv\in \bU$ such that $\sin \angle( \bar{\bu}_{2}, \bv) \leq c < 1$ with probability  $1 - O(e^{-n})$.
\end{corollary}
%\begin{proof}
%    For any $c \in (0, 1)$, we could choose constants $C_2$, $C_3$ in \Cref{lem:subspace_angle_2} such that $\sin \angle(\overline{\bU}, \bU) \leq \frac{c^{2}}{4} < 1$. Then by \Cref{lem:vector_angle_2block}, we construct $\bv$ such that $\sin \angle(\bar{\bu}_{2}, \bv) \leq c$.
%\end{proof}

\begin{lemma}[Lemma $23$ in \cite{chin2015stochastic}]
If $\sin \angle(\bar{\bu}_{2}, \bv) < c \leq 0.5$, then we can identify at least $(1 - 4c^{2}/3)n$ vertices from each block correctly.
\end{lemma}

%\begin{proof}
%    Define two sets $V_{1}^{\prime} = \{i|\bv(i) >0\}$ and $V_{2}^{\prime} = \{i|\bv(i) \leq 0\}$. Without loss of generality, we assume $|V_{1}^{\prime}| \leq \frac{n}{2}$. We write
   %\begin{align}
    %    \bv = c_{1}\bar{u}_{2} + \bepsilon = \Big[\bepsilon(1) + \frac{c_{1}}{\sqrt{n}}, \cdots, \bepsilon(n/2) + \frac{c_{1}}{\sqrt{n}},\,\, \bepsilon(n/2 + 1) - \frac{c_{1}}{\sqrt{n}} , \cdots,  \bepsilon(n) -\frac{c_{1}}{\sqrt{n}} \Big]^{\sT}\,,
    %\end{align}
    %where $\bar{\bu}_{2} \perp \bepsilon$ with $\|\bepsilon\|_2 < c$ and $c_{1} \geq \sqrt{1 - c^{2}}$. Thus there are at most $\frac{c^{2}}{1 - c^{2}}n$ coordinates of $\bepsilon$ can be bigger than $\frac{\sqrt{1 - c^{2}}}{\sqrt{n}}$, otherwise $\|\bepsilon\|_2 \geq c$. Notice that $c_{1} \geq \sqrt{1 - c^{2}}$, then at least $1 - \frac{c^{2}}{1 - c^{2}} > 1 - \frac{4}{3}c^{2}$ fraction of vertices with $\bar{\bu}_{2}(i) = \frac{1}{\sqrt{n}}$ will have $\bv(i) > 0$. As a result, $V_{1}^{\prime}$ contains at least $(1 - \frac{4}{3}c^{2})n$ vertices from the true block $V_1$.
%\end{proof}

The proof of \Cref{lem:spectralclustering_accuracy_2} is completed when choosing $C_2$, $C_3$ in \Cref{lem:subspace_angle_2} such that $c \leq 1/4$.


\subsection{Proof of \Cref{lem:correction_accuracy_2}}
The proof strategy is similar to \Cref{subsec:local_correction} and \Cref{subsec:merging}. In \Cref{alg:2_block_partition}, we first color the hyperedges with red and blue with equal probability. By running \Cref{alg:2_block_correction} on the red graph, we obtain a $\nu$-correct partition $V_1^{\prime}, V_2^{\prime}$ of $V = V_1 \cup V_2$, i.e., $|V_{l}\cap V_{l}^{\prime}| \geq \nu n/2$ for $l = 1, 2$. In the rest of the proof, we condition on this event and  the event that the maximum red degree of a vertex is at most $\log^{2}(n)$ with probability at least $1 - o(1)$. This can be proved by Bernstein's inequality (\Cref{lem:Bernstein}).

Similarly, we consider the probability of a hyperedge $e = \{i_{1}, \cdots, i_{m}\}$ being blue conditioning on the event that $e$ is not a red hyperedge, in each underlying $m$-uniform hypergraph separately. If vertices $i_{1}, \cdots, i_{m}$ are all from the same true cluster, then the probability is $\psi_m$, otherwise $\phi_m$, where $\psi_m$ and $\phi_m$ are defined in \Cref{eqn:psi_m} and \Cref{eqn:phi_m}, and the presence of those hyperedges are represented by random variables $\zeta_e^{(a_m)} \sim \mathrm{Bernoulli}\left(\psi_m\right)$, $\xi_e^{(b_m)} \sim \mathrm{Bernoulli}\left(\phi_m\right)$, respectively.

%For any vertex $u\in V_{l}$ with $l \in \{1, 2\}$, we want to compute the number of hyperedges, which contains $u$ with the remaining vertices located in vertex set $V^{\prime}_j$ with $j\in \{1, 2\}$.
Following a similar argument in \Cref{subsec:local_correction}, the row sum of $u$ can be written as
\begin{align}
    S^{\prime}_{lj}(u):= \sum_{m\in \mathcal{M}}(m-1)\cdot \left\{ \sum_{e\in \, \mathcal{E}^{(a_m)}_{l, j}} \zeta_e^{(a_m)} + \sum_{e\in \mathcal{E}^{(b_m)}_{l, j}} \xi_e^{(b_m)}\right\}\,,\quad u\in V_l\,, \notag 
\end{align}
 where $\mathcal{E}^{(a_m)}_{l, j}:= E_m([V_l]^{1}, [V_l\cap V^{\prime}_j]^{m-1})$ denotes the set of $m$-hyperedges with $1$ vertex from $[V_l]^{1}$ and the other $m-1$ vertices from $[V_l\cap V^{\prime}_j]^{m-1}$, while
 \begin{align}
        \mathcal{E}^{(b_m)}_{l, j} := E_m \Big([V_l]^{1}, \,\, [V^{\prime}_j]^{m-1} \setminus [V_l\cap V^{\prime}_j]^{m-1} \Big) \notag 
    \end{align}
    denotes the set of $m$-hyperedges with $1$ vertex in $[V_l]^{1}$ while the remaining $m-1$ vertices in $[V^{\prime}_j]^{m-1}\setminus[V_l\cap V^{\prime}_j]^{m-1}$, with their cardinalities
    \begin{align}
       |\mathcal{E}^{(a_m)}_{l, j}| \leq \binom{|V_l\cap V^{\prime}_j|}{m-1}\,,\quad |\mathcal{E}^{(b_m)}_{l, j}| \leq \left[ \binom{|V^{\prime}_j|}{m-1} - \binom{|V_l\cap V^{\prime}_j|}{m-1} \right]\,. \notag 
    \end{align}
%We multiply $(m-1)$ inside $S^{\prime}_{lj}(u)$ since we are counting row sum of $u$. 
According to the fact $|V_{l}\cap V_{l}^{\prime}| \geq \nu n/2$, $|V_l| = n/2$, $|V^{\prime}_l| = n/2$ for $l = 1, 2$, we have
\begin{align}
    |\mathcal{E}^{(a_m)}_{l, l}| \geq \binom{ \frac{\nu n}{2}}{m-1}\,,\quad |\mathcal{E}^{(a_m)}_{l, j}| \leq \binom{\frac{(1 - \nu) n}{2}}{m-1}\,,\,\, j\neq l\,. \notag 
\end{align}
To simplify the calculation, we take the lower and upper bound of $|\mathcal{E}^{(a_m)}_{l, l}|$ and $|\mathcal{E}^{(a_m)}_{l, j}|(j\neq l)$ respectively. Taking expectation with respect to $\zeta_e^{(a_m)}$ and $\xi_e^{(b_m)}$, for any $u \in V_l$, we have
\begin{subequations}
    \begin{align}
    \E S^{\prime}_{ll}(u) &= \sum_{m\in \mathcal{M}} (m-1)\cdot\left[ \binom{\frac{\nu n}{2}}{m-1} (\psi_{m} - \phi_{m}) + \binom{\frac{n}{2}}{m-1} \phi_{m}\right], \notag \\
    \E S^{\prime}_{lj}(u) &= \sum_{m\in \mathcal{M}} (m-1)\cdot\left[ \binom{\frac{(1 - \nu) n}{2}}{m-1} (\psi_{m} - \phi_{m}) +  \binom{\frac{n}{2}}{m-1} \phi_{m} \right],\,\, j\neq l\,. \notag 
\end{align}
\end{subequations}
By assumptions in \Cref{thm:mainresult_k}, $\E S^{\prime}_{ll}(u) - \E S^{\prime}_{lj}(u) = \Omega(1)$. We define
\begin{align}
    \mu_{\rm{C}} := \frac{1}{2}\sum_{m\in \mathcal{M}} (m-1)\cdot\left\{ \left[ \binom{\frac{\nu n}{2}}{m-1} + \binom{ \frac{(1 - \nu)n}{2} }{m-1} \right](\psi_{m} - \phi_{m}) + 2\binom{ \frac{n}{2} }{m-1} \phi_{m}\right\}\,. \notag 
\end{align}
After \Cref{alg:k_block_merging}, if a vertex $u\in V_l$ is mislabelled, one of the following event must happen
\begin{itemize}
    \item $S^{\prime}_{ll}(u) \leq \mu_{\rm{C}}$,
    \item $S^{\prime}_{lj}(u) \geq \mu_{\rm{C}}$, for some $j\neq l$.
\end{itemize}
By an argument similar to \Cref{lem:mislabel_probability_correction}, we can prove that
\begin{align}
    \rho_{1}^{
    \prime} = \P \left( S^{\prime}_{ll}(u) \leq \mu_{\rm{C}}\right) \leq \rho\,, \quad \rho_{2}^{
    \prime} = \P \left( S^{\prime}_{lj}(u) \geq \mu_{\rm{C}}\right) \leq \rho\,, \notag 
\end{align}
where $\rho = \exp\left( -C_{\mathcal M}\cdot \mathrm{SNR}_{\mathcal{M}} \right)$ and  
\begin{align}
    C_{\mathcal M}:=\frac{[(\nu)^{\mathcal{M}_{\max}-1} - (1 - \nu)^{\mathcal{M}_{\max}-1}]^2 }{8(\mathcal{M}_{\max} - 1)^2}, \quad \mathrm{SNR}_{\mathcal{M}}:= \frac{ \left[\sum_{m\in \mathcal{M}} (m-1)\left(\frac{a_m - b_m}{2^{m-1}} \right) \right]^2 }{\sum_{m\in \mathcal{M}} (m-1)\left(\frac{a_m - b_m}{2^{m-1}} + b_m \right)}\,. \notag 
\end{align}
As a result, the probability that either of those events happened is bounded by $\rho$. The number of mislabeled vertices in $V_1$ after \Cref{alg:k_block_correction} is at most
\begin{align*}
    R_l = \sum_{i=1}^{|V_l\setminus V_l^{\prime}|}\Gamma_{i}\, + \sum_{i=1}^{|V_l \cap V_l^{\prime} |}\Lambda_{i}\,,
\end{align*}
where $\Gamma_{i}$ (resp. $\Lambda_{i}$) are i.i.d indicator random variables with mean $\rho_1^{\prime}$ (resp. $\rho_2^{\prime}$). Then 
\begin{align*}
   \E R_l \leq  \frac{n}{2} \rho_1^{\prime} + \frac{(1 - \nu)n}{2} \rho_2^{\prime} =(1 - \nu/2)n\rho\,. 
\end{align*}
where $\nu$ is the correctness after \Cref{alg:k_block_spectral_partition}. Let $t_l:= (1 + \nu/2)n\rho$, then by Chernoff \Cref{lem:Chernoff},
\begin{align*}
   \P \left( R_l \geq n\rho \right) = \P \left[ R_l - (1 - \nu/2)n\rho \geq t_l \right] \leq \P \left( R_l - \E R_l \geq t_l \right) \leq e^{-c t_l} = O(e^{-n\rho})\,,
\end{align*}
which means that with probability $1- O(e^{- n\rho})$, the fraction of mislabeled vertices in $V_l$ is smaller than $2\rho$, i.e., the correctness of $V_l$ is at least $\gamma := \max\{\nu, 1 - 2\rho\}$. 

\section{Useful Lemmas}

% \begin{lemma}\label{lem:binomial_convex}
% The function $f: \N^{+} \mapsto \N^{+}$, defined by $f(n) = \binom{n}{m}$, is convex for any $m\geq 2$. In particular, if $n_1 + n_2 + \cdots + n_k = n$ for some $k\in \N^{+}$, then $f(n_1) + f(n_2) + \cdots + f(n_k) \leq f(n)$.
% \end{lemma}
% \begin{proof}
% By induction, it suffices to check when $k=2$, which can be verified by direction calculation $\binom{n_1}{m}+\binom{n_2}{m}\leq \binom{n_1+n_2}{m}$.
% \end{proof}

\begin{lemma}[Chernoff's inequality, Theorem $2.3.6$ in \cite{vershynin_2018}]\label{lem:Chernoff}
    Let $X_i$ be independent Bernoulli random variables with parameters $p_i$. Consider their sum $S_N = \sum_{i=1}^{N}X_i$ and denote its mean by $\mu = \E S_N$. Then for any $\delta\in (0, 1]$, 
    \begin{align*}
        \P \big( |S_N - \mu| \geq \delta \mu \big) \leq 2 \exp( -c\delta^2 \mu )\,.
    \end{align*}
\end{lemma}

\begin{lemma}[Hoeffding's inequality, Theorem $2.2.6$ in \cite{vershynin_2018}]\label{lem:Hoeffding}
    Let $X_1,\dots, X_N$ be independent  random variables with $X_i \in [a_i, b_i]$ for each $i\in \{1, \dots, N\}$. Then for any, $t \geq 0$, we have
    \begin{align*}
        \P\Bigg( \bigg|\sum_{i=1}^{N}(X_i - \E X_i ) \bigg| \geq t \Bigg) \leq 2\exp \Bigg( -\frac{2t^{2} }{ \sum_{i=1}^{N}(b_i - a_i )^{2} } \Bigg)\,.
    \end{align*}
\end{lemma}

\begin{lemma}[Bernstein's inequality, Theorem $2.8.4$ in \cite{vershynin_2018}]\label{lem:Bernstein}
    Let $X_1,\dots,X_N$ be independent mean-zero random variables such that $|X_i|\leq K$ for all $i$. Let $\sigma^2 = \sum_{i=1}^{N}\E X_i^2$. 
    Then for every $t \geq 0$, we have
    \begin{align*}
        \P \Bigg( \Big|\sum_{i=1}^{N} X_i \Big| \geq t \Bigg) \leq 2 \exp \Bigg( - \frac{t^2/2}{\sigma^2 + Kt/3} \Bigg)\,.
    \end{align*}
\end{lemma}

\begin{lemma}[Bennett's inequality, Theorem $2.9.2$ in \cite{vershynin_2018}]\label{lem:Bennett}
    Let $X_1,\dots, X_N$ be independent random variables. Assume that $|X_i - \E X_i| \leq K$ almost surely for every $i$. Then for any $t>0$, we have
    \begin{align}
        \P \Bigg( \sum_{i=1}^{N} (X_i - \E X_i) \geq t \Bigg) \leq \exp \Bigg( - \frac{\sigma^2}{K^2} \cdot h \bigg( \frac{Kt}{\sigma^2} \bigg)\Bigg)\,, \notag 
    \end{align}
    where $\sigma^2 = \sum_{i=1}^{N}\textnormal{Var}(X_i)$, and $h(u) := (1 + u)\log(1 + u) - u$.
\end{lemma}

%\todo{this one follows directly from singular value perturbation inequality, you can remove it}
%\begin{lemma}[Weyl's eigenvalue perturbation inequality]\label{lem:Weyl_eigenvalue_perturbation}
   % Let $\bA, \bE \in \R^{n \times n}$ be two real symmetric matrices, then $|\lambda_i(\bA + \bE) - \lambda_i(\bA)| \leq \|\bE\|$ for all $1 \leq i \leq n$.
%\end{lemma}


%\todo{or if you keep thm D.5, then D.6 follows directly by writing $A,E$ into a $(m+n)\times (m+n)$ matrix, see wikipedia}
\begin{lemma}[Weyl's inequality]\label{lem:Weyl_inequality}
Let $\bA, \bE \in \R^{m \times n}$ be two real $m\times n$ matrices, then $|\sigma_i(\bA + \bE) - \sigma_i(\bA)| \leq \|\bE\|$ for every $1 \leq i \leq \min\{ m, n\}$. Furthermore, if $m = n$ and $\bA, \bE \in \R^{n \times n}$ are real symmetric, then $|\lambda_i(\bA + \bE) - \lambda_i(\bA)| \leq \|\bE\|$ for all $1 \leq i \leq n$.
\end{lemma}

\begin{lemma}[Davis-Kahan's $\sin \bTheta$ Theorem, Theorem 2.2.1 in \cite{chen2020spectral}]\label{lem:Davis_Kahan_sin}
Let $\overline{\bM}$ and $\bM = \overline{\bM} + \bE$ be two real symmetric $n \times n$ matrices, with eigenvalue decompositions given respectively by
\begin{subequations}
    \begin{align*}
        \overline{\bM} =&\, \sum\limits_{i=1}^{n}\overline{\lambda}_i \overline{\bu}_{i} {\overline{\bu}_{i}}^{\sT} = \begin{bmatrix}
                \overline{\bU} & \overline{\bU}_{\perp}
            \end{bmatrix}
            \begin{bmatrix}
                \overline{\bLambda} & \bzero\\
                \bzero & \overline{\bLambda}_{\perp}
            \end{bmatrix}
            \begin{bmatrix}
                {\overline{\bU}}^{\sT}\\
                {\overline{\bU}}^{\sT}_{\perp}
            \end{bmatrix}\,,\\
            \bM =&\, \sum\limits_{i=1}^{n}\lambda_i \bu_i \bu_i^\sT = \begin{bmatrix}
                \bU & \bU_{\perp}
            \end{bmatrix}
            \begin{bmatrix}
                \bLambda & \bzero\\
                \bzero & \bLambda_{\perp}
            \end{bmatrix}
            \begin{bmatrix}
               \bU^{\sT}\\
                \bU_{\perp}^{\sT}
            \end{bmatrix}\,.
        \end{align*}
    \end{subequations}
    Here, $\{ \overline{\lambda}_{i}\}_{i=1}^{n}$(resp. $\{\lambda_{i}\}_{i=1}^{n}$) stand for the eigenvalues of $\overline{\bM}$(resp. $\bM$), and $\overline{\bu}_i$(resp. $\bu_i$) denotes the eigenvector associated $\overline{\lambda}_{i}$(resp. $\lambda_i$). Additionally, for some fixed integer $r\in[n]$, we denote
    \begin{align}
        \overline{\bLambda}:= \mathrm{diag}\{\overline{\lambda}_1, \dots, \overline{\lambda}_r\}, \quad \overline{\bLambda}_{\perp}:=&\, \mathrm{diag}\{\overline{\lambda}_{r+1}, \dots, \overline{\lambda}_{n}\}, \notag \\
        \overline{\bU}:=  [\overline{\bu}_1, \dots, \overline{\bu}_r] \in \R^{n \times r}, \quad \overline{\bU}_{\perp}:=&\, [\overline{\bu}_{r+1}, \dots, \overline{\bu}_{n}] \in \R^{n \times (n - r)}. \notag 
    \end{align}
    The matrices $\bLambda$, $\bLambda_{\perp}$, $\bU$, $\bU_{\perp}$ are defined analogously. Assume that
    \begin{align*}
        \mathrm{eigenvalues}(\overline{\bLambda}) \subseteq [\alpha, \beta]\,,\quad \mathrm{eigenvalues}(\bLambda_{\perp}) \subseteq (-\infty, \alpha - \Delta] \cup [\beta + \Delta, \infty),\quad \alpha, \beta \in \R\,, \Delta >0\,,
    \end{align*}
    and the projection matrices are given by $P_\bU := \bU\bU^{\sT}$, $P_{\overline{\bU}} := \overline{\bU} \,\overline{\bU}^{\sT}$, then one has $\|P_\bU - P_{\overline{\bU}}\| \leq (2\|\bE\|/\Delta)$. In particular, suppose that $|\overline{\lambda}_1| \geq |\overline{\lambda}_2| \geq \cdots \geq |\overline{\lambda}_r| \geq |\overline{\lambda}_{r+1}| \geq \cdots |\overline{\lambda}_{n}|$ (resp. $|\lambda_1| \geq \cdots \geq |\lambda_{n}|$). If $\|\bE\|\leq (1 - 1/\sqrt{2})(|\overline{\lambda}|_{r} - |\overline{\lambda}|_{r+1})$, then one has
    \begin{equation*}
         \|P_\bU - P_{\overline{\bU}}\| \leq \frac{2\|\bE\| }{ |\overline{\lambda}_{r}| - |\overline{\lambda}_{r+1}|} \,.
    \end{equation*}
\end{lemma}

\begin{lemma}[Wedin's $\sin \bTheta$ Theorem, Theorem 2.3.1 in \cite{chen2020spectral}]\label{lem:Wedin_sin}
    Let $\overline{\bM}$ and $\bM = \overline{\bM} + \bE$ be two $n_1 \times n_2$ real matrices and $n_1\geq n_2$, with SVDs given respectively by 
    \begin{subequations}
    \begin{align}
        \overline{\bM} =&\, \sum\limits_{i=1}^{n_1}\overline{\sigma}_i \overline{\bu}_{i} {\overline{\bv}_{i}}^{\sT} = \begin{bmatrix}
                \overline{\bU} & \overline{\bU}_{\perp}
            \end{bmatrix}
            \begin{bmatrix}
                \overline{\bSigma} & \bzero & \bzero\\
                \bzero & \overline{\bSigma}_{\perp} & \bzero
            \end{bmatrix}
            \begin{bmatrix}
                {\overline{\bV}}^{\sT}\\
                \overline{\bV}^{\sT}_{\perp}
            \end{bmatrix} \notag \\
            \bM =&\, \sum\limits_{i=1}^{n_1}\sigma_i \bu_i \bv_i^\sT = \begin{bmatrix}
                \bU & \bU_{\perp}
            \end{bmatrix}
            \begin{bmatrix}
                \bSigma & \bzero & \bzero\\
                \bzero & \bSigma_{\perp} & \bzero
            \end{bmatrix}
            \begin{bmatrix}
               \bV^{\sT}\\
                \bV_{\perp}^{\sT}
            \end{bmatrix} .\notag
        \end{align}
    \end{subequations}
    Here, $\overline{\sigma}_1 \geq \dots \geq \overline{\sigma}_{n_1}$ (resp. $\sigma_1 \geq \dots \geq \sigma_{n_1}$) stand for the singular values of $\overline{\bM}$(resp. $\bM$), $\overline{\bu}_i$(resp. $\bu_i$) denotes the left singular vector associated with the singular value $\overline{\sigma}_i$(resp. $\sigma_i$), and $\overline{\bv}_i$(resp. $\bv_i$) denotes the right singular vector associated with the singular value $\overline{\sigma}_i$(resp. $\sigma_i$). In addition, for any fixed integer $r\in[n]$, we denote 
    \begin{align*}
        \bSigma:= \mathrm{diag}\{\sigma_1, \dots, \sigma_r\}, \quad \bSigma_{\perp}:=&\, \mathrm{diag}\{\sigma_{r+1}, \dots, \sigma_{n_1}\},\\
        \bU:= [\bu_1, \dots, \bu_r] \in \R^{n_1 \times r}, \quad \bU_{\perp}:=&\, [\bu_{r+1}, \dots, \bu_{n_1}] \in \R^{n_1 \times (n_1 - r)},\\
        \bV:= [\bv_1, \dots, \bv_r] \in \R^{n_2 \times r}, \quad \bV_{\perp}:=&\, [\bv_{r+1}, \dots, \bv_{n_2}] \in \R^{n_2 \times (n_2 - r)}.
    \end{align*}
    The matrices $\overline{\bSigma}$, $\overline{\bSigma}_{\perp}$, $\overline{\bU}$, $\overline{\bU}_{\perp}$, $\overline{\bV}$, $\overline{\bV}_{\perp}$ are defined analogously. If $\bE = \bM - \overline{\bM}$ satisfies $\|\bE\|\leq \overline{\sigma}_r - \overline{\sigma}_{r+1}$, then with the projection matrices $P_\bU := \bU\bU^{\sT}$, one has
    \begin{equation*}
        \max \left\{\|P_\bU - P_{\overline{\bU}}\|, \|P_\bV - P_{\overline{\bV}}\| \right\} \leq \frac{\sqrt{2} \max \left\{\|\bE^\sT \overline{\bU}\|, \|\bE \overline{\bV}\| \right\} }{ \overline{\sigma}_r - \overline{\sigma}_{r+1} - \|\bE\|} \,.
    \end{equation*}
    In particular, if $\|\bE\|\leq (1 - 1/\sqrt{2})(\overline{\sigma}_{r} - \overline{\sigma}_{r+1})$, then one has
    \begin{equation*}
        \max \left\{\|P_\bU - P_{\overline{\bU}}\|, \|P_\bV - P_{\overline{\bV}}\| \right\} \leq \frac{\sqrt{2} \|\bE\| }{ \overline{\sigma}_{r} - \overline{\sigma}_{r+1}} \,.
    \end{equation*}
\end{lemma}

\begin{comment}
\subsection{Proof of Lemma \ref{lem: 2 block vector angle}}
\todo{remove this part and cite \cite{chin2015stochastic}}
\begin{proof}
    The corresponding $v$ is constructed as follows. Recall that 
    \[
    \sin \angle(W, \bar{W}) := \|P_{W} - P_{\bar{W}}\| \leq c \leq \frac{1}{4},
    \]
    and the first unit two eigenvectors of $\E A$ are given by \[\bar{v}_{1} = [\frac{1}{\sqrt{n}}, \cdots, \frac{1}{\sqrt{n}}]^{T},\] and 
    \[\bar{v}_{2} = [\frac{1}{\sqrt{n}}, \cdots,\frac{1}{\sqrt{n}}, -\frac{1}{\sqrt{n}}, \cdots, -\frac{1}{\sqrt{n}}]^{T}.\] 
    
     Define $u_{i} :=P_{W}\bar{v}_{i}$ and $x_{i} = u_{i} - \bar{v}_{i}$, then $\|x_{i}\|_2 \leq c$, since $u_{i} \in P_{W}$ and $\bar{v}_{i} \in P_{\bar{W}}$ for $i = 1, 2$. Also,
    \begin{align}
        |u_{1}^{T}u_{2}| = |\bar{v}_{1}^{T}x_{2} + \bar{v}_{2}^{T}x_{1}+ \bar{x}_{1}^{T}x_{2} + \bar{v}_{1}^{T}v_{2}| \leq |\bar{v}_{1}^{T}x_{2}| + |\bar{v}_{2}^{T}x_{1}|+ |\bar{x}_{1}^{T}x_{2}| \leq 2c + c^{2}
    \end{align}
    We define $u_{\perp} = u_{2} - u_{1}\frac{u_{1}^{T}u_{2}}{\|u_{1}\|^{2}_{2}}$. Clearly, $\|u_{\perp}\|_2 \leq \|u_2\|_2\leq \|\bar{v}_2\|_2\leq 1$ and  when $c\leq \frac{1}{4}$,
    \begin{align}
        |u_{\perp}^{T}\bar{v}_{2}| &= \left| u_{2}^{T}\bar{v}_{2} - \frac{u_{1}^{T}u_{2}}{\|u_{1}\|^{2}_{2}}u_{1}^{T}\bar{v}_{2} \right| \geq \left| u_{2}^{T}\bar{v}_{2} \right| - \left|\frac{u_{1}^{T}u_{2}}{\|u_{1}\|^{2}_{2}}u_{1}^{T}\bar{v}_{2} \right|\\
        & \geq (1 - c) - \frac{(2c + c^{2})c}{(1 - c)^{2}} \geq 1 - 2c.
    \end{align}
    Taking $v=\frac{u}{\|u\|_2}$, we obtain a unit vector $v$ with $v \perp u_{1}$ and $|v^{T}\bar{v}_{2}| \geq 1 - 2c$, which gives \[\sin \angle(v, \bar{v}_{2}) \leq \sqrt{1 - (1 - 2c)^{2}}\leq 2 \sqrt{c}.\]
\end{proof}
\end{comment}
\begin{thebibliography}{10}

\bibitem{abbe2018community}
Emmanuel Abbe.
\newblock Community detection and stochastic block models: Recent developments.
\newblock {\em Journal of Machine Learning Research}, 18(177):1--86, 2018.

\bibitem{abbe2016exact}
Emmanuel Abbe, Afonso~S Bandeira, and Georgina Hall.
\newblock Exact recovery in the stochastic block model.
\newblock {\em IEEE Transactions on Information Theory}, 62(1):471--487, 2016.

\bibitem{abbe2015community}
Emmanuel Abbe and Colin Sandon.
\newblock Community detection in general stochastic block models: Fundamental
  limits and efficient algorithms for recovery.
\newblock In {\em 2015 IEEE 56th Annual Symposium on Foundations of Computer
  Science}, pages 670--688. IEEE, 2015.

\bibitem{abbe2018proof}
Emmanuel Abbe and Colin Sandon.
\newblock Proof of the achievability conjectures for the general stochastic
  block model.
\newblock {\em Communications on Pure and Applied Mathematics},
  71(7):1334--1406, 2018.

\bibitem{ahn2016community}
Kwangjun Ahn, Kangwook Lee, and Changho Suh.
\newblock Community recovery in hypergraphs.
\newblock In {\em Communication, Control, and Computing (Allerton), 2016 54th
  Annual Allerton Conference on}, pages 657--663. IEEE, 2016.

\bibitem{ahn2018hypergraph}
Kwangjun Ahn, Kangwook Lee, and Changho Suh.
\newblock Hypergraph spectral clustering in the weighted stochastic block
  model.
\newblock {\em IEEE Journal of Selected Topics in Signal Processing},
  12(5):959--974, 2018.

\bibitem{angelini2015spectral}
Maria~Chiara Angelini, Francesco Caltagirone, Florent Krzakala, and Lenka
  Zdeborov{\'a}.
\newblock Spectral detection on sparse hypergraphs.
\newblock In {\em Communication, Control, and Computing (Allerton), 2015 53rd
  Annual Allerton Conference on}, pages 66--73. IEEE, 2015.

\bibitem{battiston2020networks}
Federico Battiston, Giulia Cencetti, Iacopo Iacopini, Vito Latora, Maxime
  Lucas, Alice Patania, Jean-Gabriel Young, and Giovanni Petri.
\newblock Networks beyond pairwise interactions: structure and dynamics.
\newblock {\em Physics Reports}, 874:1--92, 2020.

\bibitem{benson2016higher}
Austin~R Benson, David~F Gleich, and Jure Leskovec.
\newblock Higher-order organization of complex networks.
\newblock {\em Science}, 353(6295):163--166, 2016.

\bibitem{bordenave2018nonbacktracking}
Charles Bordenave, Marc Lelarge, and Laurent Massouli{\'e}.
\newblock Nonbacktracking spectrum of random graphs: Community detection and
  nonregular {R}amanujan graphs.
\newblock {\em Annals of probability}, 46(1):1--71, 2018.

\bibitem{bu2010music}
Jiajun Bu, Shulong Tan, Chun Chen, Can Wang, Hao Wu, Lijun Zhang, and Xiaofei
  He.
\newblock Music recommendation by unified hypergraph: combining social media
  information and music content.
\newblock In {\em Proceedings of the 18th ACM international conference on
  Multimedia}, pages 391--400, 2010.

\bibitem{caceres2021sparse}
Elena C{\'a}ceres, Anderson Misobuchi, and Rafael Pimentel.
\newblock Sparse {SYK} and traversable wormholes.
\newblock {\em Journal of High Energy Physics}, 2021(11):1--32, 2021.

\bibitem{chen2020spectral}
Yuxin Chen, Yuejie Chi, Jianqing Fan, and Cong Ma.
\newblock Spectral methods for data science: A statistical perspective.
\newblock {\em Foundations and Trends{\textregistered} in Machine Learning},
  14(5):566--806, 2021.

\bibitem{chien2018community}
I~Chien, Chung-Yi Lin, and I-Hsiang Wang.
\newblock Community detection in hypergraphs: Optimal statistical limit and
  efficient algorithms.
\newblock In {\em International Conference on Artificial Intelligence and
  Statistics}, pages 871--879, 2018.

\bibitem{chien2018minimax}
I~Eli Chien, Chung-Yi Lin, and I-Hsiang Wang.
\newblock On the minimax misclassification ratio of hypergraph community
  detection.
\newblock {\em IEEE Transactions on Information Theory}, 65(12):8095--8118,
  2019.

\bibitem{chin2015stochastic}
Peter Chin, Anup Rao, and Van Vu.
\newblock Stochastic block model and community detection in sparse graphs: A
  spectral algorithm with optimal rate of recovery.
\newblock In {\em Conference on Learning Theory}, pages 391--423, 2015.

\bibitem{cole2020exact}
Sam Cole and Yizhe Zhu.
\newblock Exact recovery in the hypergraph stochastic block model: A spectral
  algorithm.
\newblock {\em Linear Algebra and its Applications}, 593:45--73, 2020.

\bibitem{cook2018size}
Nicholas Cook, Larry Goldstein, and Tobias Johnson.
\newblock Size biased couplings and the spectral gap for random regular graphs.
\newblock {\em The Annals of Probability}, 46(1):72--125, 2018.

\bibitem{cooper2020adjacency}
Joshua Cooper.
\newblock Adjacency spectra of random and complete hypergraphs.
\newblock {\em Linear Algebra and its Applications}, 596:184--202, 2020.

\bibitem{dumitriu2021spectra}
Ioana Dumitriu and Yizhe Zhu.
\newblock Spectra of random regular hypergraphs.
\newblock {\em Electronic Journal of Combinatorics}, 28(3):P3.36, 2021.

\bibitem{feige2005spectral}
Uriel Feige and Eran Ofek.
\newblock Spectral techniques applied to sparse random graphs.
\newblock {\em Random Structures \& Algorithms}, 27(2):251--275, 2005.

\bibitem{friedman1995second}
Joel Friedman and Avi Wigderson.
\newblock On the second eigenvalue of hypergraphs.
\newblock {\em Combinatorica}, 15(1):43--65, 1995.

\bibitem{ghoshdastidar2014consistency}
Debarghya Ghoshdastidar and Ambedkar Dukkipati.
\newblock Consistency of spectral partitioning of uniform hypergraphs under
  planted partition model.
\newblock In {\em Advances in Neural Information Processing Systems}, pages
  397--405, 2014.

\bibitem{ghoshdastidar2017consistency}
Debarghya Ghoshdastidar and Ambedkar Dukkipati.
\newblock Consistency of spectral hypergraph partitioning under planted
  partition model.
\newblock {\em The Annals of Statistics}, 45(1):289--315, 2017.

\bibitem{ghoshdastidar2017uniform}
Debarghya Ghoshdastidar and Ambedkar Dukkipati.
\newblock Uniform hypergraph partitioning: Provable tensor methods and sampling
  techniques.
\newblock {\em The Journal of Machine Learning Research}, 18(1):1638--1678,
  2017.

\bibitem{govindu2005tensor}
Venu~Madhav Govindu.
\newblock A tensor decomposition for geometric grouping and segmentation.
\newblock In {\em 2005 IEEE Computer Society Conference on Computer Vision and
  Pattern Recognition (CVPR'05)}, volume~1, pages 1150--1157. IEEE, 2005.

\bibitem{harris2021}
Kameron~Decker Harris and Yizhe Zhu.
\newblock Deterministic tensor completion with hypergraph expanders.
\newblock {\em SIAM Journal on Mathematics of Data Science}, 3(4):1117--1140,
  2021.

\bibitem{holland1983stochastic}
Paul~W Holland, Kathryn~Blackmond Laskey, and Samuel Leinhardt.
\newblock Stochastic blockmodels: First steps.
\newblock {\em Social networks}, 5(2):109--137, 1983.

\bibitem{jain2014provable}
Prateek Jain and Sewoong Oh.
\newblock Provable tensor factorization with missing data.
\newblock In {\em Advances in Neural Information Processing Systems}, pages
  1431--1439, 2014.

\bibitem{jin2021sharp}
Jiashun Jin, Tracy Ke, and Jiajun Liang.
\newblock Sharp impossibility results for hyper-graph testing.
\newblock {\em Advances in Neural Information Processing Systems}, 34, 2021.

\bibitem{ke2020community}
Zheng~Tracy Ke, Feng Shi, and Dong Xia.
\newblock Community detection for hypergraph networks via regularized tensor
  power iteration.
\newblock {\em arXiv preprint arXiv:1909.06503}, 2019.

\bibitem{kim2018stochastic}
Chiheon Kim, Afonso~S Bandeira, and Michel~X Goemans.
\newblock Stochastic block model for hypergraphs: Statistical limits and a
  semidefinite programming approach.
\newblock {\em arXiv preprint arXiv:1807.02884}, 2018.

\bibitem{le2015estimating}
Can~M Le and Elizaveta Levina.
\newblock Estimating the number of communities in networks by spectral methods.
\newblock {\em arXiv preprint arXiv:1507.00827}, 2015.

\bibitem{le2017concentration}
Can~M Le, Elizaveta Levina, and Roman Vershynin.
\newblock Concentration and regularization of random graphs.
\newblock {\em Random Structures \& Algorithms}, 51(3):538--561, 2017.

\bibitem{lee2020robust}
Jeonghwan Lee, Daesung Kim, and Hye~Won Chung.
\newblock Robust hypergraph clustering via convex relaxation of truncated
  {MLE}.
\newblock {\em IEEE Journal on Selected Areas in Information Theory}, 2020.

\bibitem{lei2020consistent}
Jing Lei, Kehui Chen, and Brian Lynch.
\newblock Consistent community detection in multi-layer network data.
\newblock {\em Biometrika}, 107(1):61--73, 2020.

\bibitem{lei2015consistency}
Jing Lei and Alessandro Rinaldo.
\newblock Consistency of spectral clustering in stochastic block models.
\newblock {\em The Annals of Statistics}, 43(1):215--237, 2015.

\bibitem{li2013news}
Lei Li and Tao Li.
\newblock News recommendation via hypergraph learning: encapsulation of user
  behavior and news content.
\newblock In {\em Proceedings of the sixth ACM international conference on Web
  search and data mining}, pages 305--314, 2013.

\bibitem{lin2017fundamental}
Chung-Yi Lin, I~Eli Chien, and I-Hsiang Wang.
\newblock On the fundamental statistical limit of community detection in random
  hypergraphs.
\newblock In {\em Information Theory (ISIT), 2017 IEEE International Symposium
  on}, pages 2178--2182. IEEE, 2017.

\bibitem{lung2018hypergraph}
Rodica~Ioana Lung, No{\'e}mi Gask{\'o}, and Mihai~Alexandru Suciu.
\newblock A hypergraph model for representing scientific output.
\newblock {\em Scientometrics}, 117(3):1361--1379, 2018.

\bibitem{massoulie2014community}
Laurent Massouli{\'e}.
\newblock Community detection thresholds and the weak {R}amanujan property.
\newblock In {\em Proceedings of the forty-sixth annual ACM symposium on Theory
  of computing}, pages 694--703. ACM, 2014.

\bibitem{michoel2012alignment}
Tom Michoel and Bruno Nachtergaele.
\newblock Alignment and integration of complex networks by hypergraph-based
  spectral clustering.
\newblock {\em Physical Review E}, 86(5):056111, 2012.

\bibitem{mossel2015reconstruction}
Elchanan Mossel, Joe Neeman, and Allan Sly.
\newblock Reconstruction and estimation in the planted partition model.
\newblock {\em Probability Theory and Related Fields}, 162(3-4):431--461, 2015.

\bibitem{mossel2016consistency}
Elchanan Mossel, Joe Neeman, and Allan Sly.
\newblock Consistency thresholds for the planted bisection model.
\newblock {\em Electron. J. Probab}, 21(21):1--24, 2016.

\bibitem{mossel2018proof}
Elchanan Mossel, Joe Neeman, and Allan Sly.
\newblock A proof of the block model threshold conjecture.
\newblock {\em Combinatorica}, 38(3):665--708, 2018.

\bibitem{newman2002random}
Mark~EJ Newman, Duncan~J Watts, and Steven~H Strogatz.
\newblock Random graph models of social networks.
\newblock {\em Proceedings of the national academy of sciences}, 99(suppl
  1):2566--2572, 2002.

\bibitem{ng2002spectral}
Andrew~Y Ng, Michael~I Jordan, and Yair Weiss.
\newblock On spectral clustering: Analysis and an algorithm.
\newblock In {\em Advances in neural information processing systems}, pages
  849--856, 2002.

\bibitem{nguyen2015tensor}
Nam~H Nguyen, Petros Drineas, and Trac~D Tran.
\newblock Tensor sparsification via a bound on the spectral norm of random
  tensors.
\newblock {\em Information and Inference: A Journal of the IMA}, 4(3):195--229,
  2015.

\bibitem{Pal_2021}
Soumik Pal and Yizhe Zhu.
\newblock Community detection in the sparse hypergraph stochastic block model.
\newblock {\em Random Structures \& Algorithms}, 59(3):407--463, 2021.

\bibitem{sen2018optimization}
Subhabrata Sen.
\newblock Optimization on sparse random hypergraphs and spin glasses.
\newblock {\em Random Structures \& Algorithms}, 53(3):504--536, 2018.

\bibitem{shi2000normalized}
Jianbo Shi and Jitendra Malik.
\newblock Normalized cuts and image segmentation.
\newblock {\em IEEE Transactions on pattern analysis and machine intelligence},
  22(8):888--905, 2000.

\bibitem{soma2019spectral}
Tasuku Soma and Yuichi Yoshida.
\newblock Spectral sparsification of hypergraphs.
\newblock In {\em Proceedings of the Thirtieth Annual ACM-SIAM Symposium on
  Discrete Algorithms}, pages 2570--2581. SIAM, 2019.

\bibitem{tian2009hypergraph}
Ze~Tian, TaeHyun Hwang, and Rui Kuang.
\newblock A hypergraph-based learning algorithm for classifying gene expression
  and array{CGH} data with prior knowledge.
\newblock {\em Bioinformatics}, 25(21):2831--2838, 2009.

\bibitem{vershynin_2018}
Roman Vershynin.
\newblock {\em High-Dimensional Probability: An Introduction with Applications
  in Data Science}.
\newblock Cambridge Series in Statistical and Probabilistic Mathematics.
  Cambridge University Press, 2018.

\bibitem{wen2019learning}
Longyin Wen, Dawei Du, Shengkun Li, Xiao Bian, and Siwei Lyu.
\newblock Learning non-uniform hypergraph for multi-object tracking.
\newblock In {\em Proceedings of the AAAI conference on artificial
  intelligence}, volume~33, pages 8981--8988, 2019.

\bibitem{yuan2021testing}
Mingao Yuan, Ruiqi Liu, Yang Feng, and Zuofeng Shang.
\newblock Testing community structures for hypergraphs.
\newblock {\em arXiv preprint arXiv:1810.04617}, 2018.

\bibitem{yuan2021community}
Mingao Yuan, Bin Zhao, and Xiaofeng Zhao.
\newblock Community detection in censored hypergraph.
\newblock {\em arXiv preprint arXiv:2111.03179}, 2021.

\bibitem{zhang2016minimax}
Anderson~Y Zhang and Harrison~H Zhou.
\newblock Minimax rates of community detection in stochastic block models.
\newblock {\em The Annals of Statistics}, 44(5):2252--2280, 2016.

\bibitem{zhang2021exact}
Qiaosheng Zhang and Vincent~YF Tan.
\newblock Exact recovery in the general hypergraph stochastic block model.
\newblock {\em arXiv preprint arXiv:2105.04770}, 2021.

\bibitem{zhen2021community}
Yaoming Zhen and Junhui Wang.
\newblock Community detection in general hypergraph via graph embedding.
\newblock {\em arXiv preprint arXiv:2103.15035}, 2021.

\bibitem{zhou2007learning}
Denny Zhou, Jiayuan Huang, and Bernhard Sch{\"o}lkopf.
\newblock Learning with hypergraphs: Clustering, classification, and embedding.
\newblock In {\em Advances in neural information processing systems}, pages
  1601--1608, 2007.

\bibitem{zhou2021sparse}
Zhixin Zhou and Yizhe Zhu.
\newblock Sparse random tensors: Concentration, regularization and
  applications.
\newblock {\em Electronic Journal of Statistics}, 15(1):2483--2516, 2021.

\end{thebibliography}
